\documentclass[a4paper,11pt]{article}
\usepackage[T1]{fontenc}
\usepackage{lmodern}
\usepackage[margin=23mm,headheight=15pt]{geometry}
\usepackage{microtype,amsmath,amssymb,amsthm,mathtools}
\usepackage{booktabs,longtable,array,tabularx,enumitem,fancyhdr,xcolor,listings,needspace}
\usepackage{titlesec}
\usepackage{xurl}
\usepackage[unicode,breaklinks]{hyperref}
\definecolor{teal}{HTML}{174C53}
\definecolor{pale}{HTML}{F1F5F5}
\hypersetup{pdftitle={Chapter I: Machine-Native Final Companion},pdfauthor={Komorebi P.},colorlinks=true,linkcolor=teal,urlcolor=teal}
\titleformat{\section}{\Large\bfseries\color{teal}}{\thesection}{0.7em}{}
\titleformat{\subsection}{\large\bfseries\color{teal}}{\thesubsection}{0.7em}{}
\titleformat{\subsubsection}{\normalsize\bfseries\color{teal}}{\thesubsubsection}{0.7em}{}
\setcounter{tocdepth}{2}
\setlength{\parindent}{0pt}\setlength{\parskip}{0.45em}
\setlist{itemsep=0.18em,topsep=0.4em}
\pagestyle{fancy}\fancyhf{}\fancyhead[L]{\small PSC I \enspace Machine-native companion}\fancyhead[R]{\small Komorebi P.}\fancyfoot[C]{\thepage}
\lstset{basicstyle=\ttfamily\scriptsize,breaklines=true,breakatwhitespace=false,columns=fullflexible,keepspaces=true,showstringspaces=false,backgroundcolor=\color{pale},frame=single,rulecolor=\color{teal!25},framerule=0.3pt,xleftmargin=1mm,xrightmargin=1mm,aboveskip=0.6em,belowskip=0.6em}
\allowdisplaybreaks\emergencystretch=3em\sloppy
\newcommand{\bbk}{\Bbbk}\newcommand{\A}{\mathcal A}\newcommand{\B}{\mathcal B}\newcommand{\Ccal}{\mathcal C}\newcommand{\Pcal}{\mathcal P}\newcommand{\Qcal}{\mathcal Q}\newcommand{\Rcal}{\mathcal R}\newcommand{\Scal}{\mathcal S}
\newcommand{\Hull}{\operatorname{Hull}}\newcommand{\srank}{\operatorname{srank}}\newcommand{\psrank}{\operatorname{psrank}}\newcommand{\End}{\operatorname{End}}\newcommand{\Hom}{\operatorname{Hom}}\newcommand{\im}{\operatorname{im}}\newcommand{\Ker}{\operatorname{Ker}}\newcommand{\Jac}{\operatorname{Jac}}\newcommand{\Irr}{\operatorname{Irr}}\newcommand{\Span}{\operatorname{span}}\newcommand{\alg}{\operatorname{alg}}\newcommand{\Reach}{\operatorname{Reach}}\newcommand{\Supp}{\operatorname{Supp}}\newcommand{\id}{\mathrm{id}}\newcommand{\cl}[1]{\overline{#1}}
\newcommand{\field}[2]{\par\smallskip\noindent\textbf{#1.} #2\par}

% Operation Komorebi: shared publication profile.
\usepackage{komorebi-publication}
\renewcommand{\KomorebiTitle}{Projector-Resolved Seed Calculus I}
\renewcommand{\KomorebiSubtitle}{Typed operator systems, seed rank, and exact packet coordinates}
\renewcommand{\KomorebiSeries}{Chapter I}
\renewcommand{\KomorebiKind}{Computational companion}
\renewcommand{\KomorebiEdition}{}
\renewcommand{\KomorebiShortTitle}{Chapter I}
\renewcommand{\KomorebiTitleSize}{31}
\renewcommand{\KomorebiTitleLeading}{36}
\renewcommand{\KomorebiProfile}{classical}
\renewcommand{\KomorebiOrientation}{portrait}
\KomorebiSetup
\usepackage{titlesec}
% Shared publication layout. Mathematical definitions and proofs are unchanged.
\RequirePackage{setspace,caption,needspace,etoolbox}
\definecolor{pubNavy}{HTML}{3C263F}
\definecolor{pubTeal}{HTML}{7B4165}
\definecolor{pubCopper}{HTML}{A68A56}
\definecolor{pubInk}{HTML}{22282C}
\definecolor{pubMuted}{HTML}{615963}
\definecolor{pubWash}{HTML}{F5F0F4}
\newcommand{\CoreLayout}{%
  \geometry{letterpaper,left=.87in,right=.87in,top=.88in,bottom=.86in,headheight=14pt,headsep=18pt,footskip=26pt,ignorehead,ignorefoot}%
  \color{pubInk}\setstretch{1.06}%
  \setlength{\parindent}{1.05em}\setlength{\parskip}{2.5pt plus .5pt minus .3pt}%
  \setlength{\emergencystretch}{3em}\raggedbottom
  \setlength{\abovedisplayskip}{9pt plus 2pt minus 2pt}%
  \setlength{\belowdisplayskip}{9pt plus 2pt minus 2pt}%
  \setlength{\abovedisplayshortskip}{5pt plus 2pt minus 1pt}%
  \setlength{\belowdisplayshortskip}{6pt plus 2pt minus 1pt}%
  \pagestyle{empty}\fancypagestyle{plain}{\fancyhf{}\renewcommand{\headrulewidth}{0pt}\renewcommand{\footrulewidth}{0pt}}%
  \titleformat{\section}{\normalfont\bfseries\fontsize{14}{18}\selectfont\color{pubNavy}\raggedright}{\textcolor{pubTeal}{\thesection}}{.65em}{}%
  \titleformat{\subsection}{\normalfont\bfseries\fontsize{11.8}{15}\selectfont\color{pubNavy}\raggedright}{\thesubsection}{.65em}{}%
  \titleformat{\subsubsection}{\normalfont\itshape\fontsize{11.3}{14}\selectfont\color{pubNavy}\raggedright}{\thesubsubsection}{.7em}{}%
  \titleformat{\part}{\normalfont\bfseries\fontsize{18}{23}\selectfont\color{pubNavy}\raggedright}{\textcolor{pubTeal}{Part \thepart}}{.7em}{}%
  \titlespacing*{\part}{0pt}{16pt}{18pt}%
  \titlespacing*{\section}{0pt}{18pt plus 3pt minus 2pt}{8pt}%
  \titlespacing*{\subsection}{0pt}{13pt plus 2pt minus 1pt}{6pt}%
  \titlespacing*{\subsubsection}{0pt}{11pt}{4pt}%
  \ifcsname chapter\endcsname
    \titleformat{\chapter}[display]{\normalfont\color{pubNavy}}{\sffamily\fontsize{8.5}{11}\selectfont\color{pubTeal}\MakeUppercase{\chaptertitlename}\enspace\thechapter}{9pt}{\fontsize{24}{29}\selectfont\bfseries\raggedright}%
    \titlespacing*{\chapter}{0pt}{6pt}{22pt}%
  \fi
  \captionsetup{format=hang,justification=raggedright,singlelinecheck=false,font={small,sf,stretch=1.06},labelfont={bf,sf,color=pubInk},textfont={color=pubMuted},labelsep=period,skip=8pt}%
  \hypersetup{colorlinks=true,linkcolor=pubTeal,citecolor=pubTeal,urlcolor=pubTeal,breaklinks=true}%
  \clubpenalty=10000\widowpenalty=10000\displaywidowpenalty=10000
  \AtBeginEnvironment{longtable}{\ifdim\csname f@size\endcsname pt>10pt\relax\small\fi}
  \frenchspacing
}
\newcommand{\CoreOpening}[3]{%
  \thispagestyle{empty}\begingroup\setlength{\parindent}{0pt}%
  {\sffamily\fontsize{8.5}{11}\selectfont\color{pubTeal}\MakeUppercase{#1}\par}%
  \vspace{18pt}%
  {\fontsize{28}{34}\selectfont\color{pubNavy}\raggedright#2\par}%
  \vspace{12pt}%
  {\sffamily\fontsize{11}{15.5}\selectfont\color{pubMuted}\raggedright#3\par}%
  \vspace{19pt}\noindent\textcolor{pubCopper}{\rule{42pt}{.8pt}}\par
  \vspace{16pt}\endgroup\normalfont\color{pubInk}%
}

\newcommand{\CoreContext}[2]{%
 \par\begingroup\small\sffamily\color{pubMuted}\setlength{\parindent}{0pt}%
 \textbf{#1.}\enspace #2\par\endgroup\vspace{11pt}%
}

\BeforeBeginEnvironment{tabularx}{\noindent}


% Publication-design overrides. Mathematical bodies retain their source text.
\CoreLayout
\colorlet{accent}{pubTeal}\colorlet{teal}{pubTeal}\colorlet{ink}{pubInk}\colorlet{muted}{pubMuted}\colorlet{pale}{pubWash}\colorlet{light}{pubWash}
\makeatletter
\@ifundefined{KomorebiCover}{}{\renewcommand{\KomorebiCover}{%
\clearpage\hypersetup{pageanchor=false}\thispagestyle{empty}\begingroup\setlength{\parindent}{0pt}%
{\sffamily\fontsize{8.5}{11}\selectfont\color{pubTeal}REALIZATION CALCULUS\par}%
\vspace{27pt}%
{\sffamily\fontsize{9}{12}\selectfont\color{pubTeal}\MakeUppercase{\KomorebiSeries}\par}%
\vspace{22pt}%
{\normalfont\fontsize{29}{35}\selectfont\color{pubNavy}\raggedright\KomorebiTitle\par}%
\vspace{17pt}%
{\normalfont\fontsize{17}{23}\selectfont\color{pubMuted}\raggedright\KomorebiSubtitle\par}%
\vspace{22pt}\noindent\textcolor{pubCopper}{\rule{42pt}{.8pt}}\par%
\vspace{20pt}%
\@ifundefined{KomorebiCoverBlurb}{}{\normalfont\fontsize{11.4}{16.2}\selectfont\color{pubMuted}\raggedright\KomorebiCoverBlurb\par}%
\vfill%
{\normalfont\fontsize{14}{18}\selectfont\color{pubInk}Komorebi P.\par}%
\vspace{9pt}{\sffamily\fontsize{9}{13}\selectfont\color{pubMuted}\KomorebiKind\par\KomorebiEdition\par}%
\endgroup\clearpage\pagenumbering{arabic}\hypersetup{pageanchor=true}\color{pubInk}\normalfont\pagestyle{empty}}}
\makeatother

\begin{document}
\KomorebiCover


\section*{Computational records and interfaces}\phantomsection\addcontentsline{toc}{section}{Computational records and interfaces}
This computational companion represents Chapter I as a collection of mathematical records that can be inspected and queried, with a finite subset supported by executable interfaces. The statement inventory is \nolinkurl{formal/coverage/COVERAGE_108.json}. The executable runtime has its own entrypoint under \nolinkurl{source/machine_v1.0.1/}. That path identifies the exact archived payload. The catalogue and runtime have the distinct scopes explained below.

The catalogue contains 108 mathematical records, including 69 theorem, proposition, corollary, and lemma records with source proofs. The formal correspondence is reported alongside each record. The preserved runtime interface has twenty typed signatures, and its attached historical Lean tranche has eight mappings. Those inherited interfaces and receipts retain their original scope.

The mathematical source supplies the statements and proofs, including the split-semisimple hypotheses and explicit map qualifications. The supplementary worked models remain in the mathematical text. Stable identifiers connect the mathematical, formal, and computational material. Statement hashes and source locators distinguish the catalogue from the archived runtime payloads.

The verification record distinguishes evidence for the mathematical correspondence from the runtime receipts. Its source inventory and formal correspondence map are available under \nolinkurl{formal/coverage/}, and compiler and replay receipts are under \nolinkurl{formal/evidence/}. The Lean guide explains the formal sources, their hypotheses, and how to interpret those checks.

\textbf{How to use this volume.} Read Sections 1--3 for mathematical and evidence semantics. Sections 4--6 give the executable contracts and an exact worked example. Section 7 describes reproduction. Appendix A reproduces every canonical source statement and every attached proof body, with source locators and machine metadata. Later appendices give signatures, action coverage, operation vocabulary, and the historical Lean map.
\clearpage\tableofcontents\clearpage
\section{The mathematical object behind the interface}
\subsection{Operator systems, generated hulls, and seed rank}
A seed problem consists of an acting algebra, a module, and a family of seed vectors. The algebra must be fixed before the generation question is asked. For a unital algebra $A$ acting on a module $M$ and a finite seed family $S=\{s_1,\ldots,s_r\}$, the algebraic generated hull is
\[
\Hull_A(S)=\sum_{j=1}^{r}As_j.
\]
A one-seed system is cyclic. Seed rank measures the minimum cardinality of a generating family, under the regime declared in the source. Enlarging the acting algebra changes the question. Replacing algebraic span by a closure changes it again. The machine record therefore keeps the acting regime and hypotheses beside the statement.

Chapter I begins with typed operator grammars and their realizations, then passes through free presentations, generator number, packet coordinates, block algebras, transfer graphs, and realization. The resulting record graph covers the whole expanded foundation. It is not limited to the finite matrix examples that the runtime executes.

For finite families of nonnegative integers, the ambient convention is $\max\varnothing=0$. The height of an empty finite poset is zero. A nonempty poset with no strict comparable pair also has height zero, so height counts strict steps. These conventions clarify the empty cases without adding stable records.

The source separates algebraic, linear, and topological claims. A finite-dimensional matrix computation can exercise a linear identity. It cannot establish the Banach space result in the countable setting or a generator formula over every division ring. Those claims retain their own source proofs and dependencies.
\subsection{Exact packet coordinates}
For a finite index set $I$, a packet system on $X$ consists of linear maps $P_\alpha$ satisfying
\[
P_\alpha^2=P_\alpha,\qquad
P_\beta P_\alpha=0\quad(\beta\ne\alpha),\qquad
\sum_{\alpha\in I}P_\alpha=I_X.
\]
The packet spaces are the ranges $X_\alpha=P_\alpha X$. Analysis and synthesis are
\[
\mathcal A_Px=(P_\alpha x)_{\alpha\in I},\qquad
\mathcal S_P((x_\alpha)_\alpha)=\sum_\alpha x_\alpha.
\]
Their inverse laws are exact. No assumption of inner product orthogonality is needed. This distinction is visible in the oblique two-dimensional example later in this companion.

Machine execution represents each packet as an ambient vector. It does not construct a basis of every range or replace the ambient representation by smaller coordinate arrays. This choice makes exact identities easy to compare with the source formulas while retaining the packet index structure.
\subsection{Transfer blocks and composition}
For an endomorphism $T$, the transfer block from packet $\alpha$ to packet $\beta$ is
\[
T_{\beta\alpha}=P_\beta TP_\alpha.
\]
Its source is $X_\alpha$ and its target is $X_\beta$. In the finite runtime it is stored as a full ambient matrix. Completeness gives
\[
T=\sum_{\alpha,\beta}T_{\beta\alpha},\qquad
(ST)_{\gamma\alpha}=\sum_\beta S_{\gamma\beta}T_{\beta\alpha}.
\]
The order matters: the rightmost factor acts first. The runtime uses the same convention. Its nested block array is indexed first by target and then by source.

The graph records source dependencies for these identities. It does not infer every implicit mathematical dependency. The packet-system hypotheses remain required even when a compact source statement inherits them from the surrounding section.
\subsection{Intertwining and reconstruction}
For packet systems $P$ on $X$ and $Q$ on $Y$, an intertwiner $U:X\to Y$ satisfies $UT=SU$. Resolving both sides yields
\[
\sum_\beta U_{\gamma\beta}T_{\beta\alpha}
=
\sum_\delta S_{\gamma\delta}U_{\delta\alpha},
\qquad U_{\gamma\alpha}=Q_\gamma UP_\alpha.
\]
All local equations together are equivalent to the global equation under the packet hypotheses. The runtime returns the global Boolean, every local residual, and the conjunction of the local Booleans. It provides the data needed to inspect agreement or failure in a concrete exact case.
\subsection{The boundary of finite execution}
Countable packet sums require a topology and a convergence condition. Quotient descent requires invariance. Coordinate transport requires a suitable invertible map. Packet-pure seed counts differ from unrestricted seed counts because one mixed-support seed can contribute to several source components. These are mathematical hypotheses, not formatting conventions.

The source records seven explicit boundaries. Their presence is part of the interface. Runtime action names do not imply that every associated theorem has an implementation, and metadata listing a boundary does not itself execute that boundary check.

\section{Canonical machine representation}
\subsection{Entrypoint and package identity}
The mathematical records are located at \nolinkurl{formal/coverage/COVERAGE_108.json}. Appendix A is generated directly from those records. Each includes the statement, source proof where present, source location, hashes, formal status, and declaration correspondence. The archived mathematical source identified by those hashes is the statement authority.

The preserved executable payload is located at \nolinkurl{source/machine_v1.0.1/}. Runtime commands in this volume assume that working directory. Its entrypoint, \nolinkurl{PSC_I.machine.json}, points to the preserved record collection, statement AST, dependency graph, action registry, and evidence index. A runtime result therefore retains the exact interface scope explained below. It is not a certificate for every current manuscript statement.

The preserved runtime source digest is:
\begin{lstlisting}
37bf8f03dbbf66ce41f8b07a8e7be8fa5fca44176bdd56dd25b483f72736ff77
\end{lstlisting}
That digest identifies the archived runtime source, including explicit nonzero-packet assumptions for two corollaries and the empty-family and poset-height conventions. The mathematical inventory has its own statement and proof hashes, printed in Appendix A. Runtime receipts remain attached to their original bytes. They do not transfer automatically to the mathematical inventory.
\subsection{Record identity and source provenance}
Each current record has a stable ID, kind, title, statement text, optional proof text, source locator, hashes, formal status, declaration targets, and a scope note. These fields identify the actual current statement and the proof text that accompanies it. The preserved executable payload additionally supplies its own dependency annotations, semantic profiles, evidence coordinates, and fingerprints for its versioned source.

There are 28 definitions, 28 theorems, 22 corollaries, 16 propositions, three lemmas, three classical inputs, seven boundaries, and one remark. The 69 theorem/proposition/corollary/lemma records have attached proof bodies. Definitions and boundary statements are not counted as proved theorems merely because they occupy records.

A fingerprint in the preserved runtime provides a content identity under that payload's generating convention. It is not a proof certificate. The release gate checks the source digest, ID count, graph, AST coverage, selected API behavior, and exact fixtures. It does not establish all mathematical content through independent formal reconstruction.
\subsection{Statement AST and typed signatures}
The preserved runtime statement AST has one entry per stable ID in its versioned source. The statuses are 20 \nolinkurl{TYPED_CORE}, 81 \nolinkurl{STRUCTURED_SYMBOLIC}, and seven \nolinkurl{TEXTUAL_SEMANTIC}. A typed-core record points to a signature declaring parameters, requirements where supplied, and a return description.

The AST segments statements into prose and mathematical surfaces. Some expressions are recognized as relations or direct-sum expressions. Others retain \nolinkurl{latex_expr} leaves. Even a typed-core record can contain these leaves. Consequently, ``typed core'' refers to the attached interface signature, not to a complete formal elaboration or a proof term.

The signature for block intertwining describes a logical equivalence between the global and local equations. Its source proof remains attached separately. The runtime can evaluate a finite instance. These are three different representations with three different evidentiary roles.
\subsection{Dependencies and closure}
The preserved runtime dependency graph contains 108 nodes and 97 internal edges. The release gate confirms that the recorded internal dependencies resolve and that the graph is acyclic. Direct dependencies are stored on each record. The API can also return a transitive closure.

Closure queries return sorted identifiers. They do not produce a proof order, carry out the proof, or resolve external classical inputs. A consumer can use the graph to inspect a prerequisite set and then obtain source statements and evidence for each record.
\subsection{Operation profiles and actions}
Each preserved runtime record carries source-derived operation metadata, with a transitive profile alongside the direct profile. The ontology distinguishes structural operations from hypothesis classes. These annotations support planning and inspection. They should not be interpreted as a verified extraction of every semantic operation used by an arbitrary proof.

The action registry contains 37 action names associated with 75 records. Eight names are implemented in the runtime. Other entries are marked planner-only or fixture-only. The source has broader mathematical scope than the executable subset.

\section{Evidence semantics and certification scope}
\subsection{Independent evidence coordinates}
The preserved runtime evidence object separates source, KIR, Python, and Lean. Its source coordinate identifies that payload's source, and its KIR coordinate records validation of that version. Python evidence reports exact finite fixtures for 26 records. Its Lean coordinate points to eight historical source mappings.

The catalogue uses the formal correspondence map instead. Each of the 108 records carries a status and scope note, and each proof-level statement has its own correspondence assessment. The verification index distinguishes the execution evidence attached to each component.

No coordinate upgrades another automatically. Exact arithmetic removes floating-point ambiguity from a finite computation. It does not quantify over arbitrary fields, modules, dimensions, or topological spaces. A mapped Lean declaration proves its own formal statement under its dependencies. Its relation to the source statement must still be read from the mapping.
\subsection{Formal correspondence and checks}
The correspondence map classifies the 69 proof-level records as 69 covered, 0 partial, and 0 not mapped. The classifications describe the complete statements at their recorded scopes. They do not count independent theorems or executable actions. Every statement, source proof, and correspondence scope is reproduced in Appendix A. The exact declaration lists are in \nolinkurl{formal/coverage/COVERAGE_108.json}.

The recorded audit checks 116 formal source modules plus the retained boundary suite. Including that suite, it inventories 3513 source-owned entries and checks 3500 logical declarations. Of these, 3499 are checked by logical replay and one declaration sharing an upstream name receives a fresh-name kernel body check after exact type/value comparison. It inventories 13 runtime helpers separately and records 12 rejected negative controls. The allowed transitive axiom set is \nolinkurl{Classical.choice}, \nolinkurl{Quot.sound}, \nolinkurl{propext}. The receipt is \nolinkurl{formal/evidence/RELEASE_VERIFICATION.json}. Its detailed compiler and logical-replay evidence is under \nolinkurl{formal/evidence/whole_closure/}.

The recorded computational checks pass for 96 exact finite or symbolic visual assertions, 33 companion tests, and 10 regenerated PDF/PNG figure pairs. These bounded computations have their own scope and do not replace the formal declarations.

The pinned upstream dependency cache is reused. Logical replay uses Lean's official kernel. It is not an independently implemented checker. The source-to-manuscript correspondence is a separate reviewed mathematical artifact.


\subsection{Historical Lean linkage}
The attached seven-module formalization records Lean 4.34.0-rc2 and a pinned Mathlib manifest. Its eight mappings include definitions, generalized theorems, one-directional correspondence, and a related sufficient direction. The support-pattern algebra mapping is explicitly partial. It must not be counted as a full formal certificate for every direction of the source criterion.

A separate predecessor report records 35 of 35 obligations for its identified source. That lineage is not the canonical correspondence for the expanded mathematical inventory. The predecessor count, the eight attached runtime mappings, and the expanded correspondence map describe different scopes.
\subsection{Recorded runtime checks}
The documented verifier and exact fixture suite were run from the archived runtime payload. Both passed. The verification log is attached in the companion's \nolinkurl{build/} directory. All 65 files extracted from the original machine archive retain their original bytes.

An initial invocation used an interpreter without SymPy and failed when the blockification action imported that dependency. The successful replay used an existing environment containing SymPy. No runtime action implementation was changed to resolve that environment failure. The two audited statement corrections were applied separately and have their own provenance record.

Typesetting this companion does not execute Lean. Lean checks are reported in the verification index and explained in the formal companion. The historical runtime status file is retained with its original source identity.
\subsection{What PASS means}
A release-gate pass means the named package checks succeeded. A fixture pass means a concrete exact computation satisfied its assertion. A transport-gate pass means all operation flags that this runtime checked were accepted. None of these strings, taken alone, means that the entire 108-record foundation has been formally proved by the runtime.

The full source proof catalog is included so that a reader can inspect mathematical arguments beyond the runtime. Editorial reproduction of those arguments is not an adversarial proof audit or a new formal verification claim.

\section{Executable action contracts}
\subsection{Input and output conventions}
The API is a Python command that emits JSON on successful calls. Exact arithmetic actions use SymPy matrices. Integers and rational strings such as \nolinkurl{"3/7"} preserve rational values. JSON floating-point numbers should be avoided when the intended claim is an exact rational identity. Returned scalar entries are strings.

Payload shapes are conventional rather than comprehensively validated at the API boundary. Supply nonempty compatible square projector matrices and compatible operators or vectors. Malformed inputs can raise Python exceptions. The command parser does not present a universal contract for structured errors.

Most computational actions do not call packet validation automatically. Validate the packet family before interpreting block or analysis output through the theorems. With invalid projectors, a matrix expression may still be computable while its theorem interpretation fails.
\subsection{The six exact finite actions}
\begin{longtable}{@{}p{0.29\textwidth}p{0.28\textwidth}p{0.36\textwidth}@{}}
\toprule Action & Payload & Returned information\\\midrule\endhead
\nolinkurl{VALIDATE_PACKET_SYSTEM} & \nolinkurl{projectors} & Idempotence, annihilation, completeness, and their conjunction.\\
\nolinkurl{ANALYZE_PACKETS} & \nolinkurl{projectors}, \nolinkurl{vector} & Ambient packet vectors and a reconstruction Boolean.\\
\nolinkurl{BLOCKIFY_OPERATOR} & \nolinkurl{projectors}, \nolinkurl{operator} & Target/source array of ambient transfer matrices.\\
\nolinkurl{RECONSTRUCT_OPERATOR} & \nolinkurl{blocks} & Sum of every supplied block. Compatibility with a packet family is not independently checked.\\
\nolinkurl{COMPOSE_TRANSFER_BLOCKS} & \nolinkurl{projectors}, \nolinkurl{S}, \nolinkurl{T} & Array of block products and equality with directly resolved $ST$.\\
\nolinkurl{LOWER_INTERTWINER} & \nolinkurl{P}, \nolinkurl{Q}, \nolinkurl{T}, \nolinkurl{S}, \nolinkurl{U} & Global equation, local equations, residuals, and local conjunction.\\\bottomrule
\end{longtable}
\subsection{The two semantic gate actions}
\nolinkurl{CHECK_REALIZATION_PROFILE} and \nolinkurl{TRANSPORT_THEOREM} share the same runtime gate. Their payload contains a record ID and a profile. The gate reads that record's direct \nolinkurl{action_required_operations} and checks the corresponding profile flags.

The accepted flags are true, \nolinkurl{VERIFIED}, or \nolinkurl{PRESERVED}. False and \nolinkurl{NOT_PRESERVED} produce a negative list. Other or missing values are unknown. Any negative or unknown required operation blocks the gate.

This archived implementation does not evaluate the profile's \nolinkurl{verified_hypotheses}, does not use the transitive operation closure as its required set, and does not execute the returned boundary guards. It also does not construct a target theorem, a proof object, or an independently verified realization functor. The preservation flags are declarations supplied by the caller.

Thus the action named \nolinkurl{TRANSPORT_THEOREM} is an operation-profile gate. A consumer must separately justify target hypotheses, source-to-target interpretation, and any constructions used by the proof. Unknown flags fail closed only within the operation list actually inspected by the function.
\subsection{Planner-only and fixture-only actions}
The full registry in Appendix C makes execution status explicit. A planner-only action can identify a mathematically meaningful next step without implementing it. A fixture-only association records that a finite check exercises related mathematics. Calling an unimplemented action through \nolinkurl{execute} exits with an unsupported-action error.

\section{An exact worked execution}
\subsection{Oblique packets and reconstruction}
Use
\[
P_0=\begin{pmatrix}1&-1\\0&0\end{pmatrix},\qquad
P_1=\begin{pmatrix}0&1\\0&1\end{pmatrix},\qquad
x=\begin{pmatrix}3\\5\end{pmatrix}.
\]
These are complementary idempotents but are not self-adjoint under the usual Euclidean inner product. The example therefore exercises the algebraic packet hypotheses without adding orthogonality.
\begin{lstlisting}
python runtime/machine_api.py execute VALIDATE_PACKET_SYSTEM tests/packet_input.json
python runtime/machine_api.py execute ANALYZE_PACKETS tests/analyze_input.json
\end{lstlisting}

\begin{lstlisting}
{
  "complete": true,
  "idempotent": true,
  "pairwise_annihilating": true,
  "valid": true
}
\end{lstlisting}
The analysis output is exact:

\begin{lstlisting}
{
  "packets": [
    [
      [
        "-2"
      ],
      [
        "0"
      ]
    ],
    [
      [
        "5"
      ],
      [
        "5"
      ]
    ]
  ],
  "reconstructs": true
}
\end{lstlisting}

The two vectors $(-2,0)^t$ and $(5,5)^t$ sum to $x$. Each belongs to its own projector range. The runtime's zero-based packet indexing corresponds to the displayed $P_0,P_1$.
\subsection{Resolving and rebuilding an operator}
Take
\[
T=\begin{pmatrix}2&3\\5&7\end{pmatrix}.
\]
The command \nolinkurl{BLOCKIFY_OPERATOR} returns four ambient $2\times2$ matrices. In target/source order, they are
\[
T_{00}=\begin{pmatrix}-3&3\\0&0\end{pmatrix},\quad
T_{01}=\begin{pmatrix}0&-7\\0&0\end{pmatrix},\quad
T_{10}=\begin{pmatrix}5&-5\\5&-5\end{pmatrix},\quad
T_{11}=\begin{pmatrix}0&12\\0&12\end{pmatrix}.
\]
Their sum is $T$. The exact reconstruction API returns:

\begin{lstlisting}
{
  "operator": [
    [
      "2",
      "3"
    ],
    [
      "5",
      "7"
    ]
  ]
}
\end{lstlisting}

The input and output JSON files for this reconstruction are supplied in \nolinkurl{examples/}. The reconstruction action performs the block sum. Its interpretation as recovery of the original operator uses the validated packet system and the way those blocks were formed.
\subsection{Composing resolved operators}
Now set
\[
S=\begin{pmatrix}11&13\\17&19\end{pmatrix},\qquad
ST=\begin{pmatrix}87&124\\129&184\end{pmatrix}.
\]
The runtime forms the blocks of $S$ and $T$, then sums over the middle packet index. It compares the resulting array with the directly resolved product. The supplied worked output has \nolinkurl{equals_direct: true}. The full block arrays are retained in \nolinkurl{examples/composition.json}.

The comparison checks a concrete instance of the multiplication identity. The general identity follows from inserting the complete family between the factors, as in the source proof reproduced in Appendix A.
\subsection{An intertwining computation}
Let
\[
U=\begin{pmatrix}1&1\\0&1\end{pmatrix},\qquad
Q_i=UP_iU^{-1},\qquad
S'=UTU^{-1}=\begin{pmatrix}7&3\\5&2\end{pmatrix}.
\]
The target packets are $Q_0=\left(\begin{smallmatrix}1&-2\\0&0\end{smallmatrix}\right)$ and $Q_1=\left(\begin{smallmatrix}0&2\\0&1\end{smallmatrix}\right)$. The global identity is $UT=S'U$.

The recorded execution of \nolinkurl{LOWER_INTERTWINER} returns true for both the global identity and the conjunction of the four local equations. Each local residual is the zero matrix. One complete local result is:

\begin{lstlisting}
{
  "holds": true,
  "lhs": [
    [
      "-3",
      "3"
    ],
    [
      "0",
      "0"
    ]
  ],
  "residual": [
    [
      "0",
      "0"
    ],
    [
      "0",
      "0"
    ]
  ],
  "rhs": [
    [
      "-3",
      "3"
    ],
    [
      "0",
      "0"
    ]
  ],
  "source": 0,
  "target": 0
}
\end{lstlisting}

This example deliberately changes packet coordinates as well as the operator. The complete input and four outputs are retained in \nolinkurl{examples/intertwining_input.json} and \nolinkurl{examples/intertwining.json}.
\subsection{Inspecting the dependency set}
The API exposes the dependencies used by the intertwining record linked to the source:
\begin{lstlisting}
python runtime/machine_api.py deps PSC-I-T-BLOCK-INTERTWINING --transitive
\end{lstlisting}

\begin{lstlisting}
{
  "dependencies": [
    "PSC-I-D-ANALYSIS-SYNTHESIS",
    "PSC-I-D-PACKET-SYSTEM",
    "PSC-I-D-TRANSFER-BLOCK",
    "PSC-I-T-BLOCK-MULTIPLICATION",
    "PSC-I-T-EXACT-PACKET-COORDINATES",
    "PSC-I-T-OPERATOR-RECONSTRUCTION"
  ],
  "record": "PSC-I-T-BLOCK-INTERTWINING"
}
\end{lstlisting}

This is a dependency closure, not a newly generated proof. Reading each named record supplies its exact source statement, proof, and evidence coordinates.

\section{Transport profiles in use}
\subsection{Finite packet multiplication}
The packet-linear profile declares preservation of the direct operations required by packet multiplication. The command
\begin{lstlisting}
python runtime/machine_api.py transport PSC-I-T-BLOCK-MULTIPLICATION profiles/packet_linear.json
\end{lstlisting}
returns:

\begin{lstlisting}
{
  "boundary_guards": [],
  "not_preserved": [],
  "record": "PSC-I-T-BLOCK-MULTIPLICATION",
  "required_operations": [
    "composition",
    "finite_sum",
    "packet_idempotents",
    "transfer_blocks"
  ],
  "status": "PASS",
  "unknown": []
}
\end{lstlisting}

The pass applies to the listed operation flags. It does not construct the target packet system or independently check a concrete functor's preservation laws.
\subsection{A countable result is blocked}
Using the same profile for the countable unconditional packet theorem returns:

\begin{lstlisting}
{
  "boundary_guards": [],
  "not_preserved": [
    "completion",
    "topology"
  ],
  "record": "PSC-I-T-COUNTABLE-UNCONDITIONAL-PACKETS",
  "required_operations": [
    "completion",
    "packet_idempotents",
    "topology"
  ],
  "status": "BLOCKED",
  "unknown": []
}
\end{lstlisting}

The profile does not preserve the topological and completion operations required by this gate. A finite linear profile therefore cannot be used to pass this countable request.

The source theorem additionally requires a real or complex Banach space, unconditional reconstruction, and a uniform multiplier bound. These hypotheses remain visible in its statement. Supplying permissive operation flags would not prove those hypotheses. The richer source contract and the narrower executable gate must be kept distinct.
\subsection{Changing a profile without changing evidence}
A profile file is a declaration supplied to the runtime. Editing its flags changes the gate's output. It does not create mathematical preservation evidence. A downstream system that wants a transport certificate must retain the arguments or formal certificates justifying those flags and the target hypotheses.

To use the gate, identify the exact record, inspect its source and dependency context, establish a target realization, justify the relevant preservation properties and hypotheses, then use the gate as an explicit consistency check. Its output can be attached to the resulting evidence record as one component.

\section{Reproducing the computational material}
\subsection{Documented payload checks}
From \nolinkurl{source/machine_v1.0.1/}, use an interpreter containing SymPy:
\begin{lstlisting}
python runtime/verify_release.py
python runtime/exact_fixture_suite.py
python runtime/nonzero_packet_regression.py
\end{lstlisting}
The first script checks entrypoint identity, record IDs, graph acyclicity and resolution, AST coverage, consistency of the action registry, source hash, selected API queries and actions, structured transport response, and fixture replay. The second emits the detailed fixture results. The latter does not replace a theorem prover.
\subsection{Companion build}
Within the original source archive, \nolinkurl{build_companion.py} regenerates the catalogue from the mathematical inventory and builds the standalone PDF. The archived runtime builder, \nolinkurl{build_historical_companion.py}, reproduces the runtime checks and examples separately from that inventory. These archive workflows are distinct from the collected monograph build. A TeX installation with the packages declared in the preamble is required for either PDF build.

The original standalone source package includes the machine archive as provenance. The build accepts \nolinkurl{PSC_MACHINE_ARCHIVE} when that archive is stored elsewhere. The document source itself remains readable independently of the builder.
\subsection{A source change is a new evidence task}
If a source statement changes, its source and statement identity change. The KIR, AST, dependency annotations, operation profile, formal map, and tests must be reviewed for impact. A historical pass cannot be assigned to a new statement solely because the stable ID stayed the same.

The catalogue is generated from the archived mathematical source and records each statement and proof hash. The archived runtime payloads retain their own evidence. The additional traditional models and visual appendix are outside the 108-record inventory. They are not silently counted as extra machine records.
\subsection{Runtime interfaces and mathematical statements}
The source archive retains the earlier computational PDFs as provenance. This companion supplies the statement catalogue and the archived runtime's exact transport-gate specification. Planner metadata remains distinct from executed proof machinery.

\clearpage\appendix
\section{Mathematical record catalogue}\label{app:catalog}
This catalogue is generated from the archived mathematical inventory's 108 stable records. Every statement and attached proof comes from that inventory. The source locator and statement hash identify the exact archived wording. Cross-references are linked to stable records in this catalog.

The formal-status line reports the reviewed correspondence map, not the older machine payload's eight-entry Lean map. Scope notes are part of each correspondence: ``partial'' is not a formal certificate for the complete source statement. The full declaration lists and source hashes are preserved in \nolinkurl{formal/coverage/COVERAGE_108.json}.

The executable runtime contracts and historical maps remain separate evidence for their stated scopes. Reproducing a source proof in this catalog does not give that proof a Lean certificate.
\subsection{Foundations of seed generation}
\Needspace{12\baselineskip}\subsubsection*{1. Coefficient and module regimes}
\phantomsection\label{rec:PSC-I-D-COEFFICIENT-REGIME}
\field{Stable ID}{\nolinkurl{PSC-I-D-COEFFICIENT-REGIME}}
\field{Kind}{definition}
\textbf{Statement.}\par
All modules are left modules unless stated otherwise.  The algebraic regime consists of a unital ring $A$ and a unital left module $M$ over $A$.  When linear dimensions, tensor products, quivers, or operator blocks are used, $A$ is a unital algebra over $\bbk$, where $\bbk$ is a field, and $M$ is a vector space over $\bbk$.  The topological regime consists of a Hausdorff topological vector space over a complete valued field, together with a unital algebra of continuous linear operators.  Every infinite sum, completion, quotient topology, and closure below is interpreted only in the regime explicitly declared in its statement.
\field{Formal correspondence}{documentary. Regime declarations are distributed through Lean typeclass hypotheses, not one encoded regime object.}
\field{Source locator}{\nolinkurl{editable/traditional/02_foundations.tex}; statement lines 7--10.}
\field{Statement SHA-256}{\nolinkurl{3350ba8d143f60c74ec0f3c1ea24b7208cc0346c6f292b8590f68a9520efd7cb}}
\Needspace{12\baselineskip}\subsubsection*{2. Canonical theorem record}
\phantomsection\label{rec:PSC-I-D-THEOREM-RECORD}
\field{Stable ID}{\nolinkurl{PSC-I-D-THEOREM-RECORD}}
\field{Kind}{definition}
\textbf{Statement.}\par
A canonical theorem record is the tuple
\[
(\mathrm{ID},\mathrm{status},\mathrm{hypotheses},\mathrm{statement},
\mathrm{dependencies},\mathrm{proof},\mathrm{exports},\mathrm{boundary}).
\]
A downstream paper cites the stable record ID.  A domain specialization may add domain hypotheses, but it does not repeat a generic proof from this paper.
\field{Formal correspondence}{documentary. Publication metadata. JSON records are not Lean proof declarations.}
\field{Source locator}{\nolinkurl{editable/traditional/02_foundations.tex}; statement lines 14--22.}
\field{Statement SHA-256}{\nolinkurl{cd1763a1091beb27f77276664bfaa37d65f73f2b14b41b36f4b9692284853761}}
\Needspace{12\baselineskip}\subsubsection*{3. Symbolic scope}
\phantomsection\label{rec:PSC-I-B-SYMBOLIC-SCOPE}
\field{Stable ID}{\nolinkurl{PSC-I-B-SYMBOLIC-SCOPE}}
\field{Kind}{boundary}
\textbf{Statement.}\par
This edition proves the algebraic and topological statements printed below.  It does not assert historical novelty, independent external certification, automatic discovery of useful packet systems, nonlinear obstruction theory, multilevel refinement, arithmetic congruences, quantum gravity, automorphy, or local or global Langlands reciprocity.
\field{Formal correspondence}{scope boundary. Scope exclusion, not an additional positive theorem.}
\field{Source locator}{\nolinkurl{editable/traditional/02_foundations.tex}; statement lines 25--28.}
\field{Statement SHA-256}{\nolinkurl{de436ff66bb1b236ea21bf9d3cf19c10b67981488cc3ab97eadc5b05708a39ba}}
\Needspace{12\baselineskip}\subsubsection*{4. Typed operator grammar}
\phantomsection\label{rec:PSC-I-D-TYPED-GRAMMAR}
\field{Stable ID}{\nolinkurl{PSC-I-D-TYPED-GRAMMAR}}
\field{Kind}{definition}
\textbf{Statement.}\par
Let $Q=(Q_0,Q_1,s,t)$ be a quiver with finite vertex set.  A typed operator grammar over $\bbk$ is a pair
\[
\mathfrak G=(Q,I),
\]
where $I\triangleleft \bbk Q$ is an ideal on both sides generated by linear combinations over $\bbk$ of paths having common source and common target.  Write
\[
\Lambda_{\mathfrak G}:=\bbk Q/I.
\]
A realization of $\mathfrak G$ is a family $(X_v)_{v\in Q_0}$ and linear maps
\[
X_a:X_{s(a)}\longrightarrow X_{t(a)}
\qquad(a\in Q_1)
\]
that annihilate every relation in $I$.  A morphism $f:X\to Y$ of realizations is a family of linear maps $f_v:X_v\to Y_v$ satisfying
\[
f_{t(a)}X_a=Y_af_{s(a)}
\qquad(a\in Q_1).
\]
\field{Formal correspondence}{definition represented. Finite vertex set, arbitrary arrow set, source/target maps, independent vector spaces and typed arrow maps; relations are finite linear combinations of paths with fixed endpoints. The more general free-algebra relation format is specialized both to homogeneous typed path relations and to all representatives of an actual TwoSidedIdeal.}
\field{Source locator}{\nolinkurl{editable/traditional/02_foundations.tex}; statement lines 37--57.}
\field{Statement SHA-256}{\nolinkurl{282b22656cc22122458c5c399315c5468190e48cc5578ef5258b9a5246a905c9}}
\Needspace{12\baselineskip}\subsubsection*{5. Collapsed one-space realization}
\phantomsection\label{rec:PSC-I-D-COLLAPSED-REALIZATION}
\field{Stable ID}{\nolinkurl{PSC-I-D-COLLAPSED-REALIZATION}}
\field{Kind}{definition}
\textbf{Statement.}\par
For a typed realization $X$, define its realization on a single space by
\[
|X|:=\bigoplus_{v\in Q_0}X_v.
\]
Every vertex idempotent acts as a coordinate projector on $|X|$, and every path acts as a block operator from its source space to its target space.
\field{Formal correspondence}{definition represented. Total space is the finite product, canonically the finite direct sum. Vertex operators are single-after-projection maps and arrows are single-after-arrow-after-projection. Induction on the genuinely source/target-indexed TypedPath proves the actual block-composition action, including the zero-length vertex path.}
\field{Source locator}{\nolinkurl{editable/traditional/02_foundations.tex}; statement lines 60--67.}
\field{Statement SHA-256}{\nolinkurl{79fcb85b9e889dae8d8305bb688123c8a54c9d9a18a5b86d37765b76b25a82aa}}
\Needspace{12\baselineskip}\subsubsection*{6. Typed realization and module equivalence}
\phantomsection\label{rec:PSC-I-T-TYPED-MODULE-EQUIVALENCE}
\field{Stable ID}{\nolinkurl{PSC-I-T-TYPED-MODULE-EQUIVALENCE}}
\field{Kind}{theorem}
\textbf{Statement.}\par
The category of realizations of $\mathfrak G$ is equivalent to the category of unital left modules over $\Lambda_{\mathfrak G}$.  Explicitly, a realization $X$ determines the collapsed space $|X|$ of \hyperref[rec:PSC-I-D-COLLAPSED-REALIZATION]{\nolinkurl{PSC-I-D-COLLAPSED-REALIZATION}}
with the vertex idempotent $e_v$ acting as the coordinate projector onto $X_v$ and each arrow acting in the block between its declared source and target.  Conversely, a unital module $M$ over $\Lambda_{\mathfrak G}$ determines the realization
\[
M_v=e_vM,
\qquad
M_a(x)=ax.
\]
The two constructions are inverse up to canonical natural isomorphism.

\begin{proof}
Because $Q_0$ is finite, $1=\sum_{v\in Q_0}e_v$ in $\bbk Q$, the $e_v$ are pairwise orthogonal, and the same identities hold in $\Lambda_{\mathfrak G}$.  Given a realization, define the action of a path by composition of its arrow maps and extend linearly.  Every element of $I$ acts as zero, so the action factors through $\Lambda_{\mathfrak G}$.  The vertex idempotents give
\[
|X|=\bigoplus_v e_v|X|=\bigoplus_vX_v.
\]
Conversely, if $M$ is unital, then
\[
M=1M=\sum_ve_vM.
\]
Orthogonality makes the sum direct.  For an arrow $a:v\to w$, the identities $a=e_wae_v$ imply
\[
a(e_vM)\subseteq e_wM.
\]
The relations in $I$ act as zero because $M$ is a module over $\Lambda_{\mathfrak G}$.  Applying both constructions returns the original vertex spaces and arrow maps, or the original module decomposition, through the evident identity maps.  These identifications commute with morphisms, giving an equivalence of categories.
\end{proof}
\field{Formal correspondence}{covered. The grammar algebra is constructed by FreeAlgebra and RingQuot, with complete orthogonal vertex idempotents and exact source/target arrow corners. Collapse is full, faithful and essentially surjective, using actual vertex ranges and their analysis/synthesis equivalence. Action objects are then proved equivalent to ordinary ModuleCat for all unital left modules; the coefficient action is recovered through algebraMap. A separate quotient comparison uses the actual TwoSidedIdeal ring congruence and proves both inverse identities. No category equivalence, split decomposition, path action or module structure is supplied as an assumption.}
\field{Source locator}{\nolinkurl{editable/traditional/02_foundations.tex}; statement lines 70--80; proof lines 84--98.}
\field{Statement SHA-256}{\nolinkurl{a7fb9fcf8a2ff9a5b6ca94fba965afd801c3acb93f1fe6e0323b9fd3fab8470b}}
\Needspace{12\baselineskip}\subsubsection*{7. Realized operator algebra}
\phantomsection\label{rec:PSC-I-D-REALIZED-ALGEBRA}
\field{Stable ID}{\nolinkurl{PSC-I-D-REALIZED-ALGEBRA}}
\field{Kind}{definition}
\textbf{Statement.}\par
For a realization $X$, the realized operator algebra is
\[
\A_X:=\im\!\left(\Lambda_{\mathfrak G}\longrightarrow\End_{\bbk}(|X|)\right).
\]
It records the actual operators in the realization and may be a proper quotient of the abstract grammar algebra.
\field{Formal correspondence}{definition represented. The realized algebra is literally the range of the constructed algebra homomorphism into Module.End. It is not a supplied equivalent algebra or an abstract surrogate.}
\field{Source locator}{\nolinkurl{editable/traditional/02_foundations.tex}; statement lines 114--121.}
\field{Statement SHA-256}{\nolinkurl{43b505da15a63ccf3534f295a12aaadd6cd3d7d4f7e2ab5dba25d3d5cf63d0bc}}
\Needspace{12\baselineskip}\subsubsection*{8. Action factorization and hull invariance}
\phantomsection\label{rec:PSC-I-P-ACTION-FACTORIZATION}
\field{Stable ID}{\nolinkurl{PSC-I-P-ACTION-FACTORIZATION}}
\field{Kind}{proposition}
\textbf{Statement.}\par
Let $S\subseteq |X|$.  The generated submodule under the grammar algebra equals the generated submodule under the realized operator algebra:
\[
\Lambda_{\mathfrak G}S=\A_XS.
\]
Consequently, every statement about seed rank depends only on the realized action, not on redundant relations in a chosen presentation of the grammar.

\begin{proof}
The action map $\rho:\Lambda_{\mathfrak G}\twoheadrightarrow\A_X$ is surjective by definition and satisfies $\lambda s=\rho(\lambda)s$.  Hence every grammar orbit vector lies in $\A_XS$.  The reverse inclusion holds because every element of $\A_X$ is $\rho(\lambda)$ for some $\lambda$.
\end{proof}
\field{Formal correspondence}{covered. The two actual spans are compared by span induction and range-surjectivity. The cardinal seed ranks are equal via the identity semilinear bijection into the actual image algebra. Arbitrary seed sets and zero modules are included; there is no finite-generation assumption.}
\field{Source locator}{\nolinkurl{editable/traditional/02_foundations.tex}; statement lines 124--131; proof lines 135--137.}
\field{Statement SHA-256}{\nolinkurl{472b9117cf7ed5ed76713d4b131141c2e9f00b9b6dd827bcb2515705d21e94f1}}
\Needspace{12\baselineskip}\subsubsection*{9. Linear seed system}
\phantomsection\label{rec:PSC-I-D-SEED-SYSTEM}
\field{Stable ID}{\nolinkurl{PSC-I-D-SEED-SYSTEM}}
\field{Kind}{definition}
\textbf{Statement.}\par
A linear seed system is a tuple
\[
\mathfrak S=(X,\A,S,\Pcal,\tau),
\]
where $X$ is a unital left module over $\A$, $\A$ is fixed independently of the target element, $S=(s_1,\ldots,s_r)$ is a finite distinguished seed family, $\Pcal$ is either empty or a declared packet system, and $\tau$ is either the algebraic convention or a specified Hausdorff topology in which the declared operator action is continuous.
\field{Formal correspondence}{documentary. Algebraic and continuous packet cases are represented by concrete operator algebras, packet objects and arbitrary seed sets, which include every finite distinguished family after passage to its range. The record is a regime/data definition, not an assertion that one bundled structure covers the optional packet-free and topological regimes. Optional fields are represented in separate structures.}
\field{Source locator}{\nolinkurl{editable/traditional/02_foundations.tex}; statement lines 145--152.}
\field{Statement SHA-256}{\nolinkurl{f1b9660868ad56e506bd33c2ab474312509199727cc8c23a64d4334641974e26}}
\Needspace{12\baselineskip}\subsubsection*{10. Generated hull}
\phantomsection\label{rec:PSC-I-D-GENERATED-HULL}
\field{Stable ID}{\nolinkurl{PSC-I-D-GENERATED-HULL}}
\field{Kind}{definition}
\textbf{Statement.}\par
For $S\subseteq X$, define
\[
\Hull_{\A}(S):=\sum_{s\in S}\A s.
\]
In the topological regime define
\[
\Hull_{\A,\tau}(S):=\cl{\Hull_{\A}(S)}^{\,\tau}.
\]
For a singleton $S=\{s\}$, write $\Hull_{\A}(s)$.
\field{Formal correspondence}{definition represented. The scalar span of algebra-orbit vectors is proved equal as a set to the module span. The general topological operator hull is an actual Submodule and its topologicalClosure is used in the topological development.}
\field{Source locator}{\nolinkurl{editable/traditional/02_foundations.tex}; statement lines 155--166.}
\field{Statement SHA-256}{\nolinkurl{859743fdac21ebb8ba120148cbf03f5511974f6e3356b5f733a3beea316b55f5}}
\Needspace{12\baselineskip}\subsubsection*{11. Natural presentation}
\phantomsection\label{rec:PSC-I-D-NATURAL-PRESENTATION}
\field{Stable ID}{\nolinkurl{PSC-I-D-NATURAL-PRESENTATION}}
\field{Kind}{definition}
\textbf{Statement.}\par
A seed presentation is natural relative to a declared grammar when the primitive operators and relations are fixed independently of the target, the relation system is finite or effectively enumerable, the topology is intrinsic to the realization, no primitive operator directly stores the complete target, and the presentation exposes a reusable invariant, theorem, or computation not encoded by a primitive chosen for the target.
\field{Formal correspondence}{documentary. A methodological criterion rather than a decidable Lean predicate.}
\field{Source locator}{\nolinkurl{editable/traditional/02_foundations.tex}; statement lines 171--174.}
\field{Statement SHA-256}{\nolinkurl{94104ee5b6710359c3fa14a5febde34de8455668f83dfc5f49cf3e5028770033}}
\Needspace{12\baselineskip}\subsubsection*{12. Algebra generation and module generation are distinct}
\phantomsection\label{rec:PSC-I-P-ALGEBRA-MODULE-GENERATION}
\field{Stable ID}{\nolinkurl{PSC-I-P-ALGEBRA-MODULE-GENERATION}}
\field{Kind}{proposition}
\textbf{Statement.}\par
Every nonzero unital algebra $A$, regarded as a left module over itself, has
\[
\srank_A(A)=1,
\]
with cyclic seed $1_A$.  This conclusion is independent of the number of algebra generators required to present $A$ as an algebra.

\begin{proof}
For every $a\in A$, one has $a=a1_A$.  Hence $A1_A=A$.  The number of operations needed to generate $A$ as an algebra is a different invariant and does not enter the module calculation.
\end{proof}
\field{Formal correspondence}{covered. The regular action is multiplication on A. The singleton unit is proved generating by surjectivity of scalar evaluation; the nontrivial regular module then has seedRank exactly one. No assertion about algebra presentations is used as a premise.}
\field{Source locator}{\nolinkurl{editable/traditional/02_foundations.tex}; statement lines 177--184; proof lines 188--190.}
\field{Statement SHA-256}{\nolinkurl{9ba2f9c1f32694b86a7d9ac63ef11d2b3c46ee5c38307f3abcc3c12048cf91d0}}
\Needspace{12\baselineskip}\subsubsection*{13. Target-dependent operators}
\phantomsection\label{rec:PSC-I-B-TARGET-DEPENDENT-OPERATORS}
\field{Stable ID}{\nolinkurl{PSC-I-B-TARGET-DEPENDENT-OPERATORS}}
\field{Kind}{boundary}
\textbf{Statement.}\par
For a fixed seed $s$ and target $x$, the existence of a specially defined operator $A_x$ satisfying $A_xs=x$ does not establish a natural seed calculus unless $A_x$ belongs to a grammar fixed independently of the target satisfying \hyperref[rec:PSC-I-D-NATURAL-PRESENTATION]{\nolinkurl{PSC-I-D-NATURAL-PRESENTATION}}.
\field{Formal correspondence}{scope boundary. Target-independence restriction remains documentary.}
\field{Source locator}{\nolinkurl{editable/traditional/02_foundations.tex}; statement lines 192--195.}
\field{Statement SHA-256}{\nolinkurl{aece204c76d6b378ec937375eebd02a43d0a48dbed1f94054b4ed514ef573dd1}}
\Needspace{12\baselineskip}\subsubsection*{14. Seed evaluation map}
\phantomsection\label{rec:PSC-I-D-SEED-EVALUATION}
\field{Stable ID}{\nolinkurl{PSC-I-D-SEED-EVALUATION}}
\field{Kind}{definition}
\textbf{Statement.}\par
Let $A$ be a unital ring, let $M$ be a left module over $A$, and let $S=(s_1,\ldots,s_r)\in M^r$.  Define the linear map over $A$
\[
\varepsilon_S:A^{\oplus r}\longrightarrow M,
\qquad
(a_1,\ldots,a_r)\longmapsto\sum_{j=1}^{r}a_js_j.
\]
\field{Formal correspondence}{definition represented. multiSeedMap is the actual Finsupp linear combination map; finite families give the stated finite free module.}
\field{Source locator}{\nolinkurl{editable/traditional/02_foundations.tex}; statement lines 215--223.}
\field{Statement SHA-256}{\nolinkurl{db3e882c18dfdf5770d2779811295f6d287e9ffa95dbfb79bec513f96cba185d}}
\Needspace{12\baselineskip}\subsubsection*{15. Finite free presentation}
\phantomsection\label{rec:PSC-I-T-FREE-PRESENTATION}
\field{Stable ID}{\nolinkurl{PSC-I-T-FREE-PRESENTATION}}
\field{Kind}{theorem}
\textbf{Statement.}\par
For the map in \hyperref[rec:PSC-I-D-SEED-EVALUATION]{\nolinkurl{PSC-I-D-SEED-EVALUATION}},
\[
\im\varepsilon_S=\Hull_A(S),
\qquad
N_S:=\Ker\varepsilon_S\leq A^{\oplus r},
\]
and the induced map is an isomorphism of left modules over $A$
\[
\boxed{A^{\oplus r}/N_S\cong\Hull_A(S).}
\]
Consequently, $S$ generates $M$ if and only if $\varepsilon_S$ is surjective.

\begin{proof}
The image consists exactly of finite sums $\sum_ja_js_j$, which is $\Hull_A(S)$.  Because $\varepsilon_S$ is linear over $A$, its kernel is a left submodule.  The first isomorphism theorem gives
\[
A^{\oplus r}/N_S\cong\im\varepsilon_S=\Hull_A(S).
\]
The final equivalence is the equality $\im\varepsilon_S=M$.
\end{proof}
\field{Formal correspondence}{covered. Actual range/span equality, quotient equivalence and generation/surjectivity equivalence are present.}
\field{Source locator}{\nolinkurl{editable/traditional/02_foundations.tex}; statement lines 226--239; proof lines 244--250.}
\field{Statement SHA-256}{\nolinkurl{6e2db91e9df6b69a04e531d9b552fe7148f878663bcd3600c36c94ac62592ef8}}
\Needspace{12\baselineskip}\subsubsection*{16. Minimal free-cover characterization}
\phantomsection\label{rec:PSC-I-C-MINIMAL-FREE-COVER}
\field{Stable ID}{\nolinkurl{PSC-I-C-MINIMAL-FREE-COVER}}
\field{Kind}{corollary}
\textbf{Statement.}\par
For every finitely generated left module $M$ over $A$,
\[
\boxed{
\srank_A(M)=\min\{r\geq0:\text{there exists a surjection }A^{\oplus r}\twoheadrightarrow M\}.
}
\]

\begin{proof}
A seed family of size $r$ gives a surjection by \hyperref[rec:PSC-I-T-FREE-PRESENTATION]{\nolinkurl{PSC-I-T-FREE-PRESENTATION}}.  Conversely, the images of the standard basis vectors under any surjection $A^{\oplus r}\twoheadrightarrow M$ form a generating family of size at most $r$.
\end{proof}
\field{Formal correspondence}{covered. seedFinrank\_isLeast\_free\_surjection proves the finite minimum, not an executable basis algorithm.}
\field{Source locator}{\nolinkurl{editable/traditional/02_foundations.tex}; statement lines 267--275; proof lines 279--281.}
\field{Statement SHA-256}{\nolinkurl{8e4ae31565265679e9c8bf397291baf65b396ad21bb88e7f00b8d17328b1cb00}}
\Needspace{12\baselineskip}\subsubsection*{17. One-seed presentation}
\phantomsection\label{rec:PSC-I-C-ONE-SEED-PRESENTATION}
\field{Stable ID}{\nolinkurl{PSC-I-C-ONE-SEED-PRESENTATION}}
\field{Kind}{corollary}
\textbf{Statement.}\par
For $s\in M$, the annihilator
\[
\operatorname{Ann}_A(s):=\{a\in A:as=0\}
\]
is a left ideal and
\[
\boxed{A/\operatorname{Ann}_A(s)\cong\Hull_A(s).}
\]

\begin{proof}
The evaluation map for one seed is $a\mapsto as$.  Its kernel is $\operatorname{Ann}_A(s)$, which is therefore a left ideal, and \hyperref[rec:PSC-I-T-FREE-PRESENTATION]{\nolinkurl{PSC-I-T-FREE-PRESENTATION}} gives the quotient isomorphism.
\end{proof}
\field{Formal correspondence}{covered. oneSeedMap has the stated annihilator kernel and its quotient is the literal cyclic range.}
\field{Source locator}{\nolinkurl{editable/traditional/02_foundations.tex}; statement lines 283--293; proof lines 297--299.}
\field{Statement SHA-256}{\nolinkurl{92a43462490ec1465fe4129738284f9c3f53e2d8bb7b63f980c1d03a0fbbb2ed}}
\Needspace{12\baselineskip}\subsubsection*{18. Two-sided annihilator criterion}
\phantomsection\label{rec:PSC-I-P-ANNIHILATOR-TWO-SIDED}
\field{Stable ID}{\nolinkurl{PSC-I-P-ANNIHILATOR-TWO-SIDED}}
\field{Kind}{proposition}
\textbf{Statement.}\par
For $s\in M$, the left ideal $L=\operatorname{Ann}_A(s)$ is an ideal on both sides if and only if
\[
L\,\Hull_A(s)=0.
\]
If $s$ is cyclic, these conditions are equivalent to
\[
\operatorname{Ann}_A(s)=\operatorname{Ann}_A(M).
\]

\begin{proof}
The left ideal $L$ is an ideal on both sides exactly when $ab\in L$ for every $a\in L$ and $b\in A$.  This is equivalent to
\[
0=(ab)s=a(bs)
\qquad(a\in L,\ b\in A).
\]
By additivity, the last identity is equivalent to $L\Hull_A(s)=0$.

If $s$ is cyclic, then $\Hull_A(s)=M$.  Hence the preceding condition is $L\subseteq\operatorname{Ann}_A(M)$.  The reverse inclusion always holds because an operator annihilating all of $M$ annihilates $s$.  Thus equality follows.  Conversely, $\operatorname{Ann}_A(M)$ is an ideal on both sides, so equality implies that $L$ is an ideal on both sides.
\end{proof}
\field{Formal correspondence}{covered. vectorAnnihilator is the actual kernel of a\ensuremath{\mapsto}a\ensuremath{\bullet}s. Right multiplication closure is equivalent to annihilating Submodule.span A \{s\}; Ideal.IsTwoSided supplies the standard two-sided-ideal interface. Under singleton generation, the actual ideal equals Module.annihilator A M. Noncommutative order is a*b and (a*b)\ensuremath{\bullet}s=a\ensuremath{\bullet}(b\ensuremath{\bullet}s).}
\field{Source locator}{\nolinkurl{editable/traditional/02_foundations.tex}; statement lines 302--312; proof lines 316--325.}
\field{Statement SHA-256}{\nolinkurl{dd991bd6ce28ee7f94302ed9dd3fb00719cf26de3b6ccf740f441b118bdd8e54}}
\Needspace{12\baselineskip}\subsubsection*{19. Completion of a dense free presentation}
\phantomsection\label{rec:PSC-I-T-COMPLETED-PRESENTATION}
\field{Stable ID}{\nolinkurl{PSC-I-T-COMPLETED-PRESENTATION}}
\field{Kind}{theorem}
\textbf{Statement.}\par
Let $X$ be a complete Hausdorff locally convex space, let $\A\subseteq\End(X)$ act continuously, and let $S=(s_1,\ldots,s_r)$.  Put
\[
E:=\A^{\oplus r}/N_S
\]
algebraically and let $j:E\to X$ be the injective map induced by $\varepsilon_S$.  Give $E$ the locally convex topology transported from $j(E)$, equivalently the topology generated by
\[
\widetilde p(e):=p(j(e))
\]
for the continuous seminorms $p$ on $X$.  Every $a\in\A$ acts continuously on $E$ in the transported topology and therefore extends to $\widehat E$.  If $\Hull_{\A}(S)$ is dense in $X$, then $j$ extends uniquely to a topological isomorphism of modules over $\A$
\[
\boxed{\widehat E\cong X.}
\]

\begin{proof}
By \hyperref[rec:PSC-I-T-FREE-PRESENTATION]{\nolinkurl{PSC-I-T-FREE-PRESENTATION}}, $j$ is an algebraic isomorphism of modules over $\A$ from $E$ onto $\Hull_{\A}(S)$.  The transported seminorms make it a topological isomorphism onto that subspace.  For $a\in\A$ and a continuous seminorm $p$ on $X$, continuity of $a:X\to X$ supplies continuous seminorms controlling $p(aj(e))$, so the induced action on $E$ is continuous and extends uniquely to $\widehat E$.  The subspace is dense by hypothesis.  A Cauchy net in $E$ is mapped to a Cauchy net in $X$, which converges because $X$ is complete.  This defines a continuous extension $\widehat j:\widehat E\to X$.  For any $x\in X$, choose a net $(e_\lambda)$ in $E$ with $j(e_\lambda)\to x$.  The transported topology makes $(e_\lambda)$ Cauchy, so it converges to some $\xi\in\widehat E$, and continuity gives $\widehat j(\xi)=x$.  Thus $\widehat j$ is surjective.  If $\widehat j(\xi)=0$, represent $\xi$ by a Cauchy net $(e_\lambda)$.  Then $j(e_\lambda)\to0$, so every defining seminorm $\widetilde p(e_\lambda)$ tends to zero.  Thus $e_\lambda\to0$ in the completion and $\xi=0$.  Hence $\widehat j$ is bijective.  Its inverse is continuous by the universal property of completion, or directly because the topology on $E$ was transported from $X$.
\end{proof}
\field{Formal correspondence}{covered. The actual finite free evaluation A\textasciicircum{}(Fin r)\ensuremath{\to}X is used, its actual kernel is quotiented, and UniformSpace.comap pulls back the target uniformity to that quotient. Dense range is derived from dense module span. AbstractCompletion.compareEquiv supplies a uniform equivalence, linearity over A is proved by completion density, and uniqueness follows from equality on the dense image. Constant action uniform continuity is inherited through the actual quotient embedding; scalar action on Completion is the usual extension. The local-convexity pullback lemma uses ordered scalars (in particular \ensuremath{\mathbb{R}}, including underlying real spaces for complex spaces); the stronger uniform-module completion proof does not require local convexity.}
\field{Source locator}{\nolinkurl{editable/traditional/02_foundations.tex}; statement lines 333--347; proof lines 351--353.}
\field{Statement SHA-256}{\nolinkurl{23708a85485e7162e31b64ee0a8fb316bec231a1aadaa0e00f0a08f70dfe1721}}
\Needspace{12\baselineskip}\subsubsection*{20. Seed rank}
\phantomsection\label{rec:PSC-I-D-SEED-RANK}
\field{Stable ID}{\nolinkurl{PSC-I-D-SEED-RANK}}
\field{Kind}{definition}
\textbf{Statement.}\par
For a unital left module $M$ over $A$, define
\[
\srank_A(M):=\min\{|S|:AS=M\}.
\]
For a topological module $(X,\tau)$, define
\[
\srank_{\A,\tau}(X):=\min\{|S|:\cl{\A S}^{\,\tau}=X\}.
\]
Set $\srank_A(0)=0$.  When the value is finite, it is the minimal number of module generators.
\field{Formal correspondence}{definition represented. The operator hull is the scalar span of actual action values; for an acting algebra this is AS. The cardinal infimum ranges over all densely generating subsets, and existence of a minimum is proved from identity membership and cardinal well ordering. Empty generators and infinite cardinals are included. The direct moduleSeedRank definition uses Dense of the literal acting-ring module span, and a proved equality identifies it with the integer-linear operator-family definition even for a nonfaithful action.}
\field{Source locator}{\nolinkurl{editable/traditional/02_foundations.tex}; statement lines 370--381.}
\field{Statement SHA-256}{\nolinkurl{81a23b39cf2b6194ae8b0bc84ac270ee5ef86e4703b5fb5411f5fa9519aafae9}}
\Needspace{12\baselineskip}\subsubsection*{21. Cyclic characterization}
\phantomsection\label{rec:PSC-I-P-CYCLIC-CHARACTERIZATION}
\field{Stable ID}{\nolinkurl{PSC-I-P-CYCLIC-CHARACTERIZATION}}
\field{Kind}{proposition}
\textbf{Statement.}\par
For $M\neq0$,
\[
\srank_A(M)=1
\quad\Longleftrightarrow\quad
M\text{ contains a cyclic vector}.
\]

\begin{proof}
Both conditions say that there is $m\in M$ with $Am=M$.
\end{proof}
\field{Formal correspondence}{covered. The nontriviality hypothesis agrees with the current publication.}
\field{Source locator}{\nolinkurl{editable/traditional/02_foundations.tex}; statement lines 384--392; proof lines 396--398.}
\field{Statement SHA-256}{\nolinkurl{de62f3fddcf5cefb429536723a8933a119585faa47b6f2e51745ea055ae6a9cc}}
\Needspace{12\baselineskip}\subsubsection*{22. Monotonicity under enlargement of the acting algebra}
\phantomsection\label{rec:PSC-I-P-ACTING-ALGEBRA-MONOTONICITY}
\field{Stable ID}{\nolinkurl{PSC-I-P-ACTING-ALGEBRA-MONOTONICITY}}
\field{Kind}{proposition}
\textbf{Statement.}\par
Let $B\subseteq A\subseteq\End_{\bbk}(X)$ be unital operator algebras with the same identity.  Then
\[
\boxed{\srank_A(X)\leq\srank_B(X).}
\]
The same statement holds for topological seed rank when both algebras act continuously.

\begin{proof}
For every $S\subseteq X$, one has $BS\subseteq AS$.  Hence every generating family over $B$ is also a generating family over $A$.  Taking closures gives the topological statement.
\end{proof}
\field{Formal correspondence}{covered. For unital operator algebras, actual family inclusion implies operator-hull inclusion. Closure monotonicity transports a minimum generating subset. This exactly supplies the topological clause, combined with the previously mapped algebraic semilinear rank theorem.}
\field{Source locator}{\nolinkurl{editable/traditional/02_foundations.tex}; statement lines 400--407; proof lines 411--413.}
\field{Statement SHA-256}{\nolinkurl{115d0049078122c1d6c43bb8ec53237a6efae32ef13330b8e31d8bf77e68eb56}}
\Needspace{12\baselineskip}\subsubsection*{23. Quotient monotonicity}
\phantomsection\label{rec:PSC-I-T-QUOTIENT-MONOTONICITY}
\field{Stable ID}{\nolinkurl{PSC-I-T-QUOTIENT-MONOTONICITY}}
\field{Kind}{theorem}
\textbf{Statement.}\par
Let $\varphi:A\to B$ be a unital ring homomorphism and let
\[
\Phi:M\twoheadrightarrow N
\]
be a surjective additive map satisfying
\[
\Phi(am)=\varphi(a)\Phi(m)
\qquad(a\in A,\ m\in M).
\]
Then
\[
\boxed{\srank_B(N)\leq\srank_A(M).}
\]
If $M$ and $N$ are topological modules, $\Phi$ is continuous and surjective, and the actions are continuous, then
\[
\srank_{B,\tau_N}(N)\leq\srank_{A,\tau_M}(M).
\]

\begin{proof}
Let $S$ generate $M$.  For every $n\in N$, choose $m\in M$ with $\Phi(m)=n$ and write $m=\sum_ia_is_i$.  Then
\[
n=\sum_i\varphi(a_i)\Phi(s_i),
\]
so $\Phi(S)$ generates $N$.

For the topological statement, let $S$ generate $M$ densely.  Given $n=\Phi(m)$, choose a net $m_\lambda\in AS$ converging to $m$.  Continuity gives
\[
\Phi(m_\lambda)\longrightarrow n,
\]
and every $\Phi(m_\lambda)$ lies in $B\Phi(S)$.  Thus $B\Phi(S)$ is dense in $N$.
\end{proof}
\field{Formal correspondence}{covered. The new theorem starts with the printed arbitrary unital ring homomorphism, additive map, scalar compatibility, continuity, and surjectivity. Span induction maps the actual source module span into the target module span; continuous surjectivity transports density. An attained minimum generating set maps to an image of no larger cardinality. No ring-homomorphism surjectivity or faithfulness is assumed. The direct module-span rank is proved equal to the earlier operator-family rank using integer-linear acting operators. This supplies the exact topological clause alongside the separately accepted algebraic theorem.}
\field{Source locator}{\nolinkurl{editable/traditional/02_foundations.tex}; statement lines 421--440; proof lines 445--457.}
\field{Statement SHA-256}{\nolinkurl{ccde1353d26e9f7bc7aef341c35be5304729daa69f8f42274ab70e2f9d7016a7}}
\Needspace{12\baselineskip}\subsubsection*{24. Isomorphism invariance}
\phantomsection\label{rec:PSC-I-C-ISOMORPHISM-INVARIANCE}
\field{Stable ID}{\nolinkurl{PSC-I-C-ISOMORPHISM-INVARIANCE}}
\field{Kind}{corollary}
\textbf{Statement.}\par
An intertwining module isomorphism preserves algebraic seed rank.  A topological intertwining isomorphism preserves topological seed rank.

\begin{proof}
Apply \hyperref[rec:PSC-I-T-QUOTIENT-MONOTONICITY]{\nolinkurl{PSC-I-T-QUOTIENT-MONOTONICITY}} to the isomorphism and its inverse.
\end{proof}
\field{Formal correspondence}{covered. A continuous semilinear module equivalence with continuous inverse preserves the direct module-span cardinal rank, by applying the exact semilinear surjection inequality in both directions. The same-ring case uses identity ring maps. The earlier bidirectional operator-family theorem covers intertwining operator-algebra isomorphisms, and the new rank equality bridges these definitions. This supplies the exact topological clause alongside separately accepted algebraic invariance.}
\field{Source locator}{\nolinkurl{editable/traditional/02_foundations.tex}; statement lines 464--467; proof lines 471--473.}
\field{Statement SHA-256}{\nolinkurl{fc01d1773cdeebcaba8346b323df6df33d48e0f6f7adda763b14307b2e902ca1}}
\Needspace{12\baselineskip}\subsubsection*{25. Short exact sequence bounds}
\phantomsection\label{rec:PSC-I-T-EXACT-SEQUENCE-BOUNDS}
\field{Stable ID}{\nolinkurl{PSC-I-T-EXACT-SEQUENCE-BOUNDS}}
\field{Kind}{theorem}
\textbf{Statement.}\par
Let
\[
0\longrightarrow L\xrightarrow{\iota}M\xrightarrow{\pi}N\longrightarrow0
\]
be a short exact sequence of left modules over $A$.  Then
\[
\boxed{
\srank_A(N)\leq\srank_A(M)
\leq\srank_A(L)+\srank_A(N).
}
\]

\begin{proof}
The lower inequality is quotient monotonicity. For the upper inequality, choose generating families $(\ell_i)_{i\in I}$ of $L$ and $(n_j)_{j\in J}$ of $N$ with the least possible cardinalities. Choose a lift $m_j\in M$ of each $n_j$. For any $m\in M$, a finite subset $F\subseteq J$ and coefficients $a_j\in A$ give
\[
\pi(m)=\sum_{j\in F}a_jn_j.
\]
The vector $m-\sum_{j\in F}a_jm_j$ belongs to $\iota(L)$, so it is a finite linear combination of the $\iota(\ell_i)$. The union of the lifted family and the image of the generating family of $L$ therefore generates $M$. Its cardinality is at most $|I|+|J|$, giving the upper bound. This argument uses finite sums for each vector even when either generating family is infinite.
\end{proof}
\field{Formal correspondence}{covered. Actual linear f,g with injectivity, range f=ker g and surjectivity; canonical seedRank=top.spanRank. Minimum generators of N are lifted and combined with images of minimum generators of L. The proof derives generation of M and the cardinal upper bound; the lower bound is surjective monotonicity. No splitting, finite generation or desired bound is assumed.}
\field{Source locator}{\nolinkurl{editable/traditional/02_foundations.tex}; statement lines 475--488; proof lines 492--498.}
\field{Statement SHA-256}{\nolinkurl{2e4e5af8593f49860709b8eb0b2d2b88577c0bc765a8b9e0c086441aaf142348}}
\Needspace{12\baselineskip}\subsubsection*{26. Shared direct-sum bounds}
\phantomsection\label{rec:PSC-I-P-SHARED-DIRECT-SUM}
\field{Stable ID}{\nolinkurl{PSC-I-P-SHARED-DIRECT-SUM}}
\field{Kind}{proposition}
\textbf{Statement.}\par
If one algebra $A$ acts diagonally on $M\oplus N$, then
\[
\boxed{
\max\{\srank_A(M),\srank_A(N)\}
\leq\srank_A(M\oplus N)
\leq\srank_A(M)+\srank_A(N).
}
\]

\begin{proof}
The coordinate projections are intertwining quotients and give the lower bound.  For the upper bound, embed minimal generating families of $M$ as $(m_i,0)$ and of $N$ as $(0,n_j)$. Their union generates by finite sums and has cardinality at most the sum of the two seed ranks.
\end{proof}
\field{Formal correspondence}{covered. The actual diagonal A-action on M\ensuremath{\times}N realizes the binary direct sum. Projection surjectivity gives both lower bounds; inclusion of generating sets in the two factors yields cardinal addition for the upper bound.}
\field{Source locator}{\nolinkurl{editable/traditional/02_foundations.tex}; statement lines 506--516; proof lines 520--522.}
\field{Statement SHA-256}{\nolinkurl{613d0bed5e9f4890c069d8d37670d0b8c9cb3a981df5bceb95686c08dc5638c6}}
\Needspace{12\baselineskip}\subsubsection*{27. Independent direct-sum formula}
\phantomsection\label{rec:PSC-I-T-INDEPENDENT-DIRECT-SUM}
\field{Stable ID}{\nolinkurl{PSC-I-T-INDEPENDENT-DIRECT-SUM}}
\field{Kind}{theorem}
\textbf{Statement.}\par
Let $M$ be a left module over $A$ and $N$ a left module over $B$.  Under the action of the product algebra
\[
(a,b)(m,n)=(am,bn),
\]
one has
\[
\boxed{
\srank_{A\times B}(M\oplus N)
=
\max\{\srank_A(M),\srank_B(N)\}.
}
\]

\begin{proof}
The coordinate projections give the lower bound. Let $r$ be the maximum on the right, regarded as a cardinal, and choose an index set $I$ of cardinality $r$. Choose minimal generating families for $M$ and $N$, index them by subsets of $I$, and fill any unused positions with zeros. Write the resulting families as $(m_i)_{i\in I}$ and $(n_i)_{i\in I}$. For any $m\in M$ and $n\in N$, their finite expressions involve a common finite subset $F\subseteq I$ after zero coefficients have been added. Thus
\[
m=\sum_{i\in F}a_im_i,
\qquad
n=\sum_{i\in F}b_in_i,
\]
and
\[
(m,n)=\sum_{i\in F}(a_i,b_i)(m_i,n_i).
\]
The family of paired vectors generates and has cardinality at most $r$. The empty family handles the case in which both modules are zero.
\end{proof}
\field{Formal correspondence}{covered. Actual product-algebra action and the maximum formula are proved.}
\field{Source locator}{\nolinkurl{editable/traditional/02_foundations.tex}; statement lines 524--538; proof lines 542--554.}
\field{Statement SHA-256}{\nolinkurl{b6f497941ccc0be69e3780d70cd1d79cb39363dedb591b81527183724fc58748}}
\Needspace{12\baselineskip}\subsubsection*{28. Tensor-product bound}
\phantomsection\label{rec:PSC-I-P-TENSOR-BOUND}
\field{Stable ID}{\nolinkurl{PSC-I-P-TENSOR-BOUND}}
\field{Kind}{proposition}
\textbf{Statement.}\par
Let $A$ and $B$ be unital algebras over $\bbk$, let $M$ be a left module over $A$, and let $N$ be a left module over $B$.  Under the natural $A\otimes_{\bbk}B$ action,
\[
\boxed{
\srank_{A\otimes B}(M\otimes_{\bbk}N)
\leq
\srank_A(M)\srank_B(N).
}
\]

\begin{proof}
Choose minimal generating families $(m_i)_{i\in I}$ for $M$ and $(n_j)_{j\in J}$ for $N$. Each vector in either module is a finite linear combination of its generators. Bilinearity therefore expresses every simple tensor as a finite sum of terms
\[
(am_i)\otimes(bn_j)=(a\otimes b)(m_i\otimes n_j).
\]
Simple tensors span the algebraic tensor product, so the family $(m_i\otimes n_j)_{(i,j)\in I\times J}$ generates it. This family has cardinality at most $|I||J|$, which is the claimed bound. The proof also covers an empty generating family, when the corresponding module and tensor product are zero.
\end{proof}
\field{Formal correspondence}{covered. The genuine tensor-product action and cardinal product upper bound are proved.}
\field{Source locator}{\nolinkurl{editable/traditional/02_foundations.tex}; statement lines 557--567; proof lines 571--577.}
\field{Statement SHA-256}{\nolinkurl{b1b31d79b3ffc3d73319922c4a3f2f1dde12a0f774ceda1ae29227e9c63a6d4e}}
\Needspace{12\baselineskip}\subsubsection*{29. Tensor rank is not generally multiplicative}
\phantomsection\label{rec:PSC-I-B-TENSOR-NONEQUALITY}
\field{Stable ID}{\nolinkurl{PSC-I-B-TENSOR-NONEQUALITY}}
\field{Kind}{boundary}
\textbf{Statement.}\par
The inequality in \hyperref[rec:PSC-I-P-TENSOR-BOUND]{\nolinkurl{PSC-I-P-TENSOR-BOUND}} can be strict.  Equality requires additional hypotheses on the acting algebras and their module multiplicities and is not exported as a universal law.
\field{Formal correspondence}{scope boundary. No universal tensor equality is exported. This boundary is not counted as a positive theorem.}
\field{Source locator}{\nolinkurl{editable/traditional/02_foundations.tex}; statement lines 584--587.}
\field{Statement SHA-256}{\nolinkurl{9328501ddd219b9e0defabc4233a3e0c6437a8890f6920385de876bc5223e204}}
\Needspace{12\baselineskip}\subsubsection*{30. Scalar-extension monotonicity}
\phantomsection\label{rec:PSC-I-P-SCALAR-EXTENSION-MONOTONICITY}
\field{Stable ID}{\nolinkurl{PSC-I-P-SCALAR-EXTENSION-MONOTONICITY}}
\field{Kind}{proposition}
\textbf{Statement.}\par
Let \(K/\bbk\) be a field extension, let \(A\) be a unital algebra over \(\bbk\), and let \(M\) be a finitely generated left module over \(A\).  Set
\[
A_K:=K\otimes_{\bbk}A,
\qquad
M_K:=K\otimes_{\bbk}M.
\]
Then
\[
\boxed{
\srank_{A_K}(M_K)\leq\srank_A(M).
}
\]
Every generating family \(S\subseteq M\) extends to the generating family \(1\otimes S\subseteq M_K\).

\begin{proof}
Let \(S=\{s_1,\ldots,s_r\}\) generate \(M\).  Every pure tensor \(\lambda\otimes m\in M_K\) has an expression
\[
m=\sum_{i=1}^r a_is_i,
\]
so
\[
\lambda\otimes m
=
\sum_{i=1}^r(\lambda\otimes a_i)(1\otimes s_i).
\]
Pure tensors span \(M_K\), hence \(1\otimes S\) generates \(M_K\) over \(A_K\).  Taking the minimum over \(S\) proves the inequality.
\end{proof}
\field{Formal correspondence}{covered. The source constructs the natural (K\ensuremath{\otimes}k A)-action on K\ensuremath{\otimes}k M using commuting tensor-factor actions and proves (c\ensuremath{\otimes}a)\ensuremath{\bullet}(d\ensuremath{\otimes}m)=(c*d)\ensuremath{\otimes}(a\ensuremath{\bullet}m). Tensor and span induction prove the literal image (1\ensuremath{\otimes}-) of any generating set generates. Minimum-cardinality generation plus cardinality of an image gives the rank inequality.}
\field{Source locator}{\nolinkurl{editable/traditional/02_foundations.tex}; statement lines 592--607; proof lines 612--624.}
\field{Statement SHA-256}{\nolinkurl{89013f407adf17424cd55db8a6f59a4a9141e34d39a1d350731d9e18d542a3e4}}
\Needspace{12\baselineskip}\subsubsection*{31. Seed-rank witness}
\phantomsection\label{rec:PSC-I-D-SEED-RANK-WITNESS}
\field{Stable ID}{\nolinkurl{PSC-I-D-SEED-RANK-WITNESS}}
\field{Kind}{definition}
\textbf{Statement.}\par
A surjective intertwiner $\Phi:M\twoheadrightarrow N$ is a witness for seed rank of order $r$ when $\srank(N)=r$.  It certifies $\srank(M)\geq r$.
\field{Formal correspondence}{documentary. A named use of the quotient lower bound, not a new proof target.}
\field{Source locator}{\nolinkurl{editable/traditional/02_foundations.tex}; statement lines 632--635.}
\field{Statement SHA-256}{\nolinkurl{81381ceea3d420ff6ba2f0b9378863598308dc392cf597cc1dc0b4d7f84f7cc0}}
\Needspace{12\baselineskip}\subsubsection*{32. Scalar quotient witness}
\phantomsection\label{rec:PSC-I-P-SCALAR-QUOTIENT-WITNESS}
\field{Stable ID}{\nolinkurl{PSC-I-P-SCALAR-QUOTIENT-WITNESS}}
\field{Kind}{proposition}
\textbf{Statement.}\par
Let $M$ be a module over $A$ with coefficient field $\bbk$.  Suppose $M$ has a quotient $N$ of dimension $r$ on which $A$ acts through a character
\[
\chi:A\longrightarrow\bbk.
\]
Then
\[
\boxed{\srank_A(M)\geq r.}
\]

\begin{proof}
For every $n\in N$, the cyclic submodule $An$ is contained in the scalar span $\bbk n$.  Therefore $q$ seed images span a subspace of dimension at most $q$.  Generating the quotient of dimension $r$ requires at least $r$ seeds.  Quotient monotonicity transfers the lower bound to $M$.
\end{proof}
\field{Formal correspondence}{covered. The character, surjective semilinear map and finite scalar dimension match the printed hypotheses.}
\field{Source locator}{\nolinkurl{editable/traditional/02_foundations.tex}; statement lines 638--648; proof lines 652--654.}
\field{Statement SHA-256}{\nolinkurl{92e35ac7e226b34fce4aa0e1e40a198c1273cc439ac42a8ea2f437f75965a632}}
\Needspace{12\baselineskip}\subsubsection*{33. Nakayama lemma}
\phantomsection\label{rec:PSC-I-CI-NAKAYAMA}
\field{Stable ID}{\nolinkurl{PSC-I-CI-NAKAYAMA}}
\field{Kind}{classicalinput}
\textbf{Statement.}\par
Let $A$ be a unital ring, let $J\subseteq\Jac(A)$, and let $M$ be a finitely generated left module over $A$.  If $M=N+JM$ for a submodule $N\leq M$, then $M=N$.
\field{Formal correspondence}{classical input integrated. The noncommutative input is integrated at the printed finitely generated left-module scope. For any left ideal I contained in Ring.jacobson R, the actual submodule equality N \ensuremath{\sqcup} I \ensuremath{\bullet} top = top implies N = top. No commutativity, Noetherian, semisimplicity, finite length, or two-sidedness of I is required for this lifting theorem. The proof derives the radical lifting result from maximal submodules and the ring-radical action inclusion, rather than assuming the desired conclusion.}
\field{Source locator}{\nolinkurl{editable/traditional/02_foundations.tex}; statement lines 669--672.}
\field{Statement SHA-256}{\nolinkurl{68901a557c34762568711f0e6c8b400485a0c4b66477b8f231e899d9a48f5878}}
\Needspace{12\baselineskip}\subsubsection*{34. Generator number is detected modulo the Jacobson radical}
\phantomsection\label{rec:PSC-I-T-MODULE-TOP}
\field{Stable ID}{\nolinkurl{PSC-I-T-MODULE-TOP}}
\field{Kind}{theorem}
\textbf{Statement.}\par
Let $A$ be a unital ring, let $J=\Jac(A)$, and let $M$ be a finitely generated left module over $A$.  A family $S\subseteq M$ generates $M$ if and only if its image generates $M/JM$ as a module over $A/J$.  When $A/J$ is semisimple Artinian, this quotient is the usual semisimple top.  No semisimplicity of $A/J$ is required for the general equality.  Consequently,
\[
\boxed{
\srank_A(M)=\srank_{A/J}(M/JM).
}
\]

\begin{proof}
If $S$ generates $M$, its image generates every quotient.  Conversely, suppose the image of $S$ generates $M/JM$ and put $N=AS$.  Then
\[
M=N+JM.
\]
Nakayama's lemma gives $M=N$.  Thus the same finite families generate before and after passage to the top, and the minimal cardinalities agree.
\end{proof}
\field{Formal correspondence}{covered. The complete generating-set/family equivalence and canonical cardinal seedRank equality are proved for arbitrary unital rings and finitely generated left modules using the literal M / (Jac(R) \ensuremath{\bullet} top), its actual scalar action over R / Jac(R), and an exact quotient-map kernel. Any lifts of a generating quotient family generate upstairs, and both minimum-cardinality inequalities are established. The natural seedFinrank equality is proved as well. No semisimplicity hypothesis is needed. The descriptive semisimple-top sentence is covered by the pinned general IsSemisimpleRing.isSemisimpleModule instance applied to this actual quotient module. This is not an algorithm for computing the radical or finding generators.}
\field{Source locator}{\nolinkurl{editable/traditional/02_foundations.tex}; statement lines 675--683; proof lines 688--694.}
\field{Statement SHA-256}{\nolinkurl{8e3170bbaf10331d7d774260fbdf4b7df0ca59f1c23116f008a1aa5ac2bf07a8}}
\Needspace{12\baselineskip}\subsubsection*{35. Local algebra formula}
\phantomsection\label{rec:PSC-I-C-LOCAL-ALGEBRA-FORMULA}
\field{Stable ID}{\nolinkurl{PSC-I-C-LOCAL-ALGEBRA-FORMULA}}
\field{Kind}{corollary}
\textbf{Statement.}\par
Let $A$ be a local ring with Jacobson radical $J$ and residue division ring $D=A/J$.  For every finitely generated left module $M$ over $A$,
\[
\boxed{
\srank_A(M)=\dim_D(M/JM),
}
\]
where the dimension is that of a left vector space over $D$.

\begin{proof}
The module $M/JM$ over $A/J$ is a left vector space over $D$.  Its minimal number of generators is its dimension over $D$.
\end{proof}
\field{Formal correspondence}{covered. The local formula is proved for noncommutative local rings and finitely generated left modules. From IsLocalRing R the source proves uniqueness of the maximal left ideal, identifies the actual Jacobson radical as maximal, constructs DivisionRing (R / Jac(R)) internally, and proves the canonical seedRank equals the left vector-space Module.rank of the actual radical quotient. The natural finite-rank equality is also proved. No commutative residue-field surrogate or caller-supplied division-ring certificate is used.}
\field{Source locator}{\nolinkurl{editable/traditional/02_foundations.tex}; statement lines 709--718; proof lines 722--724.}
\field{Statement SHA-256}{\nolinkurl{50bd885a16ef40862a42e9ddd15b7dd03758b5ee514ee98aab28acf22c1e26ed}}
\Needspace{12\baselineskip}\subsubsection*{36. Wedderburn--Artin decomposition}
\phantomsection\label{rec:PSC-I-CI-WEDDERBURN-ARTIN}
\field{Stable ID}{\nolinkurl{PSC-I-CI-WEDDERBURN-ARTIN}}
\field{Kind}{classicalinput}
\textbf{Statement.}\par
Every semisimple Artinian ring is isomorphic to a finite product
\[
A\cong\prod_{i=1}^{q}M_{d_i}(D_i)
\]
of full matrix rings over division rings.
\field{Formal correspondence}{classical input integrated. The actual arbitrary-division-ring Wedderburn-Artin existence theorem is now integrated at its semisimple-ring scope. The finite-algebra radical quotient is proved semisimple and its actual decomposition is constructed. These depend on proved pinned Mathlib results, not a custom existence axiom. Compatible finite module decompositions are separately constructed in ProductPresentation and SemisimpleExistence.}
\field{Source locator}{\nolinkurl{editable/traditional/02_foundations.tex}; statement lines 727--734.}
\field{Statement SHA-256}{\nolinkurl{effdfe288b991fcc5cdf7549800403bbb6fd698be57f017880428042cc23a058}}
\Needspace{12\baselineskip}\subsubsection*{37. Regular-module multiplicities}
\phantomsection\label{rec:PSC-I-L-REGULAR-MODULE-MULTIPLICITIES}
\field{Stable ID}{\nolinkurl{PSC-I-L-REGULAR-MODULE-MULTIPLICITIES}}
\field{Kind}{lemma}
\textbf{Statement.}\par
Let
\[
A=\prod_{i=1}^{q}M_{d_i}(D_i)
\]
and let $S_i=D_i^{d_i}$ denote the natural simple left module for the $i$th factor, extended by zero across the other factors.  Then, as a left module over $A$,
\[
\boxed{A\cong\bigoplus_{i=1}^{q}S_i^{\oplus d_i}.}
\]

\begin{proof}
For one factor $M_d(D)$, let $e_{jj}$ be the diagonal matrix units.  The left regular module decomposes as
\[
M_d(D)=\bigoplus_{j=1}^{d}M_d(D)e_{jj}.
\]
Each minimal left ideal $M_d(D)e_{jj}$ is isomorphic to the natural column module $D^d$.  Taking the direct product of factors gives the stated decomposition.
\end{proof}
\field{Formal correspondence}{covered. The actual left regular module is explicitly equivalent to its natural columns, then to the displayed finite direct sum of natural-module copies over arbitrary division rings. Coordinate ring homomorphisms define the inflated actions, other factors act by zero, and positive-dimensional natural factors are proved simple. No transpose-as-ring-map or commutative-scalar shortcut is used. Finite products for the inner repeated copies are the canonical finite direct-sum model.}
\field{Source locator}{\nolinkurl{editable/traditional/02_foundations.tex}; statement lines 737--747; proof lines 751--757.}
\field{Statement SHA-256}{\nolinkurl{cd1bf9aa7ec7ad71f24491e1969c318f330a268fb73a29873ff2e9494eb91799}}
\Needspace{12\baselineskip}\subsubsection*{38. Semisimple seed-rank formula over arbitrary division rings}
\phantomsection\label{rec:PSC-I-T-SEMISIMPLE-SEED-RANK}
\field{Stable ID}{\nolinkurl{PSC-I-T-SEMISIMPLE-SEED-RANK}}
\field{Kind}{theorem}
\textbf{Statement.}\par
Let
\[
A\cong\prod_{i=1}^{q}M_{d_i}(D_i)
\]
be semisimple Artinian and let
\[
M\cong\bigoplus_{i=1}^{q}S_i^{\oplus m_i}
\]
be a finitely generated left module over $A$.  Then
\[
\boxed{
\srank_A(M)
=
\max_{1\leq i\leq q}
\left\lceil\frac{m_i}{d_i}\right\rceil.
}
\]
If $A$ is any algebra of finite dimension over a field, the same formula applied to the decomposition of $M/\Jac(A)M$ over $A/\Jac(A)$ computes $\srank_A(M)$.

\begin{proof}
By \hyperref[rec:PSC-I-C-MINIMAL-FREE-COVER]{\nolinkurl{PSC-I-C-MINIMAL-FREE-COVER}}, $M$ is generated by $r$ elements exactly when it is a quotient of $A^{\oplus r}$.  By \hyperref[rec:PSC-I-L-REGULAR-MODULE-MULTIPLICITIES]{\nolinkurl{PSC-I-L-REGULAR-MODULE-MULTIPLICITIES}},
\[
A^{\oplus r}\cong\bigoplus_iS_i^{\oplus rd_i}.
\]
Every short exact sequence of modules over a semisimple ring splits.  Therefore a quotient isomorphic to $M$ exists if and only if
\[
m_i\leq rd_i
\qquad\text{for every }i.
\]
The smallest such $r$ is the displayed maximum.  For a nonsemisimple algebra of finite dimension, \hyperref[rec:PSC-I-T-MODULE-TOP]{\nolinkurl{PSC-I-T-MODULE-TOP}} reduces the calculation to the semisimple module $M/\Jac(A)M$ over $A/\Jac(A)$.
\end{proof}
\field{Formal correspondence}{covered. Actual product matrix ring with left multiplication and natural column modules acted on through the matching factor. Semilinear eM along eR enforces compatibility. DivisionRows proves left generation iff row spanning and m<=r*d; standard rows prove sufficiency. Product assembly, column rearrangement and seed-rank invariance prove the formula. RadicalTop is literally M/(Jac(R) smul top); existence constructs decompositions.}
\field{Source locator}{\nolinkurl{editable/traditional/02_foundations.tex}; statement lines 759--779; proof lines 784--795.}
\field{Statement SHA-256}{\nolinkurl{c604c6473ceeb6b651f37301bd3e719d69a555f04c9e4aea5fa6ad49bcd35de4}}
\Needspace{12\baselineskip}\subsubsection*{39. Algebraically closed multiplicity formula}
\phantomsection\label{rec:PSC-I-C-ALGEBRAICALLY-CLOSED-FORMULA}
\field{Stable ID}{\nolinkurl{PSC-I-C-ALGEBRAICALLY-CLOSED-FORMULA}}
\field{Kind}{corollary}
\textbf{Statement.}\par
Let $\bbk$ be algebraically closed, let $A$ be a unital algebra of finite dimension over $\bbk$, put $J=\Jac(A)$, and let $M$ be a finitely generated left module over $A$.  Assume compatible decompositions indexed by an integer $q\geq0$:
\[
A/J\cong\prod_{i=1}^{q}M_{d_i}(\bbk),
\qquad
M/JM\cong\bigoplus_{i=1}^{q}(\bbk^{d_i})^{\oplus m_i},
\]
where the $d_i$ are positive integers and the $m_i$ are nonnegative integers.  The first isomorphism is an isomorphism of algebras over $\bbk$.  The second is an isomorphism of left modules over $A/J$, using the natural column module for each matrix factor and zero action from the other factors.  Then
\[
\srank_A(M)=\max_{1\leq i\leq q}\left\lceil\frac{m_i}{d_i}\right\rceil,
\]
with $\max\varnothing=0$ when $q=0$.

\begin{proof}
By \hyperref[rec:PSC-I-T-MODULE-TOP]{\nolinkurl{PSC-I-T-MODULE-TOP}}, $\srank_A(M)=\srank_{A/J}(M/JM)$.  Apply \hyperref[rec:PSC-I-T-SEMISIMPLE-SEED-RANK]{\nolinkurl{PSC-I-T-SEMISIMPLE-SEED-RANK}} to the displayed decompositions of $A/J$ and $M/JM$, with every $D_i=\bbk$.  This gives the stated maximum.  If $q=0$, the unital $A/J$-module $M/JM$ is zero, and \hyperref[rec:PSC-I-T-MODULE-TOP]{\nolinkurl{PSC-I-T-MODULE-TOP}} gives seed rank zero.
\end{proof}
\field{Formal correspondence}{covered. Correct public map uses decomposition of A/J, not A. The supplied algebra equivalence forgets to a ring equivalence, and the module decomposition stays compatible. Specializing the exact radical-top theorem to all D\_i=k gives the current statement.}
\field{Source locator}{\nolinkurl{editable/traditional/02_foundations.tex}; statement lines 809--822; proof lines 826--828.}
\field{Statement SHA-256}{\nolinkurl{20c07c2db5ab4c4a956355fe5e45008edff100a4b279dd27a37ddff90871fcbf}}
\Needspace{12\baselineskip}\subsubsection*{40. Full matrix algebra}
\phantomsection\label{rec:PSC-I-C-FULL-MATRIX}
\field{Stable ID}{\nolinkurl{PSC-I-C-FULL-MATRIX}}
\field{Kind}{corollary}
\textbf{Statement.}\par
For the natural module $S=D^d$ over $M_d(D)$,
\[
\srank_{M_d(D)}(S)=1.
\]
More strongly, every nonzero vector of $S$ is cyclic.

\begin{proof}
The formula gives $\lceil1/d\rceil=1$.  Directly, let $v\neq0$ and choose a coordinate $v_j\neq0$.  For any target $w=(w_i)$, the matrix whose $j$th column has entries $w_iv_j^{-1}$ and whose other columns vanish sends $v$ to $w$.
\end{proof}
\field{Formal correspondence}{covered. For every division ring and positive matrix size, the natural module has canonical seed rank one and every nonzero vector generates. The proof uses entries w\_i / v\_j with the inverse on the correct right side. The positive-size convention of the surrounding natural-simple-module discussion is essential: size zero has seed rank zero. This guard remains explicit in Lean and is tested, not silently discarded.}
\field{Source locator}{\nolinkurl{editable/traditional/02_foundations.tex}; statement lines 830--837; proof lines 841--843.}
\field{Statement SHA-256}{\nolinkurl{f22f88ae7b9fe79b50f6c82420d301d784c941d6340fc700723841d724826055}}
\Needspace{12\baselineskip}\subsubsection*{41. Multiplicity obstruction}
\phantomsection\label{rec:PSC-I-C-MULTIPLICITY-OBSTRUCTION}
\field{Stable ID}{\nolinkurl{PSC-I-C-MULTIPLICITY-OBSTRUCTION}}
\field{Kind}{corollary}
\textbf{Statement.}\par
For the diagonal action of $M_d(D)$ on $S^{\oplus m}$,
\[
\boxed{
\srank_{M_d(D)}(S^{\oplus m})
=
\left\lceil\frac{m}{d}\right\rceil.
}
\]

\begin{proof}
Apply \hyperref[rec:PSC-I-T-SEMISIMPLE-SEED-RANK]{\nolinkurl{PSC-I-T-SEMISIMPLE-SEED-RANK}} to the single simple type with multiplicity $m$.
\end{proof}
\field{Formal correspondence}{covered. For arbitrary division rings and positive block size, the actual shared matrix action on m natural-module copies has canonical rank ceil(m/d). The rectangular-matrix, natural-column-product and finite-direct-sum formulations are connected by actual module equivalences. Zero multiplicity is included. This is not the independent-factor action, whose bound is a maximum across distinct factors.}
\field{Source locator}{\nolinkurl{editable/traditional/02_foundations.tex}; statement lines 845--855; proof lines 859--861.}
\field{Statement SHA-256}{\nolinkurl{df87d9465fd05e359ca1a96ba5b69e501f5e6b9b93b2d4e1b555a84476ec84ef}}
\Needspace{12\baselineskip}\subsubsection*{42. Primary decomposition for polynomial modules}
\phantomsection\label{rec:PSC-I-CI-PID-PRIMARY-DECOMPOSITION}
\field{Stable ID}{\nolinkurl{PSC-I-CI-PID-PRIMARY-DECOMPOSITION}}
\field{Kind}{classicalinput}
\textbf{Statement.}\par
Every module $X$ over $\bbk[t]$ with finite dimension over the coefficient field decomposes as
\[
X=\bigoplus_{p\in\Irr(\bbk[t])}X_p,
\]
where $\Irr(\bbk[t])$ denotes the monic irreducible polynomials, only finitely many components are nonzero, and every nonzero $X_p$ is a finite direct sum of modules $\bbk[t]/(p^e)$.
\field{Formal correspondence}{classical input integrated. For an existing polynomial module, multiplication by X constructs its evaluation operator and an actual polynomial-linear equivalence; finite coefficient dimension gives torsion. PID decomposition constructs positive elementary divisors, while actual intrinsic primary submodules are proved equal to ideal primary components. PolynomialMonic proves a bijection between monic irreducibles and nonzero prime ideals via p\ensuremath{\mapsto}(p), deriving injectivity from monicity and surjectivity from a normalized PID prime generator. It reindexes the internal sum and finite nonzero support to precisely the manuscript index set. intrinsic\_primary\_finite\_power\_sum is a literal finite DirectSum of k[X]/(p\textasciicircum{}e), with positive e, for each specified irreducible p.}
\field{Source locator}{\nolinkurl{editable/traditional/02_foundations.tex}; statement lines 883--890.}
\field{Statement SHA-256}{\nolinkurl{446dc677476b213cc4daac554de2b3f15ab816c9d0d6eeff98a67965eea843a9}}
\Needspace{12\baselineskip}\subsubsection*{43. Single-operator seed rank}
\phantomsection\label{rec:PSC-I-T-SINGLE-OPERATOR}
\field{Stable ID}{\nolinkurl{PSC-I-T-SINGLE-OPERATOR}}
\field{Kind}{theorem}
\textbf{Statement.}\par
Let $X$ have finite dimension over $\bbk$ and let $T\in\End_{\bbk}(X)$.  Regard $X$ as a module over $\bbk[t]$ by $t\cdot x=Tx$.  For every monic irreducible $p$, let $X_p$ be the primary component for $p$ and set
\[
\mu_p(T):=
\dim_{\bbk[t]/(p)}\bigl(X_p/p(T)X_p\bigr).
\]
Then
\[
\boxed{
\srank_{\bbk[T]}(X)=\max_p\mu_p(T).
}
\]
Moreover, $\mu_p(T)$ is the number of cyclic primary summands for $p$ in the decomposition into elementary divisors of $X_p$.

\begin{proof}
Let the minimal polynomial be
\[
m_T(t)=\prod_{p}p(t)^{e_p}.
\]
The acting algebra is
\[
\bbk[T]\cong\bbk[t]/(m_T)
\cong\prod_pR_p,
\qquad
R_p:=\bbk[t]/(p^{e_p}),
\]
by the Chinese remainder theorem.  The factor $R_p$ acts on $X_p$ and annihilates every other primary component.  Hence the formula for a direct sum with independent actions gives
\[
\srank_{\bbk[T]}(X)=\max_p\srank_{R_p}(X_p).
\]
Each $R_p$ is local with maximal ideal generated by the class of $p$ and residue division field $K_p=\bbk[t]/(p)$.  By \hyperref[rec:PSC-I-C-LOCAL-ALGEBRA-FORMULA]{\nolinkurl{PSC-I-C-LOCAL-ALGEBRA-FORMULA}},
\[
\srank_{R_p}(X_p)
=
\dim_{K_p}X_p/p(T)X_p
=
\mu_p(T).
\]
Finally, every summand $\bbk[t]/(p^e)$ contributes one copy of $K_p$ after quotienting by $p$, so the dimension equals the number of summands in the elementary divisor decomposition.
\end{proof}
\field{Formal correspondence}{covered. The full printed statement is now mapped. The actual power-annihilation submodule is proved equal to the ideal primary component. Its inclusion into V/pV has exact kernel pX\_p and is onto from the internally discharged Artinian hypothesis. The induced quotient equivalence is residue-field linear. Every intrinsic multiplicity is identified with every actual presentation fibre. The grouped factors are proved to be the literal X\_p, and any positive primary decomposition of X\_p has the same intrinsic count. The presentation-free maximum over all monic irreducibles is proved bounded and attained, and equals both natural-valued and cardinal-valued seed rank. No caller-supplied decomposition or seed-count conclusion is needed by the public maximum theorem. These are noncomputable mathematical constructions. The printed proof uses a different route through local-module/Nakayama machinery; that separately listed machinery is not thereby marked complete.}
\field{Source locator}{\nolinkurl{editable/traditional/02_foundations.tex}; statement lines 893--907; proof lines 914--939.}
\field{Statement SHA-256}{\nolinkurl{e96051e5d8333ea1ba1efe94f70b260ad356a41bc5ecad49525ec1df9c759967}}
\Needspace{12\baselineskip}\subsubsection*{44. Invariant-factor count}
\phantomsection\label{rec:PSC-I-C-INVARIANT-FACTOR-COUNT}
\field{Stable ID}{\nolinkurl{PSC-I-C-INVARIANT-FACTOR-COUNT}}
\field{Kind}{corollary}
\textbf{Statement.}\par
If $X\neq0$ and the rational canonical decomposition is written with nonconstant invariant factors
\[
X\cong\bigoplus_{j=1}^{r}\bbk[t]/(f_j),
\qquad
f_1\mid f_2\mid\cdots\mid f_r,
\]
then
\[
\boxed{\srank_{\bbk[T]}(X)=r.}
\]

\begin{proof}
For each irreducible $p$, the number of invariant factors divisible by $p$ equals the number of elementary divisors for $p$.  Because $f_1\mid\cdots\mid f_r$ and $f_1\neq1$ after zero summands are removed, some irreducible divisor of $f_1$ divides all $r$ invariant factors.  Hence $\max_p\mu_p(T)=r$.
\end{proof}
\field{Formal correspondence}{covered. Complete at the displayed conditional scope: an actual polynomial-module equivalence to the finite product of nonconstant factors and their divisibility chain imply least seed number, natural seed rank and cardinal seed rank equal the factor count. The theorem also handles the empty product. Finite products and finite direct sums represent the same finite decomposition. Existence of an ordered invariant-factor presentation is not asserted by this declaration, and is not an additional hypothesis-free claim of the printed conditional corollary.}
\field{Source locator}{\nolinkurl{editable/traditional/02_foundations.tex}; statement lines 951--963; proof lines 967--969.}
\field{Statement SHA-256}{\nolinkurl{c26c3a48531f1cf7f79090c82c8db185af548fc1137f1c427ac815d1d8120e89}}
\Needspace{12\baselineskip}\subsubsection*{45. Geometric-multiplicity formula}
\phantomsection\label{rec:PSC-I-C-ALGEBRAICALLY-CLOSED-SINGLE-OPERATOR}
\field{Stable ID}{\nolinkurl{PSC-I-C-ALGEBRAICALLY-CLOSED-SINGLE-OPERATOR}}
\field{Kind}{corollary}
\textbf{Statement.}\par
If $\bbk$ is algebraically closed, then
\[
\boxed{
\srank_{\bbk[T]}(X)
=
\max_{\lambda\in\operatorname{Spec}(T)}
\dim_{\bbk}\Ker(T-\lambda I).
}
\]

\begin{proof}
The irreducibles are $p=t-\lambda$.  On the primary component for $\lambda$, the quotient by $(T-\lambda I)X_\lambda$ has dimension equal to the number of Jordan blocks, which is also $\dim\Ker(T-\lambda I)$.
\end{proof}
\field{Formal correspondence}{covered. Complete at the finite-dimensional algebraically closed scope inherited from the single-operator theorem. The actual intrinsic multiplicity for X-C(a) is identified with the dimension of the actual eigenspace, without diagonalizability. The canonical natural and cardinal seed ranks equal the maximum over the actual spectrum. Off-spectrum dimensions are zero, the empty-spectrum/zero-space maximum is defined as zero, and the spectral maximum is attained on nonzero spaces. No executable spectral algorithm is claimed.}
\field{Source locator}{\nolinkurl{editable/traditional/02_foundations.tex}; statement lines 971--982; proof lines 986--988.}
\field{Statement SHA-256}{\nolinkurl{babe43b45ebd6abb23a0e31ed41c0918a4babb97578847c852f2c9d11e69515f}}
\Needspace{12\baselineskip}\subsubsection*{46. Cyclic-matrix criterion}
\phantomsection\label{rec:PSC-I-C-CYCLIC-MATRIX}
\field{Stable ID}{\nolinkurl{PSC-I-C-CYCLIC-MATRIX}}
\field{Kind}{corollary}
\textbf{Statement.}\par
For $X\ne0$ of finite dimension, the following are equivalent:
\begin{enumerate}[label=(\roman*)]
\item $T$ has a cyclic vector.
\item $\srank_{\bbk[T]}(X)=1$.
\item Every nonzero primary component for $p$ is cyclic.
\item The rational canonical form has one invariant factor.
\item The minimal polynomial equals the characteristic polynomial.
\item After scalar extension to an algebraic closure, every eigenvalue has geometric multiplicity one.
\end{enumerate}

\begin{proof}
The equivalence of (i) and (ii) is \hyperref[rec:PSC-I-P-CYCLIC-CHARACTERIZATION]{\nolinkurl{PSC-I-P-CYCLIC-CHARACTERIZATION}}.  By \hyperref[rec:PSC-I-T-SINGLE-OPERATOR]{\nolinkurl{PSC-I-T-SINGLE-OPERATOR}}, (ii) holds exactly when every $\mu_p(T)$ belonging to a nonzero primary component equals one, which is (iii).  Relatively prime primary cyclic modules combine into one cyclic module by the Chinese remainder theorem, so (iii) is equivalent to (iv).  In rational canonical form, one invariant factor means that the unique invariant factor is both the minimal and characteristic polynomial, proving (iv)$\Leftrightarrow$(v).  Scalar extension preserves the equality of minimal and characteristic polynomials.  Over an algebraic closure, that equality is equivalent to one Jordan block for each eigenvalue, which is (vi).
\end{proof}
\field{Formal correspondence}{covered. Block 06 closes the exact six-way criterion for a nonzero finite-dimensional operator over an arbitrary field. The final List.TFAE includes actual cyclicity, canonical seed rank one, cyclicity of every nonzero intrinsic primary component, an actual monic positive-degree one-invariant-factor presentation, minpoly=charpoly, and one-dimensional eigenspaces at every spectral value of the literal tensor extension to AlgebraicClosure k. Minimal-polynomial preservation and cyclicity preservation/descent are proved, not supplied as hypotheses. The one-factor presentation has a constructed basis in which the original operator is the standard companion matrix; its factor is proved to be the minimal polynomial. A separate theorem proves that any actual nonconstant invariant-factor chain has exactly one factor iff the operator is cyclic. This is a mathematical existence and equivalence theorem, not an executable normal-form algorithm or a full sorted invariant-factor construction for arbitrary noncyclic operators. The nonzero-space guard is retained in the six-way statement. Legacy conditional wrappers are preserved unchanged for compatibility, but the new public criterion needs no decomposition, multiplicity, polynomial-preservation, spectral, or seed-number certificates.}
\field{Source locator}{\nolinkurl{editable/traditional/02_foundations.tex}; statement lines 990--1001; proof lines 1005--1007.}
\field{Statement SHA-256}{\nolinkurl{d418194e254211c8f964adf5796418ad7aaff14e630caffc05cb8d935de9803a}}
\clearpage\subsection{Packets and transfer graphs}
\Needspace{12\baselineskip}\subsubsection*{47. Finite packet system}
\phantomsection\label{rec:PSC-I-D-PACKET-SYSTEM}
\field{Stable ID}{\nolinkurl{PSC-I-D-PACKET-SYSTEM}}
\field{Kind}{definition}
\textbf{Statement.}\par
A finite packet system on a vector space $X$ over $\bbk$ is a family
\[
\Pcal=(P_\alpha)_{\alpha\in I}\subseteq\End_{\bbk}(X)
\]
satisfying
\[
P_\alpha^2=P_\alpha,
\qquad
P_\alpha P_\beta=0\quad(\alpha\neq\beta),
\qquad
\sum_{\alpha\in I}P_\alpha=I_X.
\]
Set $X_\alpha=P_\alpha X$.  Zero packet spaces may be removed without changing any identity.  No orthogonality is assumed unless stated.
\field{Formal correspondence}{definition represented. LinearProjectorFamily encodes the actual finite completeness, idempotence and annihilation laws.}
\field{Source locator}{\nolinkurl{editable/traditional/03_packets_and_graphs.tex}; statement lines 10--25.}
\field{Statement SHA-256}{\nolinkurl{251dfa1ba22ba22eaf53fdcc7880aa1fe05933738f362ef3f6ae72583f539459}}
\Needspace{12\baselineskip}\subsubsection*{48. Packet support and packet-pure vectors}
\phantomsection\label{rec:PSC-I-D-PACKET-SUPPORT}
\field{Stable ID}{\nolinkurl{PSC-I-D-PACKET-SUPPORT}}
\field{Kind}{definition}
\textbf{Statement.}\par
For a finite packet system \(\Pcal=(P_\alpha)_{\alpha\in I}\) and \(x\in X\), define
\[
\Supp_{\Pcal}(x)
:=
\{\alpha\in I:P_\alpha x\neq0\}.
\]
The vector \(x\) is \emph{packet pure} if \(|\Supp_{\Pcal}(x)|\leq1\).  A finite family is packet pure if each of its members is packet pure, although different members may lie in different packets.
\field{Formal correspondence}{definition represented. packetSupport is exactly \{i | P\_i x\ensuremath{\ne}0\}. PacketPure is Set.Subsingleton of that support, equivalent for finite support to cardinality at most one. packetPure\_iff\_zero\_or\_range identifies purity with zero or membership in one actual packet range; a pure family is pointwise this predicate.}
\field{Source locator}{\nolinkurl{editable/traditional/03_packets_and_graphs.tex}; statement lines 28--37.}
\field{Statement SHA-256}{\nolinkurl{d95377c31eb89901afc468fd8f0ba293fee6865b0ae3b22eff99df183cc9eb9e}}
\Needspace{12\baselineskip}\subsubsection*{49. Finite packet aggregation}
\phantomsection\label{rec:PSC-I-L-PACKET-AGGREGATION}
\field{Stable ID}{\nolinkurl{PSC-I-L-PACKET-AGGREGATION}}
\field{Kind}{lemma}
\textbf{Statement.}\par
Let $\sigma:I\twoheadrightarrow J$ be a surjection and define
\[
Q_j:=\sum_{\sigma(\alpha)=j}P_\alpha.
\]
Then $(Q_j)_{j\in J}$ is a finite packet system and
\[
Q_jX=\bigoplus_{\sigma(\alpha)=j}X_\alpha.
\]

\begin{proof}
Pairwise annihilation of the $P_\alpha$ gives
\[
Q_j^2=Q_j,
\qquad
Q_jQ_k=0\quad(j\neq k).
\]
Summing over $j$ recovers $\sum_\alpha P_\alpha=I$.  The range statement follows from the directness forced by pairwise annihilation.
\end{proof}
\field{Formal correspondence}{covered. Q\_j is the actual filtered sum of P\_i. Computing P\_a Q\_j proves the laws. Actual range analysis/synthesis over the literal fibre subtype are inverse, and the finite product/direct-sum equivalence gives the printed conclusion.}
\field{Source locator}{\nolinkurl{editable/traditional/03_packets_and_graphs.tex}; statement lines 45--55; proof lines 59--67.}
\field{Statement SHA-256}{\nolinkurl{504cbd4eec7a5b5e416f91dcfd13e06586e5c1360990ed480424f4a85dfb9195}}
\Needspace{12\baselineskip}\subsubsection*{50. Packet analysis and synthesis}
\phantomsection\label{rec:PSC-I-D-ANALYSIS-SYNTHESIS}
\field{Stable ID}{\nolinkurl{PSC-I-D-ANALYSIS-SYNTHESIS}}
\field{Kind}{definition}
\textbf{Statement.}\par
For a finite packet system, define
\[
\mathcal A_{\Pcal}:X\longrightarrow\bigoplus_{\alpha\in I}X_\alpha,
\qquad
x\longmapsto(P_\alpha x)_\alpha,
\]
and
\[
\mathcal S_{\Pcal}:\bigoplus_{\alpha\in I}X_\alpha\longrightarrow X,
\qquad
(x_\alpha)_\alpha\longmapsto\sum_\alpha x_\alpha.
\]
\field{Formal correspondence}{definition represented. Analysis and synthesis use genuine projector ranges, not unconstrained packet slots.}
\field{Source locator}{\nolinkurl{editable/traditional/03_packets_and_graphs.tex}; statement lines 69--83.}
\field{Statement SHA-256}{\nolinkurl{f364fcc94b357d761a0bb880e5a2d5761fe7bfc687b768b0f4f7cc261a5ce820}}
\Needspace{12\baselineskip}\subsubsection*{51. Exact packet-coordinate isomorphism}
\phantomsection\label{rec:PSC-I-T-EXACT-PACKET-COORDINATES}
\field{Stable ID}{\nolinkurl{PSC-I-T-EXACT-PACKET-COORDINATES}}
\field{Kind}{theorem}
\textbf{Statement.}\par
For every finite packet system,
\[
\boxed{
\mathcal S_{\Pcal}\mathcal A_{\Pcal}=I_X,
\qquad
\mathcal A_{\Pcal}\mathcal S_{\Pcal}=I_{\oplus X_\alpha}.
}
\]
Hence
\[
\boxed{X\cong\bigoplus_{\alpha\in I}X_\alpha.}
\]

\begin{proof}
For $x\in X$,
\[
\mathcal S_{\Pcal}\mathcal A_{\Pcal}x
=
\sum_\alpha P_\alpha x
=x.
\]
If $x_\alpha\in X_\alpha$, then $P_\beta x_\alpha=\delta_{\alpha\beta}x_\alpha$, so
\[
\mathcal A_{\Pcal}\mathcal S_{\Pcal}(x_\alpha)_\alpha
=
\left(P_\beta\sum_\alpha x_\alpha\right)_\beta
=(x_\beta)_\beta.
\]
Thus analysis and synthesis are inverse linear maps.
\end{proof}
\field{Formal correspondence}{covered. Both inverse identities and the linear equivalence are proved.}
\field{Source locator}{\nolinkurl{editable/traditional/03_packets_and_graphs.tex}; statement lines 86--100; proof lines 105--121.}
\field{Statement SHA-256}{\nolinkurl{d50988e4b39d3f37422e1d0c2704aca3b8acd89beca5afd9d07ad6bc326fc24f}}
\Needspace{12\baselineskip}\subsubsection*{52. Packetwise equality}
\phantomsection\label{rec:PSC-I-C-PACKETWISE-EQUALITY}
\field{Stable ID}{\nolinkurl{PSC-I-C-PACKETWISE-EQUALITY}}
\field{Kind}{corollary}
\textbf{Statement.}\par
For $x,y\in X$,
\[
x=y
\quad\Longleftrightarrow\quad
P_\alpha x=P_\alpha y\text{ for every }\alpha.
\]

\begin{proof}
The forward implication is immediate.  For the reverse implication, sum the packet equalities and use $\sum_\alpha P_\alpha=I$.
\end{proof}
\field{Formal correspondence}{covered. Exact equality iff every projected coordinate agrees.}
\field{Source locator}{\nolinkurl{editable/traditional/03_packets_and_graphs.tex}; statement lines 128--136; proof lines 140--142.}
\field{Statement SHA-256}{\nolinkurl{61919a93863e4ac0a41653a394e44cd8f701eac05335283ac8cf73983da43859}}
\Needspace{12\baselineskip}\subsubsection*{53. Universal mapping property of finite packet coordinates}
\phantomsection\label{rec:PSC-I-C-PACKET-UNIVERSAL-PROPERTY}
\field{Stable ID}{\nolinkurl{PSC-I-C-PACKET-UNIVERSAL-PROPERTY}}
\field{Kind}{corollary}
\textbf{Statement.}\par
Let \(\Pcal=(P_\alpha)_{\alpha\in I}\) be a finite packet system on \(X\), and let \(Y\) be a vector space over \(\bbk\).
\begin{enumerate}
\item Every family of linear maps \(f_\alpha:X_\alpha\to Y\) determines a unique linear map \(f:X\to Y\) with \(f|_{X_\alpha}=f_\alpha\), namely
\[
f=\sum_\alpha f_\alpha P_\alpha.
\]
\item Every family of linear maps \(g_\alpha:Y\to X_\alpha\) determines a unique linear map \(g:Y\to X\) with \(P_\alpha g=g_\alpha\), namely
\[
g=\sum_\alpha g_\alpha.
\]
\end{enumerate}
Equivalently,
\[
\Hom(X,Y)\cong\prod_\alpha\Hom(X_\alpha,Y),
\qquad
\Hom(Y,X)\cong\prod_\alpha\Hom(Y,X_\alpha),
\]
where the products are finite.

\begin{proof}
For \(x=\sum_\alpha P_\alpha x\), any extension of the \(f_\alpha\) must satisfy
\[
f(x)=\sum_\alpha f_\alpha(P_\alpha x),
\]
which proves both existence and uniqueness of the first formula.  For the second, the sum \(g=\sum_\alpha g_\alpha\) satisfies
\[
P_\beta g
=\sum_\alpha P_\beta g_\alpha
=g_\beta
\]
because \(g_\alpha(Y)\subseteq X_\alpha\).  Uniqueness follows from \(\sum_\alpha P_\alpha=I\).
\end{proof}
\field{Formal correspondence}{covered. Outgoing extension is the finite sum of f\_i after actual rangeRestrict projections. Incoming assembly is actual rangeSynthesis after LinearMap.pi. Packet laws give restrictions/projections and completeness gives uniqueness. These universal properties induce the displayed Hom equivalences.}
\field{Source locator}{\nolinkurl{editable/traditional/03_packets_and_graphs.tex}; statement lines 146--166; proof lines 171--183.}
\field{Statement SHA-256}{\nolinkurl{3136b91ef330400f72205441080d679179476d4620287e1bd3658a36dca7a4be}}
\Needspace{12\baselineskip}\subsubsection*{54. Tensor-product packet systems}
\phantomsection\label{rec:PSC-I-T-TENSOR-PACKET-SYSTEM}
\field{Stable ID}{\nolinkurl{PSC-I-T-TENSOR-PACKET-SYSTEM}}
\field{Kind}{theorem}
\textbf{Statement.}\par
Let \(\Pcal=(P_\alpha)_{\alpha\in I}\) be a finite packet system on \(X\), and let \(\Qcal=(Q_\gamma)_{\gamma\in J}\) be a finite packet system on \(Y\).  Then
\[
\Pcal\boxtimes\Qcal
:=
(P_\alpha\otimes Q_\gamma)_{(\alpha,\gamma)\in I\times J}
\]
is a finite packet system on the algebraic tensor product \(X\otimes_{\bbk}Y\), with packet spaces
\[
(P_\alpha\otimes Q_\gamma)(X\otimes Y)
=X_\alpha\otimes Y_\gamma.
\]
Hence
\[
\boxed{
X\otimes Y
\cong
\bigoplus_{(\alpha,\gamma)\in I\times J}
X_\alpha\otimes Y_\gamma.
}
\]

\begin{proof}
For every pair of packet indices,
\[
(P_\alpha\otimes Q_\gamma)^2
=P_\alpha^2\otimes Q_\gamma^2
=P_\alpha\otimes Q_\gamma.
\]
If \((\alpha,\gamma)\neq(\beta,\delta)\), then either \(P_\alpha P_\beta=0\) or \(Q_\gamma Q_\delta=0\), so the corresponding tensor idempotents annihilate each other.  Finally,
\[
\sum_{\alpha,\gamma}P_\alpha\otimes Q_\gamma
=
\left(\sum_\alpha P_\alpha\right)
\otimes
\left(\sum_\gamma Q_\gamma\right)
=I_X\otimes I_Y.
\]
The range is generated by pure tensors \(P_\alpha x\otimes Q_\gamma y\), hence equals \(X_\alpha\otimes Y_\gamma\).  Exact reconstruction follows from \hyperref[rec:PSC-I-T-EXACT-PACKET-COORDINATES]{\nolinkurl{PSC-I-T-EXACT-PACKET-COORDINATES}}.
\end{proof}
\field{Formal correspondence}{covered. Actual TensorProduct.map projectors satisfy all laws by tensor induction. Tensor products of range inclusions have the tensor products of projections as left inverses; range\_map identifies the actual packet range. This proves the canonical embedded interpretation of the printed range equality. Exact coordinates then give the literal DirectSum.}
\field{Source locator}{\nolinkurl{editable/traditional/03_packets_and_graphs.tex}; statement lines 188--210; proof lines 215--232.}
\field{Statement SHA-256}{\nolinkurl{afea7128594f24b1249382b028d823051a35d9bd3a3473c247a7b9bd63c5ba47}}
\Needspace{12\baselineskip}\subsubsection*{55. Dual packet systems}
\phantomsection\label{rec:PSC-I-P-DUAL-PACKET-SYSTEM}
\field{Stable ID}{\nolinkurl{PSC-I-P-DUAL-PACKET-SYSTEM}}
\field{Kind}{proposition}
\textbf{Statement.}\par
Let \(X\) be finite dimensional and let \(\Pcal=(P_\alpha)\) be a finite packet system.  The algebraic dual maps
\[
P_\alpha^*:X^*\to X^*,
\qquad
P_\alpha^*(\varphi)=\varphi\circ P_\alpha,
\]
form a finite packet system on \(X^*\).  Its packet with index \(\alpha\) is canonically isomorphic to \(X_\alpha^*\), and
\[
X^*\cong\bigoplus_\alpha X_\alpha^*.
\]

\begin{proof}
Because dualization reverses products,
\[
P_\alpha^*P_\beta^*=(P_\beta P_\alpha)^*
=\delta_{\alpha\beta}P_\alpha^*.
\]
Also \(\sum_\alpha P_\alpha^*=I_{X^*}\).  Thus the dual maps form a packet system.  Restriction gives a map from \(P_\alpha^*X^*\) to \(X_\alpha^*\).  It is injective because a covector in the range of \(P_\alpha^*\) vanishes on every other packet, and it is surjective because any covector on \(X_\alpha\) extends to \(X\) through the direct sum decomposition.
\end{proof}
\field{Formal correspondence}{covered. Actual dualMap means precomposition. Composition reverses indices and orthogonality uses the reversed inequality correctly. Restriction and precomposition with the actual packet projection are inverse, giving the stated packet duals and literal finite DirectSum.}
\field{Source locator}{\nolinkurl{editable/traditional/03_packets_and_graphs.tex}; statement lines 234--246; proof lines 251--258.}
\field{Statement SHA-256}{\nolinkurl{ade34fd8142ece280ab2ae85913a943cb28a3661978659fde4a644eb5018d119}}
\Needspace{12\baselineskip}\subsubsection*{56. Finite Banach packet isomorphism}
\phantomsection\label{rec:PSC-I-T-BANACH-PACKET-ISOMORPHISM}
\field{Stable ID}{\nolinkurl{PSC-I-T-BANACH-PACKET-ISOMORPHISM}}
\field{Kind}{theorem}
\textbf{Statement.}\par
Let $X$ be Banach and let $(P_\alpha)_{\alpha\in I}$ be a finite family of bounded pairwise annihilating idempotents with sum $I_X$.  Then every range $X_\alpha=P_\alpha X$ is closed.  For any finite product norm equivalent to the maximum norm on $\bigoplus_\alpha X_\alpha$, packet analysis and synthesis are bounded inverse operators.

\begin{proof}
The range of a bounded idempotent is its space of fixed points,
\[
P_\alpha X=\Ker(I-P_\alpha),
\]
which is closed.  Because $I$ is finite, analysis and synthesis are bounded for the maximum and sum norms, respectively, and therefore for every equivalent finite product norm.  Their inverse identities are algebraic and were proved in \hyperref[rec:PSC-I-T-EXACT-PACKET-COORDINATES]{\nolinkurl{PSC-I-T-EXACT-PACKET-COORDINATES}}.
\end{proof}
\field{Formal correspondence}{covered. The actual finite projectors are continuous and explicitly bounded as IsBoundedLinearMap over an arbitrary NormedField, including a trivial valuation. Individual bounds produce one finite uniform bound; analysis has that bound and synthesis has the finite-cardinality bound. Any coordinate norm equivalent to the canonical maximum norm is represented by a linear identification and its literal two-sided norm inequalities. Composition gives bounded actual analysis and synthesis, both inverse identities, and closed packet ranges. No conclusion-shaped hypotheses occur.}
\field{Source locator}{\nolinkurl{editable/traditional/03_packets_and_graphs.tex}; statement lines 267--270; proof lines 274--280.}
\field{Statement SHA-256}{\nolinkurl{3fd1f470c1a5dcd6273739c6f6cb95e5fbb8f8c95e4fc4c76e010430a72afa3e}}
\Needspace{12\baselineskip}\subsubsection*{57. Packet-coordinate condition number}
\phantomsection\label{rec:PSC-I-D-PACKET-CONDITION-NUMBER}
\field{Stable ID}{\nolinkurl{PSC-I-D-PACKET-CONDITION-NUMBER}}
\field{Kind}{definition}
\textbf{Statement.}\par
Fix a finite product norm on $\bigoplus_\alpha X_\alpha$ equivalent to the maximum norm, as in \hyperref[rec:PSC-I-T-BANACH-PACKET-ISOMORPHISM]{\nolinkurl{PSC-I-T-BANACH-PACKET-ISOMORPHISM}}. Define
\[
\kappa(\Pcal):=
\|\mathcal A_{\Pcal}\|\,
\|\mathcal S_{\Pcal}\|.
\]
\field{Formal correspondence}{definition represented. The definition is the product of the ordinary operator norms of the actual bounded analysis and synthesis. The generic bounded additive homomorphism norm is defined as the infimum of nonnegative valid operator bounds; linearity is retained by ofBoundedLinear. Equivalent-norm coordinates have explicit actual-map equalities.}
\field{Source locator}{\nolinkurl{editable/traditional/03_packets_and_graphs.tex}; statement lines 282--290.}
\field{Statement SHA-256}{\nolinkurl{7c67c989c536ab3eecad8fdf8c19b0ed6eef7c164c569d69a0e04de101b2460d}}
\Needspace{12\baselineskip}\subsubsection*{58. Conditioned coordinate transport}
\phantomsection\label{rec:PSC-I-P-CONDITIONED-COORDINATES}
\field{Stable ID}{\nolinkurl{PSC-I-P-CONDITIONED-COORDINATES}}
\field{Kind}{proposition}
\textbf{Statement.}\par
For a finite Banach packet system with $X\ne0$,
\[
\kappa(\Pcal)\geq1.
\]
If $T\in\mathcal B(X)$ and
\[
[T]_{\Pcal}:=\mathcal A_{\Pcal}T\mathcal S_{\Pcal},
\]
then
\[
\boxed{
\|[T]_{\Pcal}\|
\leq\kappa(\Pcal)\|T\|,
\qquad
\|T\|
\leq\kappa(\Pcal)\|[T]_{\Pcal}\|.
}
\]

\begin{proof}
Since $I=\mathcal S_{\Pcal}\mathcal A_{\Pcal}$,
\[
1=\|I\|\leq\|\mathcal S_{\Pcal}\|\|\mathcal A_{\Pcal}\|.
\]
The first displayed estimate follows from submultiplicativity.  The inverse identity
\[
T=\mathcal S_{\Pcal}[T]_{\Pcal}\mathcal A_{\Pcal}
\]
gives the second estimate.
\end{proof}
\field{Formal correspondence}{covered. For nonzero X the identity has norm one. The actual inverse relation S\ensuremath{\circ}A=id and operator norm submultiplicativity give \ensuremath{\kappa}\ensuremath{\ge}1. The coordinate operator is literally A\ensuremath{\circ}T\ensuremath{\circ}S, and S\ensuremath{\circ}[T]\ensuremath{\circ}A=T gives the reverse estimate. The final wrapper constructs both bounded actual coordinate maps for every equivalent finite-product norm and proves both inequalities for every bounded linear T over any NormedField.}
\field{Source locator}{\nolinkurl{editable/traditional/03_packets_and_graphs.tex}; statement lines 293--313; proof lines 317--327.}
\field{Statement SHA-256}{\nolinkurl{6799946dd2c88a65a635b418c7d0a74e94a9af4d1dc88df66cdd3b05c4ebc180}}
\Needspace{12\baselineskip}\subsubsection*{59. Completion preserves a finite packet split}
\phantomsection\label{rec:PSC-I-C-COMPLETION-PACKET-SPLIT}
\field{Stable ID}{\nolinkurl{PSC-I-C-COMPLETION-PACKET-SPLIT}}
\field{Kind}{corollary}
\textbf{Statement.}\par
Let $M$ be a dense normed subspace of a Banach space $X$.  Suppose every bounded $P_\alpha$ leaves $M$ invariant and the restricted projectors form a packet system on $M$.  Then
\[
P_\alpha X=\cl{P_\alpha M}
\]
and the packet decomposition of $M$ extends uniquely to the completion $X$.

\begin{proof}
For $x\in X$, choose $m_n\in M$ with $m_n\to x$.  Continuity gives $P_\alpha m_n\to P_\alpha x$, so $P_\alpha X\subseteq\cl{P_\alpha M}$.  The reverse inclusion follows because $P_\alpha X$ is closed and contains $P_\alpha M$.  The bounded extensions remain pairwise annihilating idempotents whose finite sum is the identity.
\end{proof}
\field{Formal correspondence}{covered. Over an arbitrary NormedField, the actual dense-submodule restrictions determine the ambient continuous projector identities by equality on a dense set. Each ambient range equals closure of the corresponding dense-submodule image, and two continuous ambient projector extensions agreeing on the dense submodule coincide. This is stronger than the bounded-input manuscript statement; boundedness of the given ambient projectors is available from the printed premise and is not needed for the density/uniqueness argument.}
\field{Source locator}{\nolinkurl{editable/traditional/03_packets_and_graphs.tex}; statement lines 334--341; proof lines 345--347.}
\field{Statement SHA-256}{\nolinkurl{67202410fc4580e865ea492dba6cf8dff7e0f296d89ac886c96188a0169a6b70}}
\Needspace{12\baselineskip}\subsubsection*{60. Unconditional packet sum space}
\phantomsection\label{rec:PSC-I-D-UNCONDITIONAL-PACKET-SPACE}
\field{Stable ID}{\nolinkurl{PSC-I-D-UNCONDITIONAL-PACKET-SPACE}}
\field{Kind}{definition}
\textbf{Statement.}\par
Let $(X_n)_{n\geq1}$ be closed subspaces of a real or complex Banach space $X$.  Define
\[
\bigoplus_{n\geq1}^{u}X_n
\]
to be the space of sequences $x=(x_n)$ with $x_n\in X_n$ for which $\sum_nx_n$ converges unconditionally in $X$.  Equip it with
\[
\|x\|_u
:=
\sup_{F\Subset\mathbb N}
\sup_{|a_n|\leq1}
\left\|\sum_{n\in F}a_nx_n\right\|.
\]
\field{Formal correspondence}{definition represented. The carrier is exactly sequences x\_n in the specified closed subspaces for which Summable x holds; mathlib Summable is the finite-subset net (unconditional summation). The norm is exactly sSup of norms of finite sums with coefficients in the closed scalar unit ball. Type synonyms avoid inherited product topology. Finiteness is derived, not assumed. RCLike covers real and complex scalars as stated.}
\field{Source locator}{\nolinkurl{editable/traditional/03_packets_and_graphs.tex}; statement lines 352--366.}
\field{Statement SHA-256}{\nolinkurl{bcbeba58302101658c66251d7b338a6a3932fb07b059e77931ffd063b5303b4f}}
\Needspace{12\baselineskip}\subsubsection*{61. Completeness of the unconditional packet sum}
\phantomsection\label{rec:PSC-I-L-UNCONDITIONAL-COMPLETENESS}
\field{Stable ID}{\nolinkurl{PSC-I-L-UNCONDITIONAL-COMPLETENESS}}
\field{Kind}{lemma}
\textbf{Statement.}\par
The normed space $\bigoplus_{n\geq1}^{u}X_n$ in \hyperref[rec:PSC-I-D-UNCONDITIONAL-PACKET-SPACE]{\nolinkurl{PSC-I-D-UNCONDITIONAL-PACKET-SPACE}} is Banach.

\begin{proof}
Let $x^{(k)}=(x_n^{(k)})$ be Cauchy in $\|\cdot\|_u$.  For every fixed $n$,
\[
\|x_n^{(k)}-x_n^{(\ell)}\|
\leq
\|x^{(k)}-x^{(\ell)}\|_u,
\]
so $x_n^{(k)}$ converges to some $x_n\in X_n$ because $X_n$ is closed.  Given $\varepsilon>0$, choose $k_0$ such that
\[
\|x^{(k)}-x^{(\ell)}\|_u<\varepsilon
\qquad(k,\ell\geq k_0).
\]
For finite $F$ and $|a_n|\leq1$, pass $\ell\to\infty$ in the finite sum to obtain
\[
\left\|\sum_{n\in F}a_n(x_n^{(k)}-x_n)\right\|
\leq\varepsilon
\qquad(k\geq k_0).
\]
Hence $\|x^{(k)}-x\|_u\leq\varepsilon$ once $x$ is known to be unconditionally summable.

Choose such a $k\geq k_0$.  Since $x^{(k)}$ is unconditionally summable, there is a finite $F_0$ such that every finite $F\subseteq\mathbb N\setminus F_0$ and every $|a_n|\leq1$ satisfy
\[
\left\|\sum_{n\in F}a_nx_n^{(k)}\right\|<\varepsilon.
\]
The preceding uniform estimate then gives
\[
\left\|\sum_{n\in F}a_nx_n\right\|<2\varepsilon.
\]
Thus the tails of $\sum_nx_n$ are uniformly small under bounded scalar changes, which is the unconditional Cauchy criterion.  Therefore $x\in\bigoplus^uX_n$, and the earlier estimate proves $x^{(k)}\to x$ in $\|\cdot\|_u$.
\end{proof}
\field{Formal correspondence}{covered. Uniform boundedness on the complete continuous dual is justified by finite-dimensional real absolute summability of each scalar image; the isometric bidual inclusion yields finite multiplier bounds. The signed norm controls each coordinate; coordinate limits stay in closed subspaces. Each fixed finite multiplier estimate passes to the limit, uniformly in the finite set and coefficients. Unweighted tails prove Summable for the limit, and the same uniform estimate gives convergence in the signed norm. Metric completeness is then obtained without a projector hypothesis.}
\field{Source locator}{\nolinkurl{editable/traditional/03_packets_and_graphs.tex}; statement lines 369--372; proof lines 376--405.}
\field{Statement SHA-256}{\nolinkurl{339adef1c7737cb8adc8a4bc76b3a5a2f95a575ffb576a34404d8da60bfe0b62}}
\Needspace{12\baselineskip}\subsubsection*{62. Countable unconditional packet isomorphism}
\phantomsection\label{rec:PSC-I-T-COUNTABLE-UNCONDITIONAL-PACKETS}
\field{Stable ID}{\nolinkurl{PSC-I-T-COUNTABLE-UNCONDITIONAL-PACKETS}}
\field{Kind}{theorem}
\textbf{Statement.}\par
Let $X$ be a real or complex Banach space and let $(P_n)_{n\geq1}\subseteq\mathcal B(X)$ satisfy
\[
P_nP_m=\delta_{nm}P_n
\]
and
\[
x=\sum_{n\geq1}P_nx
\]
unconditionally for every $x\in X$.  Assume
\[
K:=
\sup_{F\Subset\mathbb N}
\sup_{|a_n|\leq1}
\left\|\sum_{n\in F}a_nP_n\right\|<\infty.
\]
Then analysis and synthesis define inverse Banach space isomorphisms
\[
\boxed{
X\cong\bigoplus_{n\geq1}^{u}P_nX.
}
\]
More precisely,
\[
\|\mathcal A_{\Pcal}\|\leq K,
\qquad
\|\mathcal S_{\Pcal}\|\leq1.
\]

\begin{proof}
Each $P_nX=\Ker(I-P_n)$ is closed.  For $x\in X$,
\[
\|\mathcal A_{\Pcal}x\|_u
=
\sup_{F,a}
\left\|\sum_{n\in F}a_nP_nx\right\|
\leq K\|x\|.
\]
For $y=(y_n)\in\bigoplus^uP_nX$, define
\[
\mathcal S_{\Pcal}y=\sum_ny_n.
\]
Since the finite partial sums with every coefficient equal to one are among those used in $\|y\|_u$,
\[
\|\mathcal S_{\Pcal}y\|
\leq\|y\|_u.
\]
The assumed expansion gives $\mathcal S_{\Pcal}\mathcal A_{\Pcal}x=x$.  Conversely, boundedness of $P_m$ permits termwise application to the unconditionally convergent series:
\[
P_m\mathcal S_{\Pcal}y
=
\sum_nP_my_n
=y_m.
\]
Thus $\mathcal A_{\Pcal}\mathcal S_{\Pcal}=I$.  Completeness of the coordinate space is \hyperref[rec:PSC-I-L-UNCONDITIONAL-COMPLETENESS]{\nolinkurl{PSC-I-L-UNCONDITIONAL-COMPLETENESS}}.
\end{proof}
\field{Formal correspondence}{covered. ofSupBound constructs K as the literal supremum of finite multiplier operator norms from BddAbove. The carriers use Summable sequences in actual projector ranges. Continuous projectors commute with the sum, so annihilation proves both inverse identities. The literal signed norm yields analysis bound K and synthesis bound 1. A proved linear isometry identifies the projector-specific carrier and norm with the general ClosedPackets.Space. Completeness follows either from the general theorem or from the constructed uniform equivalence with Banach X.}
\field{Source locator}{\nolinkurl{editable/traditional/03_packets_and_graphs.tex}; statement lines 407--436; proof lines 441--467.}
\field{Statement SHA-256}{\nolinkurl{7579699c5a34468a6f9df136e7c897641c5409e5f5e597cd849348bddc92a469}}
\Needspace{12\baselineskip}\subsubsection*{63. Uniform boundedness}
\phantomsection\label{rec:PSC-I-R-UNIFORM-BOUNDEDNESS}
\field{Stable ID}{\nolinkurl{PSC-I-R-UNIFORM-BOUNDEDNESS}}
\field{Kind}{remark}
\textbf{Statement.}\par
If every expansion $\sum_nP_nx$ is unconditional, the operators obtained from finite partial sums with signs form a pointwise bounded family.  The uniform boundedness principle therefore yields a finite constant $K$.  We retain this constant explicitly in \hyperref[rec:PSC-I-T-COUNTABLE-UNCONDITIONAL-PACKETS]{\nolinkurl{PSC-I-T-COUNTABLE-UNCONDITIONAL-PACKETS}} because later estimates use it to bound the coordinate map.
\field{Formal correspondence}{documentary. A uniform-boundedness explanation. No instance for the signed packet partial-sum family is mapped.}
\field{Source locator}{\nolinkurl{editable/traditional/03_packets_and_graphs.tex}; statement lines 469--472.}
\field{Statement SHA-256}{\nolinkurl{2ae822d9b50496c6462b94007b5ff1a8d083e7109e6da7d0c2453ca40b80373c}}
\Needspace{12\baselineskip}\subsubsection*{64. Infinite packet sums}
\phantomsection\label{rec:PSC-I-B-INFINITE-PACKET-TOPOLOGY}
\field{Stable ID}{\nolinkurl{PSC-I-B-INFINITE-PACKET-TOPOLOGY}}
\field{Kind}{boundary}
\textbf{Statement.}\par
For an infinite family, an identity such as $\sum_\alpha P_\alpha=I$ requires a declared convergence mode and the corresponding uniform boundedness statement.  Pointwise, strong, norm, weak, and unconditional convergence are not interchangeable.
\field{Formal correspondence}{scope boundary. Convergence modes remain separate; finite identities are not used as infinite sums.}
\field{Source locator}{\nolinkurl{editable/traditional/03_packets_and_graphs.tex}; statement lines 475--478.}
\field{Statement SHA-256}{\nolinkurl{305cbc8493dffcc2407698c4df7a3f54fecc76bb57e10c75390b0cedc5bce53d}}
\Needspace{12\baselineskip}\subsubsection*{65. Transfer block}
\phantomsection\label{rec:PSC-I-D-TRANSFER-BLOCK}
\field{Stable ID}{\nolinkurl{PSC-I-D-TRANSFER-BLOCK}}
\field{Kind}{definition}
\textbf{Statement.}\par
For $T\in\End_{\bbk}(X)$ and a finite packet system $\Pcal$, define
\[
T_{\beta\alpha}:=P_\beta TP_\alpha
\in\Hom_{\bbk}(X_\alpha,X_\beta).
\]
\field{Formal correspondence}{definition represented. Ambient compositions and their genuinely range-restricted packet Hom maps are both encoded.}
\field{Source locator}{\nolinkurl{editable/traditional/03_packets_and_graphs.tex}; statement lines 492--499.}
\field{Statement SHA-256}{\nolinkurl{1b27b671e8c43628db2f3e157df05e01a63658c32e64a6ae2a94fdc69fd5ad73}}
\Needspace{12\baselineskip}\subsubsection*{66. Exact operator reconstruction}
\phantomsection\label{rec:PSC-I-T-OPERATOR-RECONSTRUCTION}
\field{Stable ID}{\nolinkurl{PSC-I-T-OPERATOR-RECONSTRUCTION}}
\field{Kind}{theorem}
\textbf{Statement.}\par
Every $T\in\End_{\bbk}(X)$ satisfies
\[
\boxed{T=\sum_{\alpha,\beta\in I}T_{\beta\alpha}.}
\]

\begin{proof}
Insert the finite packet identity on both sides:
\[
T=
\left(\sum_\beta P_\beta\right)
T
\left(\sum_\alpha P_\alpha\right)
=
\sum_{\alpha,\beta}P_\beta TP_\alpha.
\]
\end{proof}
\field{Formal correspondence}{covered. The existing result is more general: source and target can have different finite packet systems.}
\field{Source locator}{\nolinkurl{editable/traditional/03_packets_and_graphs.tex}; statement lines 502--508; proof lines 512--522.}
\field{Statement SHA-256}{\nolinkurl{f8c54ea5f0aa3075a42556f7d9b9b29fd9399d88e9de974c6eb098fceb2a7abb}}
\Needspace{12\baselineskip}\subsubsection*{67. Packet matrix multiplication}
\phantomsection\label{rec:PSC-I-T-BLOCK-MULTIPLICATION}
\field{Stable ID}{\nolinkurl{PSC-I-T-BLOCK-MULTIPLICATION}}
\field{Kind}{theorem}
\textbf{Statement.}\par
For $S,T\in\End_{\bbk}(X)$,
\[
\boxed{
(ST)_{\gamma\alpha}
=
\sum_{\beta\in I}S_{\gamma\beta}T_{\beta\alpha}.
}
\]

\begin{proof}
Insert $I=\sum_\beta P_\beta$ between $S$ and $T$:
\[
P_\gamma STP_\alpha
=
\sum_\beta
P_\gamma SP_\beta TP_\alpha.
\]
\end{proof}
\field{Formal correspondence}{covered. The output-first block product is proved for three possibly different modules.}
\field{Source locator}{\nolinkurl{editable/traditional/03_packets_and_graphs.tex}; statement lines 529--539; proof lines 544--552.}
\field{Statement SHA-256}{\nolinkurl{a7aaf1cde8454ca239c380ca964f5b13f20d39baaf62b7a859caf9a849378314}}
\Needspace{12\baselineskip}\subsubsection*{68. Packet powers and paths}
\phantomsection\label{rec:PSC-I-C-BLOCK-POWERS}
\field{Stable ID}{\nolinkurl{PSC-I-C-BLOCK-POWERS}}
\field{Kind}{corollary}
\textbf{Statement.}\par
For $m\geq1$,
\[
(T^m)_{\beta_m\beta_0}
=
\sum_{\beta_1,\ldots,\beta_{m-1}}
T_{\beta_m\beta_{m-1}}
\cdots
T_{\beta_1\beta_0}.
\]
The same formula holds for a product of nonidentical operators with the corresponding blocks in temporal order.

\begin{proof}
Induct on $m$ using \hyperref[rec:PSC-I-T-BLOCK-MULTIPLICATION]{\nolinkurl{PSC-I-T-BLOCK-MULTIPLICATION}}.  For a product of different operators, use the same induction with the specified operator at each step.
\end{proof}
\field{Formal correspondence}{covered. The power theorem expands transferBlock(P,P,T\textasciicircum{}m) using packetPathSum. The independent operatorWord/packetWordSum induction expands every finite list of possibly different operators, with composition order matching rightmost-acts-first temporal order. Each summation is over actual intermediate packet labels.}
\field{Source locator}{\nolinkurl{editable/traditional/03_packets_and_graphs.tex}; statement lines 559--571; proof lines 575--577.}
\field{Statement SHA-256}{\nolinkurl{dbaa3f1f55517bd6366bcf856dc8eb6836a642144d805f9bed4899e5bba7038f}}
\Needspace{12\baselineskip}\subsubsection*{69. Adjoint reversal}
\phantomsection\label{rec:PSC-I-C-ADJOINT-REVERSAL}
\field{Stable ID}{\nolinkurl{PSC-I-C-ADJOINT-REVERSAL}}
\field{Kind}{corollary}
\textbf{Statement.}\par
If $X$ is Hilbert, every $P_\alpha$ is orthogonal, and $T\in\mathcal B(X)$, then
\[
\boxed{(T_{\beta\alpha})^*=(T^*)_{\alpha\beta}.}
\]

\begin{proof}
Because $P_\alpha^*=P_\alpha$,
\[
(P_\beta TP_\alpha)^*=P_\alpha T^*P_\beta.
\]
\end{proof}
\field{Formal correspondence}{covered. The result is proved for bounded self-adjoint outer maps, more generally than orthogonal packet projectors.}
\field{Source locator}{\nolinkurl{editable/traditional/03_packets_and_graphs.tex}; statement lines 579--585; proof lines 589--594.}
\field{Statement SHA-256}{\nolinkurl{8e978a295a29a3eac3d1c4bec1a8e3f4e05a2684fdbce8e80a357aa7891e8f79}}
\Needspace{12\baselineskip}\subsubsection*{70. Tensor-product transfer blocks}
\phantomsection\label{rec:PSC-I-C-TENSOR-TRANSFER-BLOCKS}
\field{Stable ID}{\nolinkurl{PSC-I-C-TENSOR-TRANSFER-BLOCKS}}
\field{Kind}{corollary}
\textbf{Statement.}\par
Let \(S\in\End(X)\) and \(T\in\End(Y)\), with packet systems \(\Pcal\) and \(\Qcal\).  Relative to the tensor product packet system of \hyperref[rec:PSC-I-T-TENSOR-PACKET-SYSTEM]{\nolinkurl{PSC-I-T-TENSOR-PACKET-SYSTEM}},
\[
\boxed{
(S\otimes T)_{(\beta,\delta),(\alpha,\gamma)}
=
S_{\beta\alpha}\otimes T_{\delta\gamma}.
}
\]

\begin{proof}
By definition,
\begin{align*}
(S\otimes T)_{(\beta,\delta),(\alpha,\gamma)}
&=(P_\beta\otimes Q_\delta)(S\otimes T)(P_\alpha\otimes Q_\gamma)\\
&=(P_\beta SP_\alpha)\otimes(Q_\delta TQ_\gamma)\\
&=S_{\beta\alpha}\otimes T_{\delta\gamma}.
\end{align*}
\end{proof}
\field{Formal correspondence}{covered. Actual transfer blocks factor by tensor-map composition in exactly the (target,source) order printed. Both sides are actual endomorphisms of X tensor Y; no factorization witness is an input.}
\field{Source locator}{\nolinkurl{editable/traditional/03_packets_and_graphs.tex}; statement lines 598--608; proof lines 613--621.}
\field{Statement SHA-256}{\nolinkurl{ec0bc2d2c85cdf55ea2f27003457a653f8f2426fe97dafd6c918cbda54874878}}
\Needspace{12\baselineskip}\subsubsection*{71. Blockwise intertwining equations}
\phantomsection\label{rec:PSC-I-T-BLOCK-INTERTWINING}
\field{Stable ID}{\nolinkurl{PSC-I-T-BLOCK-INTERTWINING}}
\field{Kind}{theorem}
\textbf{Statement.}\par
Let \(T\in\End(X)\), \(S\in\End(Y)\), let \(\Pcal=(P_\alpha)_{\alpha\in I}\) and \(\Qcal=(Q_\gamma)_{\gamma\in J}\) be finite packet systems, and let \(U:X\to Y\) be linear.  Define
\[
T_{\beta\alpha}=P_\beta TP_\alpha,
\qquad
S_{\gamma\delta}=Q_\gamma SQ_\delta,
\qquad
U_{\gamma\alpha}=Q_\gamma UP_\alpha.
\]
Then
\[
\boxed{UT=SU}
\]
if and only if, for every \(\gamma\in J\) and \(\alpha\in I\),
\[
\boxed{
\sum_{\beta\in I}U_{\gamma\beta}T_{\beta\alpha}
=
\sum_{\delta\in J}S_{\gamma\delta}U_{\delta\alpha}.
}
\]

\begin{proof}
Apply \(Q_\gamma\) on the left and \(P_\alpha\) on the right.  Inserting \(I_X=\sum_\beta P_\beta\) gives
\[
Q_\gamma UTP_\alpha
=
\sum_\beta Q_\gamma UP_\beta TP_\alpha
=
\sum_\beta U_{\gamma\beta}T_{\beta\alpha}.
\]
Inserting \(I_Y=\sum_\delta Q_\delta\) gives
\[
Q_\gamma SUP_\alpha
=
\sum_\delta Q_\gamma SQ_\delta UP_\alpha
=
\sum_\delta S_{\gamma\delta}U_{\delta\alpha}.
\]
Thus \(UT=SU\) implies every block equation.  Conversely, if every block equation holds, then all blocks of \(UT-SU\) vanish.  Exact operator reconstruction yields \(UT-SU=0\).
\end{proof}
\field{Formal correspondence}{covered. U is any linear map, not an intertwiner structure. Forward direction inserts middle projector sums; reverse direction reconstructs UT and SU from all equal blocks. Distinct source/target index sets are retained.}
\field{Source locator}{\nolinkurl{editable/traditional/03_packets_and_graphs.tex}; statement lines 623--645; proof lines 650--668.}
\field{Statement SHA-256}{\nolinkurl{558adbc4eead21c3f371918f40b9f462000c7b26a7234864ee44c103e09cf1e6}}
\Needspace{12\baselineskip}\subsubsection*{72. Full block-algebra representation}
\phantomsection\label{rec:PSC-I-T-BLOCK-ALGEBRA}
\field{Stable ID}{\nolinkurl{PSC-I-T-BLOCK-ALGEBRA}}
\field{Kind}{theorem}
\textbf{Statement.}\par
The map
\[
\Theta_{\Pcal}:\End_{\bbk}(X)
\longrightarrow
\prod_{\beta,\alpha\in I}\Hom_{\bbk}(X_\alpha,X_\beta),
\qquad
T\longmapsto(T_{\beta\alpha})_{\beta,\alpha},
\]
is an algebra isomorphism when the codomain is equipped with matrix multiplication whose entries are operators.

\begin{proof}
Injectivity follows from exact reconstruction.  Given arbitrary linear maps $B_{\beta\alpha}:X_\alpha\to X_\beta$, define
\[
T\left(\sum_\alpha x_\alpha\right)
:=
\sum_{\alpha,\beta}B_{\beta\alpha}x_\alpha.
\]
The packet decomposition is unique and the index set is finite, so this defines a linear operator whose blocks are the prescribed maps.  Thus $\Theta_{\Pcal}$ is surjective.  Compatibility with multiplication is \hyperref[rec:PSC-I-T-BLOCK-MULTIPLICATION]{\nolinkurl{PSC-I-T-BLOCK-MULTIPLICATION}}.
\end{proof}
\field{Formal correspondence}{covered. Full dependent packet Hom families and their multiplication are identified with the actual endomorphism algebra.}
\field{Source locator}{\nolinkurl{editable/traditional/03_packets_and_graphs.tex}; statement lines 675--686; proof lines 691--699.}
\field{Statement SHA-256}{\nolinkurl{a66e4216289ceb5e48d1ccd41d269b656c6c93ed462a47f68cbd0172c9360bda}}
\Needspace{12\baselineskip}\subsubsection*{73. Propagation of packet support}
\phantomsection\label{rec:PSC-I-C-PACKET-SUPPORT-PROPAGATION}
\field{Stable ID}{\nolinkurl{PSC-I-C-PACKET-SUPPORT-PROPAGATION}}
\field{Kind}{corollary}
\textbf{Statement.}\par
Let \(T\in\End(X)\) and \(x\in X\).  Then
\[
\boxed{
\Supp_{\Pcal}(Tx)
\subseteq
\left\{\beta:\exists\alpha\in\Supp_{\Pcal}(x)
\text{ with }T_{\beta\alpha}\neq0\right\}.
}
\]
Equality holds whenever, for every target \(\beta\) in the set on the right, the active contributions
\[
\bigl(T_{\beta\alpha}P_\alpha x:
\alpha\in\Supp_{\Pcal}(x),\ T_{\beta\alpha}\neq0\bigr)
\]
have nonzero sum.  In particular, equality holds if at most one active source packet feeds each such \(\beta\) and that source coordinate does not lie in the kernel of the corresponding block.

\begin{proof}
For every target packet \(\beta\),
\[
P_\beta Tx
=
\sum_{\alpha\in\Supp_{\Pcal}(x)}T_{\beta\alpha}P_\alpha x.
\]
If every block in this sum is zero, then the coordinate with index \(\beta\) vanishes.  This proves the inclusion.  Equality is exactly the assertion that the displayed active sum is nonzero for every target in the set on the right.  The final sufficient condition leaves at most one nonzero summand at each target, and kernel avoidance makes that summand nonzero.
\end{proof}
\field{Formal correspondence}{covered. The source expands each target coordinate as the finite sum of the actual transfer-block contributions and filters to active sources. This yields the support inclusion; nonzero active sums give equality. The single-active theorem uses activeSources=\{a\} plus nonzero transfer contribution, which is equivalent to at-most-one active source plus kernel avoidance on a target that is active.}
\field{Source locator}{\nolinkurl{editable/traditional/03_packets_and_graphs.tex}; statement lines 703--720; proof lines 725--733.}
\field{Statement SHA-256}{\nolinkurl{f05b8bf4828b88f6020e2e0b824e228acffcb77d7a514ced31ec90c64e6d4437}}
\Needspace{12\baselineskip}\subsubsection*{74. Dimension count}
\phantomsection\label{rec:PSC-I-C-DIMENSION-COUNT}
\field{Stable ID}{\nolinkurl{PSC-I-C-DIMENSION-COUNT}}
\field{Kind}{corollary}
\textbf{Statement.}\par
If $X$ is finite dimensional and $d_\alpha=\dim X_\alpha$, then
\[
\dim\End_{\bbk}(X)
=
\sum_{\alpha,\beta}d_\alpha d_\beta
=
\left(\sum_\alpha d_\alpha\right)^2.
\]

\begin{proof}
Take dimensions in the block algebra isomorphism and use $\dim\Hom(X_\alpha,X_\beta)=d_\alpha d_\beta$.
\end{proof}
\field{Formal correspondence}{covered. The actual End space is linearly equivalent to the heterogeneous packet Hom product. Its finite dimension equals the sum of products of actual packet-range dimensions, and finite distributivity gives the displayed square. Both printed equalities occur in one conjunction.}
\field{Source locator}{\nolinkurl{editable/traditional/03_packets_and_graphs.tex}; statement lines 736--746; proof lines 750--752.}
\field{Statement SHA-256}{\nolinkurl{85dd27ca5c20eeb56683b4b3d9b9093503e2c3c2fdc4a07c1aaf33eefab855c8}}
\Needspace{12\baselineskip}\subsubsection*{75. Block-diagonal projection and packet commutant}
\phantomsection\label{rec:PSC-I-P-BLOCK-DIAGONAL-COMMUTANT}
\field{Stable ID}{\nolinkurl{PSC-I-P-BLOCK-DIAGONAL-COMMUTANT}}
\field{Kind}{proposition}
\textbf{Statement.}\par
Define
\[
\Delta_{\Pcal}(T):=\sum_{\alpha}P_\alpha TP_\alpha.
\]
Then $\Delta_{\Pcal}$ is a linear idempotent on $\End(X)$ whose range is the block diagonal subalgebra.  For $T\in\End(X)$, the following are equivalent:
\begin{enumerate}[label=(\roman*)]
\item $TP_\alpha=P_\alpha T$ for every $\alpha$.
\item $T_{\beta\alpha}=0$ whenever $\beta\neq\alpha$.
\item $\Delta_{\Pcal}(T)=T$.
\end{enumerate}

\begin{proof}
Pairwise annihilation gives $\Delta_{\Pcal}^2=\Delta_{\Pcal}$ and identifies its range with diagonal block families.  If $T$ commutes with all $P_\alpha$, then for $\beta\neq\alpha$,
\[
P_\beta TP_\alpha=P_\beta P_\alpha T=0.
\]
Thus (i) implies (ii).  Exact reconstruction shows (ii) implies (iii).  If $T=\sum_\alpha P_\alpha TP_\alpha$, then
\[
TP_\gamma=P_\gamma TP_\gamma=P_\gamma T,
\]
so (iii) implies (i).
\end{proof}
\field{Formal correspondence}{covered. diagonalProjection is the actual K-linear map T\ensuremath{\mapsto}\ensuremath{\Sigma}\_i P\_i T P\_i. The proofs derive idempotence, equivalence of commutation with every projector and off-diagonal block vanishing, equivalence to fixed points of this projection, and equality of its actual linear range to the equality-supported subalgebra.}
\field{Source locator}{\nolinkurl{editable/traditional/03_packets_and_graphs.tex}; statement lines 754--766; proof lines 771--781.}
\field{Statement SHA-256}{\nolinkurl{613238e060365a7c00244f44abb552ada9f54933d8483382f259d3f32843eae5}}
\Needspace{12\baselineskip}\subsubsection*{76. Support relation and support space}
\phantomsection\label{rec:PSC-I-D-SUPPORT-RELATION}
\field{Stable ID}{\nolinkurl{PSC-I-D-SUPPORT-RELATION}}
\field{Kind}{definition}
\textbf{Statement.}\par
For a relation $R\subseteq I\times I$, with $(\beta,\alpha)$ interpreted as an allowed transfer $\alpha\to\beta$, define
\[
\mathfrak M_R
:=
\{T\in\End(X):T_{\beta\alpha}=0\text{ whenever }(\beta,\alpha)\notin R\}.
\]
\field{Formal correspondence}{definition represented. OperatorSupported is initially defined via actual restricted packet maps; operatorSupported\_iff\_ambient proves its exact equivalence to ambient transfer-block vanishing. supportSubmodule is the corresponding actual linear subspace of End(X). Relation dictionary: formal rel(a,b) means manuscript (b,a)\ensuremath{\in}R, with a source and b target.}
\field{Source locator}{\nolinkurl{editable/traditional/03_packets_and_graphs.tex}; statement lines 789--797.}
\field{Statement SHA-256}{\nolinkurl{babc2ddcee7669992338d1b7741bd50fab01175a280772ab40df14dcc814ebf3}}
\Needspace{12\baselineskip}\subsubsection*{77. Support-pattern algebra criterion}
\phantomsection\label{rec:PSC-I-T-SUPPORT-PATTERN-ALGEBRA}
\field{Stable ID}{\nolinkurl{PSC-I-T-SUPPORT-PATTERN-ALGEBRA}}
\field{Kind}{theorem}
\textbf{Statement.}\par
Assume $X_\alpha\neq0$ for every $\alpha\in I$.  Then:
\begin{enumerate}[label=(\roman*)]
\item $I_X\in\mathfrak M_R$ if and only if $R$ is reflexive.
\item $\mathfrak M_R\mathfrak M_R\subseteq\mathfrak M_R$ if and only if $R$ is transitive.
\item $\mathfrak M_R$ is a unital subalgebra if and only if $R$ is a preorder.
\end{enumerate}

\begin{proof}
The identity has diagonal blocks $I_{X_\alpha}$ and no blocks outside the diagonal, so it belongs to $\mathfrak M_R$ exactly when $(\alpha,\alpha)\in R$ for every $\alpha$.

Assume $R$ is transitive.  If a summand
\[
S_{\gamma\beta}T_{\beta\alpha}
\]
in the $(\gamma,\alpha)$ block of $ST$ is nonzero, then $(\gamma,\beta)$ and $(\beta,\alpha)$ lie in $R$.  Transitivity gives $(\gamma,\alpha)\in R$, so every forbidden block of $ST$ vanishes.

Conversely, suppose transitivity fails.  Choose $\alpha,\beta,\gamma$ with
\[
(\beta,\alpha)\in R,
\qquad
(\gamma,\beta)\in R,
\qquad
(\gamma,\alpha)\notin R.
\]
Choose nonzero vectors $u\in X_\alpha$, $v\in X_\beta$, and $w\in X_\gamma$.  Choose linear functionals $f\in X_\alpha^*$ and $g\in X_\beta^*$ with $f(u)=g(v)=1$.  Let $T$ have the sole block $x\mapsto f(x)v$ from $X_\alpha$ to $X_\beta$, and let $S$ have the sole block $y\mapsto g(y)w$ from $X_\beta$ to $X_\gamma$.  Then $S,T\in\mathfrak M_R$, but $ST$ has a nonzero forbidden $(\gamma,\alpha)$ block.  This proves necessity.  Combining reflexivity and transitivity gives (iii).
\end{proof}
\field{Formal correspondence}{covered. The two difficult converses use arbitrary maps between the actual nonzero packet ranges. Identity support is equivalent to reflexivity; multiplication closure is equivalent to transitivity. support\_mul\_closed\_iff transfers the block statement to actual ambient operators, and support\_pattern\_algebra\_criterion characterizes existence of a Subalgebra whose carrier is precisely that support space.}
\field{Source locator}{\nolinkurl{editable/traditional/03_packets_and_graphs.tex}; statement lines 800--808; proof lines 812--830.}
\field{Statement SHA-256}{\nolinkurl{fcd8af3b6718c7460c4e58639b8e942c24ede23b8673f9811f193cab4ccd3b20}}
\Needspace{12\baselineskip}\subsubsection*{78. \$*\$-closed bounded support patterns}
\phantomsection\label{rec:PSC-I-C-SUPPORT-PATTERN-STAR}
\field{Stable ID}{\nolinkurl{PSC-I-C-SUPPORT-PATTERN-STAR}}
\field{Kind}{corollary}
\textbf{Statement.}\par
Let $X$ be Hilbert, let the packet projectors be orthogonal, assume $X_\alpha\neq0$ for every $\alpha\in I$, and put
\[
\mathfrak M_R^{\mathrm b}:=\mathfrak M_R\cap\mathcal B(X).
\]
Then $\mathfrak M_R^{\mathrm b}$ is closed under adjoints if and only if $R$ is symmetric.  Hence $\mathfrak M_R^{\mathrm b}$ is a unital $*$ subalgebra of $\mathcal B(X)$ if and only if $R$ is an equivalence relation.

\begin{proof}
Adjoint reversal sends an allowed $(\beta,\alpha)$ block to an $(\alpha,\beta)$ block.  Symmetry is therefore sufficient.  If symmetry fails, choose a nonzero map of rank one in an allowed nonsymmetric block.  Its adjoint occupies the forbidden reverse block, proving necessity.  A reflexive, transitive, symmetric relation is an equivalence relation.
\end{proof}
\field{Formal correspondence}{covered. Supported uses literal continuous transfer blocks and is proved equivalent to the underlying algebraic support predicate. Converse symmetry and transitivity use actual bounded rank-one operators between nonzero ranges. Identity support forces reflexivity. Sufficiency constructs an actual StarSubalgebra. RCLike covers the real and complex Hilbert settings; no finite-dimensional restriction or blanket bounded-block assumption is used.}
\field{Source locator}{\nolinkurl{editable/traditional/03_packets_and_graphs.tex}; statement lines 832--839; proof lines 843--845.}
\field{Statement SHA-256}{\nolinkurl{89dfdc356ea3dbbf9ee8018c66d8307f445d19078e0d7226372aae44c1872333}}
\Needspace{12\baselineskip}\subsubsection*{79. Preorder support algebra decomposition}
\phantomsection\label{rec:PSC-I-T-PREORDER-ALGEBRA-DECOMPOSITION}
\field{Stable ID}{\nolinkurl{PSC-I-T-PREORDER-ALGEBRA-DECOMPOSITION}}
\field{Kind}{theorem}
\textbf{Statement.}\par
Let $R$ be a preorder on the finite set $I$.  Define
\[
\alpha\sim\beta
\quad\Longleftrightarrow\quad
(\beta,\alpha)\in R\text{ and }(\alpha,\beta)\in R.
\]
Let $\overline I=I/{\sim}$, which carries the induced partial order, and put
\[
X_C:=\bigoplus_{\alpha\in C}X_\alpha.
\]
Let $\mathfrak N_R\subseteq\mathfrak M_R$ consist of operators whose blocks vanish whenever source and target belong to the same equivalence class.  Then:
\begin{enumerate}[label=(\roman*)]
\item $\mathfrak N_R$ is a nilpotent ideal on both sides.
\item If $h$ is the maximal number of strict steps in a chain of $\overline I$, then
\[
\mathfrak N_R^{\,h+1}=0.
\]
\item The map to the diagonal classes induces an algebra isomorphism
\[
\boxed{
\mathfrak M_R/\mathfrak N_R
\cong
\prod_{C\in\overline I}\End_{\bbk}(X_C).
}
\]
\end{enumerate}

\begin{proof}
A nonzero block in $\mathfrak N_R$ moves from one equivalence class to a strictly larger class in the quotient partial order.  The block multiplication law shows that products concatenate such moves.  Therefore left or right multiplication by an element of $\mathfrak M_R$ preserves the property of having no blocks within a single class, so $\mathfrak N_R$ is an ideal on both sides.  A product of $h+1$ elements of $\mathfrak N_R$ would require a strict chain of $h+1$ steps in $\overline I$, which does not exist.  This proves nilpotence.

Let $P_C=\sum_{\alpha\in C}P_\alpha$.  Define
\[
\rho:\mathfrak M_R\longrightarrow\prod_C\End(X_C),
\qquad
\rho(T)=(P_CTP_C)_C.
\]
If an intermediate class $D$ contributes to $P_CSTP_C$, then the preorder contains transfers $C\to D$ and $D\to C$, forcing $C=D$ in the quotient poset.  Hence
\[
P_CSTP_C=(P_CSP_C)(P_CTP_C),
\]
so $\rho$ is an algebra homomorphism.  Every collection of endomorphisms of the $X_C$ defines an allowed operator that is block diagonal with respect to the classes, so $\rho$ is surjective.  Its kernel is exactly $\mathfrak N_R$.  The first isomorphism theorem gives the quotient formula.
\end{proof}
\field{Formal correspondence}{covered. The support algebra is constructed from the actual preorder. strictSupportIdeal is an actual TwoSidedIdeal of it, with membership proved equivalent to vanishing within every mutual-order class. Nonzero product blocks force strict paths; finite strict-path height is bounded and attained on nonempty indices. strictSupportIdeal\_pow\_eq\_bot uses Mathlib standard ideal powers, with the span-of-words surrogate bridged by strictIdeal\_power\_mem. Antisymmetrization supplies the actual quotient partial order; class\_maximalStrictSteps proves its height equals the exponent parameter. Grouped class projectors are actual sums of original projectors. The diagonal-class homomorphism is surjective with kernel equal to the strict ideal, and preorderQuotientDirectSumEquiv is the actual quotient by that ideal to the product of End of the literal direct sums of class packet ranges.}
\field{Source locator}{\nolinkurl{editable/traditional/03_packets_and_graphs.tex}; statement lines 847--875; proof lines 880--894.}
\field{Statement SHA-256}{\nolinkurl{89f93d5bdea5558d5ca9fef72742db6339af268623a37bedb7baabb777e4452a}}
\Needspace{12\baselineskip}\subsubsection*{80. Transfer graph of an algebra-generating family}
\phantomsection\label{rec:PSC-I-D-TRANSFER-GRAPH}
\field{Stable ID}{\nolinkurl{PSC-I-D-TRANSFER-GRAPH}}
\field{Kind}{definition}
\textbf{Statement.}\par
Let $\Pcal=(P_\alpha)_{\alpha\in I}$ be a finite packet system on $X$ with every $P_\alpha\in\A$.  Let $\mathcal G\subseteq\A$ satisfy
\[
\A=\alg_{\bbk}\bigl(\{P_\alpha:\alpha\in I\}\cup\mathcal G\bigr),
\]
where $\alg_{\bbk}(E)$ denotes the unital subalgebra over $\bbk$ generated by $E$.  The transfer graph $\Gamma(\mathcal G,\Pcal)$ has vertex set $I$ and an arrow
\[
\alpha\longrightarrow\beta
\]
when
\[
P_\beta GP_\alpha\neq0
\]
for at least one $G\in\mathcal G$.  Reachability includes paths of length zero.
\field{Formal correspondence}{definition represented. packetGeneratedAlgebra is Algebra.adjoin K (range P.proj \ensuremath{\cup} ops), exactly the declared algebra in the manuscript. TransferEdge is witnessed by an actual member of ops with nonzero ambient transfer block. ReachableFrom uses ReflTransGen and hence includes length zero.}
\field{Source locator}{\nolinkurl{editable/traditional/03_packets_and_graphs.tex}; statement lines 916--931.}
\field{Statement SHA-256}{\nolinkurl{f6ddc4edec01afcf027a48300c34e66e9d1209d4f56bce17e3b9d00f88cd62cf}}
\Needspace{12\baselineskip}\subsubsection*{81. Reachability support of the generated algebra}
\phantomsection\label{rec:PSC-I-T-ALGEBRA-REACHABILITY-SUPPORT}
\field{Stable ID}{\nolinkurl{PSC-I-T-ALGEBRA-REACHABILITY-SUPPORT}}
\field{Kind}{theorem}
\textbf{Statement.}\par
For every $A\in\A$,
\[
P_\beta AP_\alpha=0
\qquad\text{whenever }\beta\notin\Reach(\alpha).
\]
Equivalently,
\[
\A\subseteq\mathfrak M_{R_\Gamma},
\qquad
R_\Gamma:=\{(\beta,\alpha):\beta\in\Reach(\alpha)\}.
\]

\begin{proof}
It is enough to treat a word in the generators and packet projectors.  Insert
\[
I=\sum_{\gamma\in I}P_\gamma
\]
between consecutive nonprojector generators.  Every resulting $(\beta,\alpha)$ block is a finite sum of products
\[
P_\beta G_mP_{\gamma_{m-1}}
G_{m-1}\cdots
P_{\gamma_1}G_1P_\alpha.
\]
A nonzero product requires every consecutive block to be nonzero, hence gives a directed path
\[
\alpha\to\gamma_1\to\cdots\to\gamma_{m-1}\to\beta.
\]
Packet projectors only insert stationary steps.  Thus a word has no block from $\alpha$ to an unreachable $\beta$.  Linear combinations preserve the same zero pattern.
\end{proof}
\field{Formal correspondence}{covered. Each projector and declared generator lies in the actual support algebra of ReflTransGen(TransferEdge). Algebra.adjoin\_le therefore proves containment of the entire generated algebra. generatedAlgebra\_block\_zero is the literal ambient P\_b A P\_a=0 consequence outside reachability.}
\field{Source locator}{\nolinkurl{editable/traditional/03_packets_and_graphs.tex}; statement lines 934--947; proof lines 951--967.}
\field{Statement SHA-256}{\nolinkurl{7496f96565d405f8fe6b80491655dffe3b2ddad07b104c1131704464fb8ad031}}
\Needspace{12\baselineskip}\subsubsection*{82. Reachability obstruction for generated hulls}
\phantomsection\label{rec:PSC-I-P-REACHABILITY-OBSTRUCTION}
\field{Stable ID}{\nolinkurl{PSC-I-P-REACHABILITY-OBSTRUCTION}}
\field{Kind}{proposition}
\textbf{Statement.}\par
Let $I_0\subseteq I$ and let
\[
S\subseteq\bigoplus_{\alpha\in I_0}X_\alpha.
\]
Then
\[
\boxed{
\Hull_{\A}(S)
\subseteq
\bigoplus_{\beta\in\Reach(I_0)}X_\beta.
}
\]
In particular, no seed supported outside a source strongly connected component can generate a nonzero vector in that source component.

\begin{proof}
For $s\in S$, write $s=\sum_{\alpha\in I_0}P_\alpha s$.  If $\beta$ is unreachable from $I_0$, then for every $A\in\A$,
\[
P_\beta As
=
\sum_{\alpha\in I_0}P_\beta AP_\alpha s
=0
\]
by \hyperref[rec:PSC-I-T-ALGEBRA-REACHABILITY-SUPPORT]{\nolinkurl{PSC-I-T-ALGEBRA-REACHABILITY-SUPPORT}}.  The same holds for every finite linear combination of orbit vectors.
\end{proof}
\field{Formal correspondence}{covered. The actual reachable-support subspace is preserved by every element of the generated algebra. Span of actual algebra orbits yields the hull inclusion for arbitrary seed sets. The source SCC is the actual mutual-reachability quotient, with source status equivalent to no incoming condensation edge. generatedAlgebra\_source\_obstruction separately assembles the printed final clause: seeds outside that source cannot generate any nonzero vector supported in it. It obtains a nonzero coordinate by packet completeness and applies no\_reach\_into\_source and the unreachable-coordinate obstruction.}
\field{Source locator}{\nolinkurl{editable/traditional/03_packets_and_graphs.tex}; statement lines 977--992; proof lines 996--1005.}
\field{Statement SHA-256}{\nolinkurl{a46f6c292b4825e63689b6302077e4db2b2a8fe38f9958fc9018a623e7b8a239}}
\Needspace{12\baselineskip}\subsubsection*{83. Packet-pure seed rank}
\phantomsection\label{rec:PSC-I-D-DIRECTED-SEED-RANK}
\field{Stable ID}{\nolinkurl{PSC-I-D-DIRECTED-SEED-RANK}}
\field{Kind}{definition}
\textbf{Statement.}\par
A vector is packet pure when it is zero or lies in one packet space $X_\alpha$, equivalently when its packet support has cardinality at most one.  Define
\[
\psrank_{\A,\Pcal}(X)
:=
\min\{|S|:S\text{ is packet-pure and }\A S=X\}.
\]
For $C\subseteq I$, put
\[
P_C:=\sum_{\alpha\in C}P_\alpha,
\qquad
X_C:=P_CX,
\qquad
\A_C:=P_C\A P_C,
\]
and define the local packet pure rank
\[
r_C:=\psrank_{\A_C,\Pcal|_C}(X_C).
\]
\field{Formal correspondence}{definition represented. The cardinal infimum is over all actual pure generating subsets, its indexing type is nonempty, and its infimum is attained by well-ordering. Corner compression is e*T on range(e), hence literally e*T*e there. Closure and the identity restriction are proved. A local family retains outside indices as zero maps; local\_proj\_zero\_iff identifies these precisely for nonzero original packets, and zero deletion preserves purity and rank. The corrected zero-or-range wording matches empty-index conventions.}
\field{Source locator}{\nolinkurl{editable/traditional/03_packets_and_graphs.tex}; statement lines 1011--1031.}
\field{Statement SHA-256}{\nolinkurl{2a2a8b2073162ddfce0cf7462765d5453176d369133cc86dc699c8d038e8fee8}}
\Needspace{12\baselineskip}\subsubsection*{84. Basic packet-pure rank bounds}
\phantomsection\label{rec:PSC-I-P-PACKET-PURE-RANK-BOUNDS}
\field{Stable ID}{\nolinkurl{PSC-I-P-PACKET-PURE-RANK-BOUNDS}}
\field{Kind}{proposition}
\textbf{Statement.}\par
For every $A$ module resolved into packets with $P_\alpha\in A$,
\[
\boxed{
\srank_A(X)
\leq
\psrank_{A,\Pcal}(X)
\leq
\sum_{\alpha\in I}
\srank_{P_\alpha A P_\alpha}(X_\alpha).
}
\]
The corner algebra on the right is unital on $X_\alpha$ with identity $P_\alpha$.

\begin{proof}
Every packet pure generating family is an unrestricted generating family, which gives the first inequality.  For each $\alpha$, choose a minimal generating family $S_\alpha\subseteq X_\alpha$ for the corner action $P_\alpha A P_\alpha$.  The union $S=\bigcup_\alpha S_\alpha$ is packet pure.  Since every $x\in X$ has the unique decomposition $x=\sum_\alpha x_\alpha$ and every $x_\alpha$ lies in $(P_\alpha A P_\alpha)S_\alpha\subseteq AS$, the union generates $X$.  Its cardinality is at most the displayed sum.
\end{proof}
\field{Formal correspondence}{covered. The left inequality is direct inclusion of admissible generating sets. The upper inequality chooses attained minimal generating sets under each genuine corner algebra, embeds them into X, and proves their union generates by projector completeness and corner\_action\_ambient. The union cardinal is bounded by the cardinal sum. This works for arbitrary infinite dimensions and arbitrary cardinal ranks, including zero packets.}
\field{Source locator}{\nolinkurl{editable/traditional/03_packets_and_graphs.tex}; statement lines 1034--1048; proof lines 1052--1054.}
\field{Statement SHA-256}{\nolinkurl{11a3ad129e84e51edb000d4027f785874b59f3867639a7813ab9c8d944bf81c0}}
\Needspace{12\baselineskip}\subsubsection*{85. Edge-surjective packet system}
\phantomsection\label{rec:PSC-I-D-EDGE-SURJECTIVE}
\field{Stable ID}{\nolinkurl{PSC-I-D-EDGE-SURJECTIVE}}
\field{Kind}{definition}
\textbf{Statement.}\par
The module resolved into packets is surjective along edges relative to $\Gamma(\mathcal G,\Pcal)$ when, for every graph edge $\alpha\to\beta$, there exists $B\in\A$ such that the pure block
\[
P_\beta BP_\alpha:X_\alpha\twoheadrightarrow X_\beta
\]
is surjective.
\field{Formal correspondence}{definition represented. The edge is TransferEdge P ops, not a supplied surrogate graph. The witness T belongs to Algebra.adjoin of the projectors and declared operators; restrictBlocks is the literal map from the source range to the target range.}
\field{Source locator}{\nolinkurl{editable/traditional/03_packets_and_graphs.tex}; statement lines 1056--1063.}
\field{Statement SHA-256}{\nolinkurl{46dd6f245aadc93dd24ac5505cd37893d5474fb31208bab5babfacf0fa6d1ce4}}
\Needspace{12\baselineskip}\subsubsection*{86. Source-component formula for packet-pure seed rank}
\phantomsection\label{rec:PSC-I-T-SOURCE-COMPONENT-SEED-RANK}
\field{Stable ID}{\nolinkurl{PSC-I-T-SOURCE-COMPONENT-SEED-RANK}}
\field{Kind}{theorem}
\textbf{Statement.}\par
Assume $I$ is finite, every packet space is nonzero, and the module resolved into packets is surjective along edges.  Let $\mathfrak C_0$ be the set of source strongly connected components of the condensation graph of $\Gamma(\mathcal G,\Pcal)$.  Then
\[
\boxed{
\psrank_{\A,\Pcal}(X)
=
\sum_{C\in\mathfrak C_0}r_C.
}
\]

\begin{proof}
Let $S$ be a packet pure generating family for $X$.  Discard zero seeds, which do not change the generated hull, and fix a source component $C$.  Seeds outside $C$ cannot reach $C$ by \hyperref[rec:PSC-I-P-REACHABILITY-OBSTRUCTION]{\nolinkurl{PSC-I-P-REACHABILITY-OBSTRUCTION}}.  Moreover, for $s\in X_C$ and $A\in\A$,
\[
P_CAs=(P_CAP_C)s\in\A_Cs.
\]
Thus the seeds of $S$ lying in $X_C$ must generate $X_C$ under $\A_C$, and there are at least $r_C$ of them.  Summing over disjoint source components gives
\[
\psrank_{\A,\Pcal}(X)
\geq
\sum_{C\in\mathfrak C_0}r_C.
\]

For the reverse inequality, choose for each source component $C$ a packet pure family $S_C$ of size $r_C$ generating $X_C$ under $\A_C$, and put $S=\bigcup_CS_C$.  Then every source component lies in $\A S$.  Traverse the finite condensation graph in a topological order.  Suppose all components preceding a component $D$ that is not a source have already been shown to lie in $\A S$.  Choose an incoming condensation edge from such a component to $D$, represented by a graph edge $\alpha\to\beta$ with $\alpha$ already generated and $\beta\in D$.  Surjectivity along edges gives a pure block
\[
B_{\beta\alpha}:X_\alpha\twoheadrightarrow X_\beta.
\]
Because $X_\alpha\subseteq\A S$ and $\A S$ is invariant under $\A$, surjectivity implies $X_\beta\subseteq\A S$.  For any other vertex $\gamma\in D$, choose a directed path inside the strongly connected component from $\beta$ to $\gamma$.  Successive surjective edge blocks propagate the full packet space along the path, so $X_\gamma\subseteq\A S$.  Hence $X_D\subseteq\A S$.  Induction over the condensation graph gives $X=\A S$, proving the upper bound.
\end{proof}
\field{Formal correspondence}{covered. SCC is the quotient by mutual reflexive-transitive reachability. Source is proved equivalent to absence of incoming condensation edges. Finite minimality proves every vertex is reached from a source. Generated-algebra support proves source projection compression. The lower bound removes zero seeds, proves those actually lying in each source generate its corner module, and injects the dependent union of these nonzero seed sets into the original generating set. The upper bound chooses attained local minima; actual surjective blocks propagate range containment along actual paths. No source coverage, compression identity, generating conclusion, or graph certificate is assumed in the final theorem. The theorem in fact does not require nonzero original packets, so it includes the printed hypotheses.}
\field{Source locator}{\nolinkurl{editable/traditional/03_packets_and_graphs.tex}; statement lines 1066--1076; proof lines 1081--1098.}
\field{Statement SHA-256}{\nolinkurl{0be97dd456719faad7e8e6ba4091ee877f58e7c7c900c8ae79a7b8385c19daea}}
\Needspace{12\baselineskip}\subsubsection*{87. Packet-propagating system}
\phantomsection\label{rec:PSC-I-D-PACKET-PROPAGATING}
\field{Stable ID}{\nolinkurl{PSC-I-D-PACKET-PROPAGATING}}
\field{Kind}{definition}
\textbf{Statement.}\par
A packet system that is surjective along edges is called packet propagating when every packet space $X_\alpha$ is cyclic under its diagonal corner algebra
\[
P_\alpha\A P_\alpha.
\]
\field{Formal correspondence}{definition represented. Cyclicity is the existing module Cyclic predicate for the genuine packetCornerAlgebra acting on its actual range. There are no extra local-rank-one certificates supplied.}
\field{Source locator}{\nolinkurl{editable/traditional/03_packets_and_graphs.tex}; statement lines 1105--1111.}
\field{Statement SHA-256}{\nolinkurl{4e9bfc875f14eeec60ebfab1d4906a00230861dfe0e37ea46c90639653fbc111}}
\Needspace{12\baselineskip}\subsubsection*{88. Unweighted source-component count}
\phantomsection\label{rec:PSC-I-C-SOURCE-COMPONENT-COUNT}
\field{Stable ID}{\nolinkurl{PSC-I-C-SOURCE-COMPONENT-COUNT}}
\field{Kind}{corollary}
\textbf{Statement.}\par
Assume the module resolved into packets is packet propagating and $X_\alpha\neq0$ for every $\alpha\in I$.  If $c_0$ denotes the number of source strongly connected components, then
\[
\boxed{
\psrank_{\A,\Pcal}(X)=c_0.
}
\]

\begin{proof}
Fix a source component $C$ and a vertex $\alpha\in C$.  Choose a cyclic vector $s\in X_\alpha$ for $P_\alpha\A P_\alpha$.  It generates all of $X_\alpha$.  Strong connectivity and surjectivity along edges then propagate the whole packet space to every vertex of $C$.  Thus $r_C=1$ for every source component.  Apply \hyperref[rec:PSC-I-T-SOURCE-COMPONENT-SEED-RANK]{\nolinkurl{PSC-I-T-SOURCE-COMPONENT-SEED-RANK}}.
\end{proof}
\field{Formal correspondence}{covered. The upper bound chooses a cyclic vector at one actual representative vertex of each source SCC and propagates full packet ranges to every vertex. The lower bound proves every generating set has a nonzero source projection, then uses purity and orthogonality to inject Sources into the chosen pure generating set. Nonzero packet range hypotheses are explicit. The finite cardinal-number form is proved; no finite generation or dimension assumption is introduced.}
\field{Source locator}{\nolinkurl{editable/traditional/03_packets_and_graphs.tex}; statement lines 1114--1122; proof lines 1126--1128.}
\field{Statement SHA-256}{\nolinkurl{f7d07d0e6a080b4b92c8a401d9663246fb377ecc6d930b8ecdb94ec07ea7e5dc}}
\Needspace{12\baselineskip}\subsubsection*{89. Mixed-support seeds can merge source requirements}
\phantomsection\label{rec:PSC-I-B-PACKET-PURE-BOUNDARY}
\field{Stable ID}{\nolinkurl{PSC-I-B-PACKET-PURE-BOUNDARY}}
\field{Kind}{boundary}
\textbf{Statement.}\par
The source component formulas concern packet pure seeds.  They do not give the unrestricted seed rank.  For example, let
\[
\A=\bbk\times\bbk,
\qquad
X=\bbk\oplus\bbk
\]
with the coordinate action and the two coordinate packet projectors.  The transfer graph has two source components and
\[
\psrank_{\A,\Pcal}(X)=2,
\]
but the mixed vector $(1,1)$ is cyclic, so
\[
\srank_{\A}(X)=1.
\]
\field{Formal correspondence}{scope boundary. The unrestricted mixed-seed example is consistent with the proved product-algebra formula. The packet-pure rank calculation is not separately formalized.}
\field{Source locator}{\nolinkurl{editable/traditional/03_packets_and_graphs.tex}; statement lines 1130--1146.}
\field{Statement SHA-256}{\nolinkurl{a3610fb8adcedb4fb91006a77f357f3f18c57678d044c21b8522730a456b0006}}
\clearpage\subsection{Transport and realization}
\Needspace{12\baselineskip}\subsubsection*{90. Coordinate transport}
\phantomsection\label{rec:PSC-I-T-COORDINATE-TRANSPORT}
\field{Stable ID}{\nolinkurl{PSC-I-T-COORDINATE-TRANSPORT}}
\field{Kind}{theorem}
\textbf{Statement.}\par
Let $U:X\to Y$ be a vector space isomorphism and let $\Pcal=(P_\alpha)$ be a packet system on $X$.  Define
\[
Q_\alpha:=UP_\alpha U^{-1}.
\]
Then $\Qcal=(Q_\alpha)$ is a packet system on $Y$.  For $\widetilde T=UTU^{-1}$,
\[
\boxed{
Q_\beta\widetilde TQ_\alpha
=
U(P_\beta TP_\alpha)U^{-1}.
}
\]
Hence packet dimensions, zero patterns, and support preorders are invariant under coordinate transport.  If a declared algebra generating family $\mathcal G$ is transported to $U\mathcal G U^{-1}$, then its transfer graph and reachability relation are also preserved.

\begin{proof}
Conjugation preserves products and sums, so
\[
Q_\alpha Q_\beta
=UP_\alpha P_\beta U^{-1}
=\delta_{\alpha\beta}Q_\alpha,
\qquad
\sum_\alpha Q_\alpha=I_Y.
\]
The block identity is direct multiplication:
\[
Q_\beta\widetilde TQ_\alpha
=UP_\beta TP_\alpha U^{-1}.
\]
All listed invariants are determined by these transported blocks.
\end{proof}
\field{Formal correspondence}{covered. conjugate constructs the actual family U P\_i U\ensuremath{^{-1}} and proves all projector identities. conjugate\_transferBlock proves the exact displayed equality. packetRangeEquiv restricts U to each actual packet range and yields cardinal dimension equality with explicit universe lifts. support\_pattern\_invariant proves preservation for every relation. TransferEdge and ReachableFrom are unchanged when the actual generator set is transported by conjugation.}
\field{Source locator}{\nolinkurl{editable/traditional/04_transport.tex}; statement lines 7--22; proof lines 26--41.}
\field{Statement SHA-256}{\nolinkurl{ae5e550ba65822235a1f8ab84c76a5a3ac876b28334523c9d845ece2af2d2d5a}}
\Needspace{12\baselineskip}\subsubsection*{91. Packet-compatible operator morphism}
\phantomsection\label{rec:PSC-I-D-PACKET-MORPHISM}
\field{Stable ID}{\nolinkurl{PSC-I-D-PACKET-MORPHISM}}
\field{Kind}{definition}
\textbf{Statement.}\par
Let $(X,\A,\Pcal)$ and $(Y,\B,\Qcal)$ be packet resolved operator modules with packet index sets $I$ and $J$.  A packet compatible operator morphism is a triple
\[
(\Phi,\varphi,\sigma)
\]
consisting of a unital algebra homomorphism $\varphi:\A\to\B$, a linear map $\Phi:X\to Y$ satisfying
\[
\Phi(Ax)=\varphi(A)\Phi(x),
\]
and a map $\sigma:I\to J$ satisfying
\[
Q_j\Phi
=
\Phi P_{\sigma^{-1}(j)},
\qquad
P_{\sigma^{-1}(j)}:=
\sum_{\sigma(\alpha)=j}P_\alpha.
\]
An empty fiber is interpreted as the zero projector.
\field{Formal correspondence}{definition represented. The source and target have actual concrete unital operator subalgebras and nonzero complete orthogonal packet projectors. The linear map, unital algebra map and entirely arbitrary index map are separately retained. Fiber projectors are literal finite sums; no injectivity or surjectivity condition is imposed on the index map.}
\field{Source locator}{\nolinkurl{editable/traditional/04_transport.tex}; statement lines 48--68.}
\field{Statement SHA-256}{\nolinkurl{e0e522a401795ef9ecf2497c058baf522801bf33dd581a837deaabd9cbe887bf}}
\Needspace{12\baselineskip}\subsubsection*{92. Morphism of seeded systems}
\phantomsection\label{rec:PSC-I-D-SEED-MORPHISM}
\field{Stable ID}{\nolinkurl{PSC-I-D-SEED-MORPHISM}}
\field{Kind}{definition}
\textbf{Statement.}\par
A morphism
\[
(X,\A,S,\Pcal)\longrightarrow(Y,\B,S',\Qcal)
\]
is a packet compatible operator morphism satisfying the seed condition
\[
\Phi(S)\subseteq\Hull_{\B}(S').
\]
In the topological category all maps and actions are required to be continuous and the seed condition uses the topological hull when declared.
\field{Formal correspondence}{definition represented. The algebraic condition uses the actual target algebra span, with arbitrary seed sets. The continuous variant separately retains Hausdorff topological vector-space structure, actual continuity of every algebra operator and packet projector, and actual continuity of the carrier morphism. Its seed condition is membership in the literal topologicalClosure of AlgebraGeneratedHull. The paper specifies an algebra of continuous operators, not a separate topology on that algebra, and the code states this convention explicitly.}
\field{Source locator}{\nolinkurl{editable/traditional/04_transport.tex}; statement lines 73--84.}
\field{Statement SHA-256}{\nolinkurl{88c7e9a865f68bac1df1727bd7cfc8c66e5b11d2e8be77ada048a83aaafd5637}}
\Needspace{12\baselineskip}\subsubsection*{93. Naturality of generated hulls and transfer blocks}
\phantomsection\label{rec:PSC-I-T-NATURALITY}
\field{Stable ID}{\nolinkurl{PSC-I-T-NATURALITY}}
\field{Kind}{theorem}
\textbf{Statement.}\par
Let $(\Phi,\varphi,\sigma)$ be a packet compatible operator morphism.  For every subset $S\subseteq X$,
\[
\boxed{
\Phi(\Hull_{\A}(S))
=
\Hull_{\varphi(\A)}(\Phi(S))
\subseteq
\Hull_{\B}(\Phi(S)).
}
\]
For $T\in\A$ and $i,j\in J$,
\[
\boxed{
Q_j\varphi(T)Q_i\Phi
=
\sum_{\substack{\sigma(\beta)=j\\\sigma(\alpha)=i}}
\Phi P_\beta TP_\alpha.
}
\]
If $\Phi$ is surjective and $S$ generates $X$, then $\Phi(S)$ generates $Y$ under $\varphi(\A)$ and hence under $\B$.  If, in addition, the morphism is seeded, then $S'$ generates $Y$.

\begin{proof}
Every element of $\Phi(\Hull_{\A}(S))$ is a finite sum of vectors
\[
\Phi(As)=\varphi(A)\Phi(s),
\]
so it belongs to $\Hull_{\varphi(\A)}(\Phi(S))$.  Conversely, every generator $\varphi(A)\Phi(s)$ equals $\Phi(As)$, giving equality.

For the block formula, use intertwining and packet compatibility twice:
\begin{align*}
Q_j\varphi(T)Q_i\Phi
&=Q_j\varphi(T)\Phi P_{\sigma^{-1}(i)}\\
&=Q_j\Phi TP_{\sigma^{-1}(i)}\\
&=\Phi P_{\sigma^{-1}(j)}TP_{\sigma^{-1}(i)}\\
&=\sum_{\substack{\sigma(\beta)=j\\\sigma(\alpha)=i}}
\Phi P_\beta TP_\alpha.
\end{align*}
If $\Phi$ is surjective and $\Hull_{\A}(S)=X$, the hull equality gives $Y=\Hull_{\varphi(\A)}(\Phi(S))$.  If $\Phi(S)\subseteq\Hull_{\B}(S')$, then $\Hull_{\B}(\Phi(S))\subseteq\Hull_{\B}(S')$, so $S'$ generates $Y$.
\end{proof}
\field{Formal correspondence}{covered. Hull naturality is actual Submodule.map equality with the span of image-algebra orbits; target-algebra inclusion is derived. The block formula is the exact double sum over possibly many-to-one fibers, at the level of actual linear maps. Surjective generation is derived from the map of top, then the seeded condition implies generation by the target seeds.}
\field{Source locator}{\nolinkurl{editable/traditional/04_transport.tex}; statement lines 87--109; proof lines 114--131.}
\field{Statement SHA-256}{\nolinkurl{e7fb174db6031054bdf06996ebad915dfaf307ad38b27fceb17442c46128863b}}
\Needspace{12\baselineskip}\subsubsection*{94. Category of seeded packet systems}
\phantomsection\label{rec:PSC-I-P-SEED-SYSTEM-CATEGORY}
\field{Stable ID}{\nolinkurl{PSC-I-P-SEED-SYSTEM-CATEGORY}}
\field{Kind}{proposition}
\textbf{Statement.}\par
Packet compatible operator morphisms compose, seeded morphisms compose, and identity morphisms satisfy both definitions.  Thus packet resolved operator modules form a category $\mathbf{PktMod}$ and seeded packet systems form a category $\mathbf{SeedSys}$.

\begin{proof}
Let $(\Phi,\varphi,\sigma):(X,\A,\Pcal)\to(Y,\B,\Qcal)$ and $(\Psi,\psi,\tau):(Y,\B,\Qcal)\to(Z,\Ccal,\Rcal)$ be morphisms.  Intertwining is preserved because
\[
\Psi\Phi(Ax)
=
\Psi(\varphi(A)\Phi x)
=
(\psi\circ\varphi)(A)\Psi\Phi x.
\]
For $k$ in the target packet set,
\begin{align*}
R_k\Psi\Phi
&=\Psi Q_{\tau^{-1}(k)}\Phi\\
&=\Psi\Phi
\sum_{\tau(\sigma(\alpha))=k}P_\alpha.
\end{align*}
Thus the composite packet map is $\tau\circ\sigma$.  If the morphisms are seeded, then
\[
\Phi(S)\subseteq\Hull_{\B}(S'),
\qquad
\Psi(S')\subseteq\Hull_{\Ccal}(S'').
\]
By \hyperref[rec:PSC-I-T-NATURALITY]{\nolinkurl{PSC-I-T-NATURALITY}},
\[
\Psi\Phi(S)
\subseteq
\Psi(\Hull_{\B}(S'))
\subseteq
\Hull_{\Ccal}(\Psi(S'))
\subseteq
\Hull_{\Ccal}(S'').
\]
Identity maps satisfy the definitions with identity algebra and packet maps.
\end{proof}
\field{Formal correspondence}{covered. Composition explicitly aggregates composite fibers by a finite-sum rearrangement and composes operator intertwiners. Algebraic seed closure is proved using span transport. The continuous variant proves invariance of the actual closed algebraic hull from individual operator continuity, maps the source algebraic hull into the target closed hull, then derives map\_closedHull from carrier continuity and closure minimality. Seeded composition applies this proved closure transport; identity seed membership and all category laws are proved. The actual algebraic and topological Category instances therefore cover the full stated category claim.}
\field{Source locator}{\nolinkurl{editable/traditional/04_transport.tex}; statement lines 138--141; proof lines 147--180.}
\field{Statement SHA-256}{\nolinkurl{3cc5c23c5df9094cc9c18b7b9539bf0cadbbe2e74dd249651920fd5034f347a1}}
\Needspace{12\baselineskip}\subsubsection*{95. Packet exact sequence for an invariant subspace}
\phantomsection\label{rec:PSC-I-T-PACKET-EXACT-SEQUENCE}
\field{Stable ID}{\nolinkurl{PSC-I-T-PACKET-EXACT-SEQUENCE}}
\field{Kind}{theorem}
\textbf{Statement.}\par
Let $K\leq X$ satisfy $P_\alpha K\subseteq K$ for every $\alpha$.  Define
\[
K_\alpha:=P_\alpha K=K\cap X_\alpha
\]
and
\[
\overline P_\alpha(x+K):=P_\alpha x+K.
\]
Then:
\begin{enumerate}[label=(\roman*)]
\item The nonzero restrictions $P_\alpha|_K$ form a packet system on $K$.
\item The nonzero $\overline P_\alpha$ form a packet system on $X/K$.
\item For every $\alpha$ there is a short exact sequence
\[
0\longrightarrow K_\alpha
\longrightarrow X_\alpha
\longrightarrow \overline P_\alpha(X/K)
\longrightarrow0,
\]
and
\[
\overline P_\alpha(X/K)
\cong
X_\alpha/(X_\alpha\cap K)
\cong
(X_\alpha+K)/K.
\]
\item The full exact sequence decomposes packetwise as the direct sum of these short exact sequences.
\end{enumerate}

\begin{proof}
Restriction preserves idempotency, pairwise annihilation, and the identity sum on $K$.  If $x\in K\cap X_\alpha$, then $x=P_\alpha x\in P_\alpha K$, proving $K_\alpha=K\cap X_\alpha$.

The quotient operator is well defined because replacing $x$ by $x+k$ changes $P_\alpha x$ by $P_\alpha k\in K$.  The idempotent and annihilation identities descend, and the sum is the identity on $X/K$.  The map $X_\alpha\to X/K$ has image $(X_\alpha+K)/K$ and kernel $X_\alpha\cap K$, giving the short exact sequence and the quotient isomorphisms.  Exact packet coordinates identify the original sequence with the direct sum over $\alpha$.
\end{proof}
\field{Formal correspondence}{covered. K is an actual invariant submodule. Restriction/quotient projectors are constructed; restricted range equals K intersect X\_i. Actual packet inclusion is injective, its range equals the actual quotient-map kernel, and quotient map is onto. First/second quotient isomorphisms use denominator comaps in the correct ambient type. Pointwise coordinate maps intertwine with original inclusion/quotient and form a short exact sequence; exact direct-sum coordinates give the full decomposition.}
\field{Source locator}{\nolinkurl{editable/traditional/04_transport.tex}; statement lines 185--216; proof lines 221--225.}
\field{Statement SHA-256}{\nolinkurl{7311937b07e98adf3e6f91ad665358380cab6ce26e3daeffdbf4787d7a01de56}}
\Needspace{12\baselineskip}\subsubsection*{96. Packet descent through an invariant quotient}
\phantomsection\label{rec:PSC-I-T-QUOTIENT-PACKET-DESCENT}
\field{Stable ID}{\nolinkurl{PSC-I-T-QUOTIENT-PACKET-DESCENT}}
\field{Kind}{theorem}
\textbf{Statement.}\par
Under the hypotheses of \hyperref[rec:PSC-I-T-PACKET-EXACT-SEQUENCE]{\nolinkurl{PSC-I-T-PACKET-EXACT-SEQUENCE}}, the operators
\[
\overline P_\alpha(x+K)=P_\alpha x+K
\]
are well defined pairwise annihilating idempotents with sum $I_{X/K}$, and the quotient map satisfies
\[
\boxed{\pi P_\alpha=\overline P_\alpha\pi.}
\]
After zero packets are removed, they form the canonical quotient packet system.

\begin{proof}
All assertions are the quotient part of \hyperref[rec:PSC-I-T-PACKET-EXACT-SEQUENCE]{\nolinkurl{PSC-I-T-PACKET-EXACT-SEQUENCE}}.  The intertwining identity follows directly from the definition:
\[
\overline P_\alpha\pi(x)=P_\alpha x+K=\pi P_\alpha(x).
\]
\end{proof}
\field{Formal correspondence}{covered. Actual mapQ uses invariance of K. Quotient representative calculations prove the laws and commuting identity. Zero removal preserves completeness. Quotient surjectivity proves uniqueness of every descended operator, so canonicity is established.}
\field{Source locator}{\nolinkurl{editable/traditional/04_transport.tex}; statement lines 232--243; proof lines 247--252.}
\field{Statement SHA-256}{\nolinkurl{b51391a00350d220fc92610f351826489ca69d166e27e266a1ecc6ecc7a2ffc7}}
\Needspace{12\baselineskip}\subsubsection*{97. Restriction and quotient of transfer blocks}
\phantomsection\label{rec:PSC-I-C-QUOTIENT-PACKET-DESCENT}
\field{Stable ID}{\nolinkurl{PSC-I-C-QUOTIENT-PACKET-DESCENT}}
\field{Kind}{corollary}
\textbf{Statement.}\par
Under the hypotheses of \hyperref[rec:PSC-I-T-PACKET-EXACT-SEQUENCE]{\nolinkurl{PSC-I-T-PACKET-EXACT-SEQUENCE}}, let $T\in\End(X)$ preserve $K$.  Then every block $T_{\beta\alpha}$ maps $K_\alpha$ into $K_\beta$, and the induced quotient operator $\overline T$ satisfies
\[
\boxed{
\overline T_{\beta\alpha}\,\pi
=
\pi T_{\beta\alpha},
}
\]
where $\pi:X\twoheadrightarrow X/K$.

\begin{proof}
If $x\in K_\alpha$, then $Tx\in K$, so
\[
T_{\beta\alpha}x=P_\beta Tx\in K\cap X_\beta=K_\beta.
\]
On the quotient,
\[
\overline T_{\beta\alpha}\pi
=
\overline P_\beta\overline T\overline P_\alpha\pi
=
\pi P_\beta TP_\alpha.
\]
\end{proof}
\field{Formal correspondence}{covered. Invariance under P\_i,T,P\_j proves preservation of K and membership in the target packet. Tbar is actual K.mapQ K T hT, and representative calculation proves the quotient-block identity. Lean orientation is the symmetric equality of the printed formula.}
\field{Source locator}{\nolinkurl{editable/traditional/04_transport.tex}; statement lines 254--265; proof lines 269--282.}
\field{Statement SHA-256}{\nolinkurl{4fa0ae8effe1e5472dcde2be67829cba4ae586df9853d252baff055f3ee7fd46}}
\Needspace{12\baselineskip}\subsubsection*{98. Topological packet exact sequence}
\phantomsection\label{rec:PSC-I-C-TOPOLOGICAL-PACKET-EXACT-SEQUENCE}
\field{Stable ID}{\nolinkurl{PSC-I-C-TOPOLOGICAL-PACKET-EXACT-SEQUENCE}}
\field{Kind}{corollary}
\textbf{Statement.}\par
If $X$ is Banach, the $P_\alpha$ are bounded, and $K$ is closed and packet invariant, then every packet sequence in \hyperref[rec:PSC-I-T-PACKET-EXACT-SEQUENCE]{\nolinkurl{PSC-I-T-PACKET-EXACT-SEQUENCE}} is a short exact sequence of Banach spaces with bounded maps, and the quotient projectors are bounded.

\begin{proof}
Closedness of $K$ makes $X/K$ Banach.  The restricted projectors are bounded, and the quotient projectors are bounded because
\[
\|\overline P_\alpha(x+K)\|_{X/K}
\leq
\|P_\alpha\|\,\|x+K\|_{X/K}.
\]
The packet spaces are closed ranges of bounded idempotents.
\end{proof}
\field{Formal correspondence}{covered. The generic theorem works over any NormedField, including a trivial valuation, with explicit boundedness of the original projectors. Restriction and quotient CLMs are the actual algebraic packet maps. Near-minimal quotient representatives transfer every original bound to the quotient; the norm bound uses the usual infimum-of-bounds norm of the bounded additive projector. Actual packet inclusion and quotient maps are bounded with constant one and retain injectivity, exactness, and surjectivity. Closed K and closed bounded-idempotent ranges make all three terms complete.}
\field{Source locator}{\nolinkurl{editable/traditional/04_transport.tex}; statement lines 284--287; proof lines 292--300.}
\field{Statement SHA-256}{\nolinkurl{2590b3fce057e3d054522a58a728f46dbf3c609c95ff6ab061ff654ffd72fee1}}
\Needspace{12\baselineskip}\subsubsection*{99. Dense realization theorem}
\phantomsection\label{rec:PSC-I-T-DENSE-REALIZATION}
\field{Stable ID}{\nolinkurl{PSC-I-T-DENSE-REALIZATION}}
\field{Kind}{theorem}
\textbf{Statement.}\par
Let $X$ and $Y$ be Hausdorff topological vector spaces, let $\Phi:X\to Y$ be continuous, surjective, and intertwining with respect to $\varphi$, and let $S$ topologically generate $X$.  Then $\Phi(S)$ topologically generates $Y$.  If finite continuous packet systems satisfy
\[
Q_j\Phi=\Phi P_{\sigma^{-1}(j)},
\]
then
\[
\boxed{
Q_jY=\Phi(P_{\sigma^{-1}(j)}X)
}
\]
for every $j$.

\begin{proof}
Given $y=\Phi(x)$, choose a net $x_\lambda\in\Hull_{\A}(S)$ converging to $x$.  Then
\[
\Phi(x_\lambda)\to y,
\]
and naturality places every $\Phi(x_\lambda)$ in $\Hull_{\B}(\Phi(S))$.  Hence the latter hull is dense.  For packets, one inclusion follows from compatibility.  Conversely, if $y=Q_jy$, choose $x$ with $\Phi(x)=y$.  Then
\[
y=Q_j\Phi(x)=\Phi P_{\sigma^{-1}(j)}x.
\]
\end{proof}
\field{Formal correspondence}{covered. For the actual continuous linear intertwiner in the packet-morphism convention, the image of the source algebraic operator hull is contained in the target hull, and continuous surjectivity transports density. The actual finite grouped projector is the sum over the \ensuremath{\sigma}-fibre, including empty fibres. Surjectivity and the given intertwining identity prove exact equality of image and target range. No norm or metrizability assumptions occur.}
\field{Source locator}{\nolinkurl{editable/traditional/04_transport.tex}; statement lines 305--318; proof lines 322--331.}
\field{Statement SHA-256}{\nolinkurl{58bf2b9bfb4ab3ae9f3309b03cc231e0b59c81b7fe2b6c99e4b74884aaf15bca}}
\Needspace{12\baselineskip}\subsubsection*{100. Quotient topology and completed realization}
\phantomsection\label{rec:PSC-I-T-QUOTIENT-TOPOLOGY-COMPLETION}
\field{Stable ID}{\nolinkurl{PSC-I-T-QUOTIENT-TOPOLOGY-COMPLETION}}
\field{Kind}{theorem}
\textbf{Statement.}\par
Let $\Phi:X\twoheadrightarrow Y$ be a continuous linear quotient map between Hausdorff topological vector spaces, and put $K=\Ker\Phi$.  Then the induced map
\[
\widetilde\Phi:X/K\longrightarrow Y
\]
is a topological isomorphism.  If $D\subseteq X$ is a linear subspace whose image $\Phi(D)$ is dense in a complete space $Y$, then the completion of
\[
D/(D\cap K)
\]
in the topology transported from $Y$ is canonically isomorphic to $Y$.

\begin{proof}
The algebraic first isomorphism theorem gives a bijection $\widetilde\Phi$.  By definition of the quotient topology, a set $U\subseteq X/K$ is open exactly when its preimage in $X$ is open.  Because $\Phi$ is a quotient map, this is equivalent to $\widetilde\Phi(U)$ being open in $Y$.  Thus $\widetilde\Phi$ is a homeomorphism.  The second statement applies the dense completion argument of \hyperref[rec:PSC-I-T-COMPLETED-PRESENTATION]{\nolinkurl{PSC-I-T-COMPLETED-PRESENTATION}} to the injective map $D/(D\cap K)\to Y$.
\end{proof}
\field{Formal correspondence}{covered. The actual first-isomorphism linear equivalence is continuous both ways using IsQuotientMap, not continuity plus surjectivity alone. The restriction to D uses its actual kernel, definitionally comap(D inclusion,ker \ensuremath{\Phi}); density is exactly \ensuremath{\Phi}(D). The same explicit transported-uniformity quotient completion gives the canonical target equivalence and agrees with \ensuremath{\Phi} on represented classes. No original-domain quotient topology is silently substituted for the printed transported topology.}
\field{Source locator}{\nolinkurl{editable/traditional/04_transport.tex}; statement lines 333--344; proof lines 348--350.}
\field{Statement SHA-256}{\nolinkurl{2e986d76e93d67446708502bfae8275f63d8ab0eaf8dcc3e3630ee7f6b3fa077}}
\Needspace{12\baselineskip}\subsubsection*{101. Continuous surjections and quotient topology}
\phantomsection\label{rec:PSC-I-B-TOPOLOGICAL-QUOTIENT}
\field{Stable ID}{\nolinkurl{PSC-I-B-TOPOLOGICAL-QUOTIENT}}
\field{Kind}{boundary}
\textbf{Statement.}\par
Continuity and surjectivity transport density as in \hyperref[rec:PSC-I-T-DENSE-REALIZATION]{\nolinkurl{PSC-I-T-DENSE-REALIZATION}}.  They do not by themselves identify the given topology on $Y$ with the quotient topology from $X$.  The stronger conclusion in \hyperref[rec:PSC-I-T-QUOTIENT-TOPOLOGY-COMPLETION]{\nolinkurl{PSC-I-T-QUOTIENT-TOPOLOGY-COMPLETION}} requires the quotient map hypothesis or an equivalent open mapping condition.
\field{Formal correspondence}{scope boundary. The quotient-map requirement is preserved; no open-map theorem is inferred from continuity alone.}
\field{Source locator}{\nolinkurl{editable/traditional/04_transport.tex}; statement lines 352--355.}
\field{Statement SHA-256}{\nolinkurl{855989a0fadfaf7174d83889d9a2c4208bc722f866b35c5da99de3e7ecf1b9f7}}
\Needspace{12\baselineskip}\subsubsection*{102. Realization functor and preservation profile}
\phantomsection\label{rec:PSC-I-D-REALIZATION-FUNCTOR}
\field{Stable ID}{\nolinkurl{PSC-I-D-REALIZATION-FUNCTOR}}
\field{Kind}{definition}
\textbf{Statement.}\par
Let $\mathbf{SeedSys}_{\mathrm{abs}}$ be an abstract category of seed systems and let $\mathbf C_D$ be a domain category.  A realization functor is a functor
\[
\mathsf{Ev}_D:\mathbf{SeedSys}_{\mathrm{abs}}\longrightarrow\mathbf C_D
\]
with a declared preservation profile specifying which of the following structures are preserved:
\[
\begin{gathered}
0,\ I,\ \text{finite sums},\ \text{composition},\ \text{biproducts},\\
\text{kernels},\ \text{cokernels},\ \text{quotients},\ \text{tensor products},\\
\text{packet idempotents},\ \text{transfer blocks},\ \text{topological completions}.
\end{gathered}
\]
Preservation of a structure means preservation of its defining identity or universal property, not merely a similarity of notation.
\field{Formal correspondence}{documentary. The profile is represented by actual additive/linear functor and PreservesLimit/PreservesColimit instances, with typed syntax for equational operations. Universal constructions also use explicit structured categories, an induced functor, an underlying comparison isomorphism, and actual initial/terminal preservation. No Boolean profile field is accepted in place of preservation evidence.}
\field{Source locator}{\nolinkurl{editable/traditional/04_transport.tex}; statement lines 382--397.}
\field{Statement SHA-256}{\nolinkurl{152d8ee6b5bb1735f7655c6eeee72210d213959c00c9e6168d554eee25b7c5fe}}
\Needspace{12\baselineskip}\subsubsection*{103. Equational realization principle}
\phantomsection\label{rec:PSC-I-T-REALIZATION-EXPORT}
\field{Stable ID}{\nolinkurl{PSC-I-T-REALIZATION-EXPORT}}
\field{Kind}{theorem}
\textbf{Statement.}\par
Let $F:\mathbf C\to\mathbf D$ be an additive functor between preadditive categories.  Any equality of morphisms constructed from declared generating morphisms by finite sums, additive inverses, identities, zero maps, and composition is preserved by $F$.  Scalar multiplication by a coefficient ring is also preserved when the categories are enriched over that ring and $F$ is linear for the declared enrichment.  If a proof additionally uses a chosen biproduct, kernel, cokernel, quotient, tensor product, or completion and $F$ preserves that construction in its declared profile, then the corresponding proof diagram is transported to $\mathbf D$.

\begin{proof}
Functoriality gives
\[
F(I)=I,
\qquad
F(g\circ f)=F(g)\circ F(f),
\]
and additivity gives
\[
F(0)=0,
\qquad
F(f+g)=F(f)+F(g),
\qquad
F(-f)=-F(f).
\]
If the declared enrichment is present, linearity additionally gives $F(af)=aF(f)$ for every coefficient $a$.  A structural induction on the construction of a morphism expression therefore carries every permitted equational proof to its image under $F$.  If a step invokes a universal construction, preservation means that the image diagram satisfies the same universal property.  The uniquely induced morphism in the source proof is therefore carried to the uniquely induced morphism in the target proof.  Repeating this argument through the finite proof diagram transports the conclusion.
\end{proof}
\field{Formal correspondence}{covered. Structural induction on well-typed morphism expressions proves transport of identities, zero, addition, negation, scalar multiplication and composition. The additive-only statement uses the canonical integer enrichment, and no extra scalar-linearity assumption is inserted. Actual limit/colimit diagrams are transported under the matching preservation instance. The conditional universal-construction clause is explicit: a declared profile supplies the structured categories of the universal problem, an induced functor, compatibility with the underlying realization functor, and actual initial/terminal preservation. CategoryUniversal derives the mapped initial/terminal property, equality of the uniquely induced arrows, and the underlying commuting comparison squares. Tensor and completion preservation are legitimate conditional inputs, not falsely inferred from additivity.}
\field{Source locator}{\nolinkurl{editable/traditional/04_transport.tex}; statement lines 400--403; proof lines 408--424.}
\field{Statement SHA-256}{\nolinkurl{e0352bdb7ba75e24d4c7e6e0259802f47800c4945c13efe4b3c68edcacef678d}}
\Needspace{12\baselineskip}\subsubsection*{104. Realization export principle}
\phantomsection\label{rec:PSC-I-P-REALIZATION-EXPORT}
\field{Stable ID}{\nolinkurl{PSC-I-P-REALIZATION-EXPORT}}
\field{Kind}{proposition}
\textbf{Statement.}\par
Every canonical Paper I identity whose proof uses only the operations listed in a realization functor's verified preservation profile remains valid after applying that functor.  If the proof uses a domain specific hypothesis not encoded by the functor, that hypothesis must be verified separately before the transported statement is promoted to a theorem.

\begin{proof}
The first assertion is \hyperref[rec:PSC-I-T-REALIZATION-EXPORT]{\nolinkurl{PSC-I-T-REALIZATION-EXPORT}}.  The second is a logical boundary: a functor transports established identities and preserved universal constructions, but it does not create an unverified analytic, topological, positivity, finiteness, or integrality hypothesis.
\end{proof}
\field{Formal correspondence}{covered. This is the methodological specialization of record 103 to already established Paper I identities. The formal equations are transported from a proved source equality, not assumed as the desired target equality. The universal-property step uses the explicit structured-category profile of CategoryUniversal. No new analytic/topological/positivity/finiteness/integrality assumption is generated.}
\field{Source locator}{\nolinkurl{editable/traditional/04_transport.tex}; statement lines 432--435; proof lines 439--441.}
\field{Statement SHA-256}{\nolinkurl{bf39b86447c8332d60216e48d23129b6ae5b096611e3febd9d09a699f1a9a557}}
\Needspace{12\baselineskip}\subsubsection*{105. Semantic realization criterion}
\phantomsection\label{rec:PSC-I-D-SEMANTIC-REALIZATION}
\field{Stable ID}{\nolinkurl{PSC-I-D-SEMANTIC-REALIZATION}}
\field{Kind}{definition}
\textbf{Statement.}\par
A domain theorem is a semantic realization of a canonical seed calculus theorem only when:
\begin{enumerate}[label=(\roman*)]
\item An explicit realization functor or intertwining construction is given.
\item The preservation profile covers every construction used in the canonical proof.
\item Every additional analytic, finiteness, convergence, positivity, or integrality hypothesis is verified in the domain.
\item The resulting conclusion is stated in the domain's native objects.
\end{enumerate}
Reuse of the words ``seed'', ``packet'', or ``transfer'' without these data is only syntactic reuse.
\field{Formal correspondence}{documentary. A methodology and scope checklist, not automatic formal certification.}
\field{Source locator}{\nolinkurl{editable/traditional/04_transport.tex}; statement lines 443--453.}
\field{Statement SHA-256}{\nolinkurl{7bfb521bf39d1f519716ff37b9ffcf3ae76ce4f62fe121586daa29e66e7223df}}
\Needspace{12\baselineskip}\subsubsection*{106. Reserved example slots}
\phantomsection\label{rec:PSC-I-ALG-EXAMPLE-SLOTS}
\field{Stable ID}{\nolinkurl{PSC-I-ALG-EXAMPLE-SLOTS}}
\field{Kind}{definition}
\textbf{Statement.}\par
The foundational calculus is illustrated by four canonical algebraic realizations:
\begin{enumerate}[label=(\roman*)]
\item A typed quiver reachability model.
\item An algebraic coordinate ring presentation.
\item A companion and repeated primary single operator model.
\item Path and cycle adjacency models.
\end{enumerate}
The worked model sections below instantiate these four interfaces. A bilateral shift example records the analytic completion mechanism also used in Paper V. Computational specifications connect the proofs with the verification companion and are not premises of the foundational theorems.
\field{Formal correspondence}{documentary. Example slots are editorial. Existing examples and algorithm specifications remain outside the 69 proved-record denominator.}
\field{Source locator}{\nolinkurl{editable/traditional/04_transport.tex}; statement lines 470--480.}
\field{Statement SHA-256}{\nolinkurl{a870d22072364c0a743aede20db499ccda1ce4c62a009a82dbaefd76f525375d}}
\Needspace{12\baselineskip}\subsubsection*{107. Paper I export object}
\phantomsection\label{rec:PSC-I-EXPORT}
\field{Stable ID}{\nolinkurl{PSC-I-EXPORT}}
\field{Kind}{definition}
\textbf{Statement.}\par
The canonical Paper I export is
\[
\boxed{
\mathsf{PSC}_I=
\left(
\begin{array}{c}
\text{typed grammars and realized operator algebras},\\
\text{seed systems, free presentations, and seed rank},\\
\text{module-top, semisimple, and single-operator formulas},\\
\text{finite and unconditional packet coordinates},\\
\text{transfer blocks and support algebras},\\
\text{reachability and packet-pure source generation},\\
\text{morphisms, invariant exact sequences, quotients, completions},\\
\text{realization functors and equational theorem transport}
\end{array}
\right).
}
\]
\field{Formal correspondence}{documentary. An export catalogue, not one additional kernel theorem or a completed formal package object.}
\field{Source locator}{\nolinkurl{editable/traditional/04_transport.tex}; statement lines 483--503.}
\field{Statement SHA-256}{\nolinkurl{5c58d6e198263b11ae3e1f98701de3cf1b70b5bc5da229d1ed705b607efe2aa7}}
\Needspace{12\baselineskip}\subsubsection*{108. Exact limits of Paper I}
\phantomsection\label{rec:PSC-I-B-FINAL-CLAIM-BOUNDARY}
\field{Stable ID}{\nolinkurl{PSC-I-B-FINAL-CLAIM-BOUNDARY}}
\field{Kind}{boundary}
\textbf{Statement.}\par
The existence of an exact packet decomposition does not imply that the decomposition is computationally discoverable, numerically well conditioned, statistically identifiable, dynamically invariant, or mathematically useful.  The condition number in \hyperref[rec:PSC-I-D-PACKET-CONDITION-NUMBER]{\nolinkurl{PSC-I-D-PACKET-CONDITION-NUMBER}} quantifies only coordinate stability once the packets are supplied.  Refinement towers, filtered transfer, nonlinear obstruction objects, automatic packet recovery, continuum limits, arithmetic lifts, and automorphic occupancy require additional theorems in later papers.
\field{Formal correspondence}{scope boundary. No discoverability, conditioning, identifiability or later-paper existence claim follows automatically from finite packet identities.}
\field{Source locator}{\nolinkurl{editable/traditional/04_transport.tex}; statement lines 532--535.}
\field{Statement SHA-256}{\nolinkurl{d059c43a34200d583173f81a335b2acb713d6ddbf0d46054763a575b1fdef0c7}}

\clearpage\section{Typed core signatures}
These twenty signatures are reproduced from \nolinkurl{kir/core_typed_signatures.json}. Unicode mathematical characters are represented with standard JSON escapes in the listings. Their meanings are unchanged. The schema does not make these strings into independently checked proof terms.
\Needspace{7\baselineskip}\subsection*{Linear seed system}
\hyperref[rec:PSC-I-D-SEED-SYSTEM]{\nolinkurl{PSC-I-D-SEED-SYSTEM}}
\begin{lstlisting}
{
  "params": [
    "A:OperatorAlgebra",
    "M:Module(A)",
    "S:FiniteFamily(M)"
  ],
  "returns": "SeedSystem(A,M,S)"
}
\end{lstlisting}
\Needspace{7\baselineskip}\subsection*{Generated hull}
\hyperref[rec:PSC-I-D-GENERATED-HULL]{\nolinkurl{PSC-I-D-GENERATED-HULL}}
\begin{lstlisting}
{
  "params": [
    "A:OperatorAlgebra",
    "M:Module(A)",
    "S:Family(M)"
  ],
  "returns": "Submodule(M)"
}
\end{lstlisting}
\Needspace{7\baselineskip}\subsection*{Seed rank}
\hyperref[rec:PSC-I-D-SEED-RANK]{\nolinkurl{PSC-I-D-SEED-RANK}}
\begin{lstlisting}
{
  "params": [
    "A:OperatorAlgebra",
    "M:Module(A)"
  ],
  "returns": "Cardinal"
}
\end{lstlisting}
\Needspace{7\baselineskip}\subsection*{Finite packet system}
\hyperref[rec:PSC-I-D-PACKET-SYSTEM]{\nolinkurl{PSC-I-D-PACKET-SYSTEM}}
\begin{lstlisting}
{
  "params": [
    "K:Field",
    "X:VectorSpace(K)",
    "I:FiniteIndexSet",
    "P:I\u2192End(X)"
  ],
  "requires": [
    "idempotent(P)",
    "pairwise_annihilating(P)",
    "sum(P)=id_X"
  ],
  "returns": "PacketSystem(X,I)"
}
\end{lstlisting}
\Needspace{7\baselineskip}\subsection*{Packet analysis and synthesis}
\hyperref[rec:PSC-I-D-ANALYSIS-SYNTHESIS]{\nolinkurl{PSC-I-D-ANALYSIS-SYNTHESIS}}
\begin{lstlisting}
{
  "params": [
    "P:PacketSystem(X,I)"
  ],
  "returns": [
    "Analyze_P:X\u2192DirectSum(I,PacketSpace(P,i))",
    "Synthesize_P:DirectSum(I,PacketSpace(P,i))\u2192X"
  ]
}
\end{lstlisting}
\Needspace{7\baselineskip}\subsection*{Exact packet-coordinate isomorphism}
\hyperref[rec:PSC-I-T-EXACT-PACKET-COORDINATES]{\nolinkurl{PSC-I-T-EXACT-PACKET-COORDINATES}}
\begin{lstlisting}
{
  "params": [
    "P:PacketSystem(X,I)"
  ],
  "returns": "Equivalence(X,DirectSum(I,PacketSpace(P,i)))",
  "certificate": "exact_inverse"
}
\end{lstlisting}
\Needspace{7\baselineskip}\subsection*{Transfer block}
\hyperref[rec:PSC-I-D-TRANSFER-BLOCK]{\nolinkurl{PSC-I-D-TRANSFER-BLOCK}}
\begin{lstlisting}
{
  "params": [
    "P:PacketSystem(X,I)",
    "T:End(X)",
    "alpha:I",
    "beta:I"
  ],
  "returns": "LinearMap(PacketSpace(P,alpha),PacketSpace(P,beta))"
}
\end{lstlisting}
\Needspace{7\baselineskip}\subsection*{Exact operator reconstruction}
\hyperref[rec:PSC-I-T-OPERATOR-RECONSTRUCTION]{\nolinkurl{PSC-I-T-OPERATOR-RECONSTRUCTION}}
\begin{lstlisting}
{
  "params": [
    "P:PacketSystem(X,I)",
    "T:End(X)"
  ],
  "returns": "Equality(T,Sum(alpha,beta,TransferBlock(P,T,beta,alpha)))"
}
\end{lstlisting}
\Needspace{7\baselineskip}\subsection*{Packet matrix multiplication}
\hyperref[rec:PSC-I-T-BLOCK-MULTIPLICATION]{\nolinkurl{PSC-I-T-BLOCK-MULTIPLICATION}}
\begin{lstlisting}
{
  "params": [
    "P:PacketSystem(X,I)",
    "S:End(X)",
    "T:End(X)",
    "alpha:I",
    "gamma:I"
  ],
  "returns": "Equality(Block(S\u2218T,gamma,alpha),Sum(beta,Block(S,gamma,beta)\u2218Block(T,beta,alpha)))"
}
\end{lstlisting}
\Needspace{7\baselineskip}\subsection*{Packet powers and paths}
\hyperref[rec:PSC-I-C-BLOCK-POWERS]{\nolinkurl{PSC-I-C-BLOCK-POWERS}}
\begin{lstlisting}
{
  "params": [
    "P:PacketSystem(X,I)",
    "T:End(X)",
    "m:Nat"
  ],
  "returns": "PacketPathExpansion(T,m)"
}
\end{lstlisting}
\Needspace{7\baselineskip}\subsection*{Blockwise intertwining equations}
\hyperref[rec:PSC-I-T-BLOCK-INTERTWINING]{\nolinkurl{PSC-I-T-BLOCK-INTERTWINING}}
\begin{lstlisting}
{
  "params": [
    "P:PacketSystem(X,I)",
    "Q:PacketSystem(Y,J)",
    "T:End(X)",
    "S:End(Y)",
    "U:LinearMap(X,Y)"
  ],
  "returns": "Iff(U\u2218T=S\u2218U,LocalBlockIntertwining(P,Q,T,S,U))"
}
\end{lstlisting}
\Needspace{7\baselineskip}\subsection*{Full block-algebra representation}
\hyperref[rec:PSC-I-T-BLOCK-ALGEBRA]{\nolinkurl{PSC-I-T-BLOCK-ALGEBRA}}
\begin{lstlisting}
{
  "params": [
    "P:PacketSystem(X,I)"
  ],
  "returns": "AlgebraEquiv(End(X),DependentBlockAlgebra(P))"
}
\end{lstlisting}
\Needspace{7\baselineskip}\subsection*{Support relation and support space}
\hyperref[rec:PSC-I-D-SUPPORT-RELATION]{\nolinkurl{PSC-I-D-SUPPORT-RELATION}}
\begin{lstlisting}
{
  "params": [
    "I:FiniteIndexSet",
    "R:Relation(I)"
  ],
  "returns": "SupportSpace(R)"
}
\end{lstlisting}
\Needspace{7\baselineskip}\subsection*{Coordinate transport}
\hyperref[rec:PSC-I-T-COORDINATE-TRANSPORT]{\nolinkurl{PSC-I-T-COORDINATE-TRANSPORT}}
\begin{lstlisting}
{
  "params": [
    "U:LinearEquiv(X,Y)",
    "P:PacketSystem(X,I)",
    "T:End(X)"
  ],
  "returns": "TransportedPacketOperator(U,P,T)"
}
\end{lstlisting}
\Needspace{7\baselineskip}\subsection*{Packet exact sequence for an invariant subspace}
\hyperref[rec:PSC-I-T-PACKET-EXACT-SEQUENCE]{\nolinkurl{PSC-I-T-PACKET-EXACT-SEQUENCE}}
\begin{lstlisting}
{
  "params": [
    "P:PacketSystem(X,I)",
    "K:InvariantSubspace(P)"
  ],
  "returns": "PacketwiseShortExactSequence(K,X,X/K)"
}
\end{lstlisting}
\Needspace{7\baselineskip}\subsection*{Packet descent through an invariant quotient}
\hyperref[rec:PSC-I-T-QUOTIENT-PACKET-DESCENT]{\nolinkurl{PSC-I-T-QUOTIENT-PACKET-DESCENT}}
\begin{lstlisting}
{
  "params": [
    "P:PacketSystem(X,I)",
    "K:InvariantSubspace(P)"
  ],
  "returns": "PacketSystem(X/K,I)"
}
\end{lstlisting}
\Needspace{7\baselineskip}\subsection*{Naturality of generated hulls and transfer blocks}
\hyperref[rec:PSC-I-T-NATURALITY]{\nolinkurl{PSC-I-T-NATURALITY}}
\begin{lstlisting}
{
  "params": [
    "Phi:PacketCompatibleOperatorMorphism",
    "S:SeedSet"
  ],
  "returns": "NaturalityCertificate(Phi,S)"
}
\end{lstlisting}
\Needspace{7\baselineskip}\subsection*{Realization functor and preservation profile}
\hyperref[rec:PSC-I-D-REALIZATION-FUNCTOR]{\nolinkurl{PSC-I-D-REALIZATION-FUNCTOR}}
\begin{lstlisting}
{
  "params": [
    "F:AdditiveFunctor(C,D)",
    "profile:PreservationProfile"
  ],
  "returns": "RealizationFunctor(C,D,profile)"
}
\end{lstlisting}
\Needspace{7\baselineskip}\subsection*{Equational realization principle}
\hyperref[rec:PSC-I-T-REALIZATION-EXPORT]{\nolinkurl{PSC-I-T-REALIZATION-EXPORT}}
\begin{lstlisting}
{
  "params": [
    "Theta:TheoremRecord",
    "F:RealizationFunctor"
  ],
  "requires": [
    "operations_used(Theta)\u2286verified_preservation(F)",
    "target_hypotheses_verified"
  ],
  "returns": "TransportCertificate(Theta,F)"
}
\end{lstlisting}
\Needspace{7\baselineskip}\subsection*{Semantic realization criterion}
\hyperref[rec:PSC-I-D-SEMANTIC-REALIZATION]{\nolinkurl{PSC-I-D-SEMANTIC-REALIZATION}}
\begin{lstlisting}
{
  "params": [
    "Theta:TheoremRecord",
    "F:RealizationFunctor"
  ],
  "returns": "SemanticRealizationCriterion(Theta,F)"
}
\end{lstlisting}
\clearpage\section{Complete action registry}
The registry below contains all 37 action names. ``Implemented'' describes the frozen runtime dispatcher. It does not assert that all hypotheses of every linked source record are automatically validated. Record associations are links for inspection.
\Needspace{6\baselineskip}\subsection*{\nolinkurl{AGGREGATE_PACKETS}}
\field{Runtime implemented}{No}
\field{Modes}{\nolinkurl{PLANNER_ONLY}}
\field{Records}{\hyperref[rec:PSC-I-L-PACKET-AGGREGATION]{\nolinkurl{PSC-I-L-PACKET-AGGREGATION}}}
\Needspace{6\baselineskip}\subsection*{\nolinkurl{ANALYZE_PACKETS}}
\field{Runtime implemented}{Yes}
\field{Modes}{\nolinkurl{RUNTIME_EXECUTABLE_EXACT_FINITE}}
\field{Records}{\hyperref[rec:PSC-I-C-PACKET-UNIVERSAL-PROPERTY]{\nolinkurl{PSC-I-C-PACKET-UNIVERSAL-PROPERTY}}; \hyperref[rec:PSC-I-C-PACKETWISE-EQUALITY]{\nolinkurl{PSC-I-C-PACKETWISE-EQUALITY}}; \hyperref[rec:PSC-I-D-ANALYSIS-SYNTHESIS]{\nolinkurl{PSC-I-D-ANALYSIS-SYNTHESIS}}; \hyperref[rec:PSC-I-T-EXACT-PACKET-COORDINATES]{\nolinkurl{PSC-I-T-EXACT-PACKET-COORDINATES}}}
\Needspace{6\baselineskip}\subsection*{\nolinkurl{ASSEMBLE_BLOCK_ALGEBRA}}
\field{Runtime implemented}{No}
\field{Modes}{\nolinkurl{PLANNER_ONLY}}
\field{Records}{\hyperref[rec:PSC-I-C-DIMENSION-COUNT]{\nolinkurl{PSC-I-C-DIMENSION-COUNT}}; \hyperref[rec:PSC-I-T-BLOCK-ALGEBRA]{\nolinkurl{PSC-I-T-BLOCK-ALGEBRA}}}
\Needspace{6\baselineskip}\subsection*{\nolinkurl{BLOCKIFY_OPERATOR}}
\field{Runtime implemented}{Yes}
\field{Modes}{\nolinkurl{RUNTIME_EXECUTABLE_EXACT_FINITE}}
\field{Records}{\hyperref[rec:PSC-I-D-TRANSFER-BLOCK]{\nolinkurl{PSC-I-D-TRANSFER-BLOCK}}; \hyperref[rec:PSC-I-T-OPERATOR-RECONSTRUCTION]{\nolinkurl{PSC-I-T-OPERATOR-RECONSTRUCTION}}}
\Needspace{6\baselineskip}\subsection*{\nolinkurl{BUILD_SEED_HULL}}
\field{Runtime implemented}{No}
\field{Modes}{\nolinkurl{EXACT_FIXTURE_ONLY}, \nolinkurl{PLANNER_ONLY}}
\field{Records}{\hyperref[rec:PSC-I-D-GENERATED-HULL]{\nolinkurl{PSC-I-D-GENERATED-HULL}}; \hyperref[rec:PSC-I-D-SEED-EVALUATION]{\nolinkurl{PSC-I-D-SEED-EVALUATION}}}
\Needspace{6\baselineskip}\subsection*{\nolinkurl{BUILD_SUPPORT_ALGEBRA}}
\field{Runtime implemented}{No}
\field{Modes}{\nolinkurl{PLANNER_ONLY}}
\field{Records}{\hyperref[rec:PSC-I-C-SUPPORT-PATTERN-STAR]{\nolinkurl{PSC-I-C-SUPPORT-PATTERN-STAR}}; \hyperref[rec:PSC-I-T-PREORDER-ALGEBRA-DECOMPOSITION]{\nolinkurl{PSC-I-T-PREORDER-ALGEBRA-DECOMPOSITION}}; \hyperref[rec:PSC-I-T-SUPPORT-PATTERN-ALGEBRA]{\nolinkurl{PSC-I-T-SUPPORT-PATTERN-ALGEBRA}}}
\Needspace{6\baselineskip}\subsection*{\nolinkurl{BUILD_TRANSFER_GRAPH}}
\field{Runtime implemented}{No}
\field{Modes}{\nolinkurl{PLANNER_ONLY}}
\field{Records}{\hyperref[rec:PSC-I-D-TRANSFER-GRAPH]{\nolinkurl{PSC-I-D-TRANSFER-GRAPH}}}
\Needspace{6\baselineskip}\subsection*{\nolinkurl{CHANGE_PACKET_COORDINATES}}
\field{Runtime implemented}{No}
\field{Modes}{\nolinkurl{EXACT_FIXTURE_ONLY}}
\field{Records}{\hyperref[rec:PSC-I-T-COORDINATE-TRANSPORT]{\nolinkurl{PSC-I-T-COORDINATE-TRANSPORT}}}
\Needspace{6\baselineskip}\subsection*{\nolinkurl{CHECK_REALIZATION_PROFILE}}
\field{Runtime implemented}{Yes}
\field{Modes}{\nolinkurl{RUNTIME_EXECUTABLE_SEMANTIC_GATE}}
\field{Records}{\hyperref[rec:PSC-I-D-REALIZATION-FUNCTOR]{\nolinkurl{PSC-I-D-REALIZATION-FUNCTOR}}; \hyperref[rec:PSC-I-D-SEMANTIC-REALIZATION]{\nolinkurl{PSC-I-D-SEMANTIC-REALIZATION}}}
\Needspace{6\baselineskip}\subsection*{\nolinkurl{COMBINE_SEED_SYSTEMS}}
\field{Runtime implemented}{No}
\field{Modes}{\nolinkurl{EXACT_FIXTURE_ONLY}, \nolinkurl{PLANNER_ONLY}}
\field{Records}{\hyperref[rec:PSC-I-P-SHARED-DIRECT-SUM]{\nolinkurl{PSC-I-P-SHARED-DIRECT-SUM}}; \hyperref[rec:PSC-I-P-TENSOR-BOUND]{\nolinkurl{PSC-I-P-TENSOR-BOUND}}; \hyperref[rec:PSC-I-T-EXACT-SEQUENCE-BOUNDS]{\nolinkurl{PSC-I-T-EXACT-SEQUENCE-BOUNDS}}; \hyperref[rec:PSC-I-T-INDEPENDENT-DIRECT-SUM]{\nolinkurl{PSC-I-T-INDEPENDENT-DIRECT-SUM}}}
\Needspace{6\baselineskip}\subsection*{\nolinkurl{COMPLETE_PACKET_SPLIT}}
\field{Runtime implemented}{No}
\field{Modes}{\nolinkurl{PLANNER_ONLY}}
\field{Records}{\hyperref[rec:PSC-I-C-COMPLETION-PACKET-SPLIT]{\nolinkurl{PSC-I-C-COMPLETION-PACKET-SPLIT}}; \hyperref[rec:PSC-I-T-COUNTABLE-UNCONDITIONAL-PACKETS]{\nolinkurl{PSC-I-T-COUNTABLE-UNCONDITIONAL-PACKETS}}}
\Needspace{6\baselineskip}\subsection*{\nolinkurl{COMPLETE_TOPOLOGICAL_REALIZATION}}
\field{Runtime implemented}{No}
\field{Modes}{\nolinkurl{PLANNER_ONLY}}
\field{Records}{\hyperref[rec:PSC-I-C-TOPOLOGICAL-PACKET-EXACT-SEQUENCE]{\nolinkurl{PSC-I-C-TOPOLOGICAL-PACKET-EXACT-SEQUENCE}}; \hyperref[rec:PSC-I-T-DENSE-REALIZATION]{\nolinkurl{PSC-I-T-DENSE-REALIZATION}}; \hyperref[rec:PSC-I-T-QUOTIENT-TOPOLOGY-COMPLETION]{\nolinkurl{PSC-I-T-QUOTIENT-TOPOLOGY-COMPLETION}}}
\Needspace{6\baselineskip}\subsection*{\nolinkurl{COMPOSE_TRANSFER_BLOCKS}}
\field{Runtime implemented}{Yes}
\field{Modes}{\nolinkurl{RUNTIME_EXECUTABLE_EXACT_FINITE}}
\field{Records}{\hyperref[rec:PSC-I-C-TENSOR-TRANSFER-BLOCKS]{\nolinkurl{PSC-I-C-TENSOR-TRANSFER-BLOCKS}}; \hyperref[rec:PSC-I-T-BLOCK-MULTIPLICATION]{\nolinkurl{PSC-I-T-BLOCK-MULTIPLICATION}}}
\Needspace{6\baselineskip}\subsection*{\nolinkurl{COMPUTE_PACKET_PURE_RANK}}
\field{Runtime implemented}{No}
\field{Modes}{\nolinkurl{EXACT_FIXTURE_ONLY}, \nolinkurl{PLANNER_ONLY}}
\field{Records}{\hyperref[rec:PSC-I-C-SOURCE-COMPONENT-COUNT]{\nolinkurl{PSC-I-C-SOURCE-COMPONENT-COUNT}}; \hyperref[rec:PSC-I-D-DIRECTED-SEED-RANK]{\nolinkurl{PSC-I-D-DIRECTED-SEED-RANK}}; \hyperref[rec:PSC-I-P-PACKET-PURE-RANK-BOUNDS]{\nolinkurl{PSC-I-P-PACKET-PURE-RANK-BOUNDS}}; \hyperref[rec:PSC-I-T-SOURCE-COMPONENT-SEED-RANK]{\nolinkurl{PSC-I-T-SOURCE-COMPONENT-SEED-RANK}}}
\Needspace{6\baselineskip}\subsection*{\nolinkurl{COMPUTE_REACHABILITY}}
\field{Runtime implemented}{No}
\field{Modes}{\nolinkurl{EXACT_FIXTURE_ONLY}, \nolinkurl{PLANNER_ONLY}}
\field{Records}{\hyperref[rec:PSC-I-P-REACHABILITY-OBSTRUCTION]{\nolinkurl{PSC-I-P-REACHABILITY-OBSTRUCTION}}; \hyperref[rec:PSC-I-T-ALGEBRA-REACHABILITY-SUPPORT]{\nolinkurl{PSC-I-T-ALGEBRA-REACHABILITY-SUPPORT}}}
\Needspace{6\baselineskip}\subsection*{\nolinkurl{COMPUTE_SEED_RANK}}
\field{Runtime implemented}{No}
\field{Modes}{\nolinkurl{EXACT_FIXTURE_ONLY}, \nolinkurl{PLANNER_ONLY}}
\field{Records}{\hyperref[rec:PSC-I-C-ALGEBRAICALLY-CLOSED-SINGLE-OPERATOR]{\nolinkurl{PSC-I-C-ALGEBRAICALLY-CLOSED-SINGLE-OPERATOR}}; \hyperref[rec:PSC-I-C-INVARIANT-FACTOR-COUNT]{\nolinkurl{PSC-I-C-INVARIANT-FACTOR-COUNT}}; \hyperref[rec:PSC-I-D-DIRECTED-SEED-RANK]{\nolinkurl{PSC-I-D-DIRECTED-SEED-RANK}}; \hyperref[rec:PSC-I-D-SEED-RANK]{\nolinkurl{PSC-I-D-SEED-RANK}}; \hyperref[rec:PSC-I-P-CYCLIC-CHARACTERIZATION]{\nolinkurl{PSC-I-P-CYCLIC-CHARACTERIZATION}}; \hyperref[rec:PSC-I-T-MODULE-TOP]{\nolinkurl{PSC-I-T-MODULE-TOP}}; \hyperref[rec:PSC-I-T-SEMISIMPLE-SEED-RANK]{\nolinkurl{PSC-I-T-SEMISIMPLE-SEED-RANK}}; \hyperref[rec:PSC-I-T-SINGLE-OPERATOR]{\nolinkurl{PSC-I-T-SINGLE-OPERATOR}}; \hyperref[rec:PSC-I-T-SOURCE-COMPONENT-SEED-RANK]{\nolinkurl{PSC-I-T-SOURCE-COMPONENT-SEED-RANK}}}
\Needspace{6\baselineskip}\subsection*{\nolinkurl{DECOMPOSE_PRIMARY}}
\field{Runtime implemented}{No}
\field{Modes}{\nolinkurl{PLANNER_ONLY}}
\field{Records}{\hyperref[rec:PSC-I-C-ALGEBRAICALLY-CLOSED-SINGLE-OPERATOR]{\nolinkurl{PSC-I-C-ALGEBRAICALLY-CLOSED-SINGLE-OPERATOR}}; \hyperref[rec:PSC-I-C-CYCLIC-MATRIX]{\nolinkurl{PSC-I-C-CYCLIC-MATRIX}}; \hyperref[rec:PSC-I-C-INVARIANT-FACTOR-COUNT]{\nolinkurl{PSC-I-C-INVARIANT-FACTOR-COUNT}}; \hyperref[rec:PSC-I-T-SINGLE-OPERATOR]{\nolinkurl{PSC-I-T-SINGLE-OPERATOR}}}
\Needspace{6\baselineskip}\subsection*{\nolinkurl{DESCEND_PACKET_QUOTIENT}}
\field{Runtime implemented}{No}
\field{Modes}{\nolinkurl{EXACT_FIXTURE_ONLY}}
\field{Records}{\hyperref[rec:PSC-I-C-QUOTIENT-PACKET-DESCENT]{\nolinkurl{PSC-I-C-QUOTIENT-PACKET-DESCENT}}; \hyperref[rec:PSC-I-T-QUOTIENT-PACKET-DESCENT]{\nolinkurl{PSC-I-T-QUOTIENT-PACKET-DESCENT}}}
\Needspace{6\baselineskip}\subsection*{\nolinkurl{DUALIZE_PACKETS}}
\field{Runtime implemented}{No}
\field{Modes}{\nolinkurl{PLANNER_ONLY}}
\field{Records}{\hyperref[rec:PSC-I-P-DUAL-PACKET-SYSTEM]{\nolinkurl{PSC-I-P-DUAL-PACKET-SYSTEM}}}
\Needspace{6\baselineskip}\subsection*{\nolinkurl{EVALUATE_PACKET_CONDITIONING}}
\field{Runtime implemented}{No}
\field{Modes}{\nolinkurl{PLANNER_ONLY}}
\field{Records}{\hyperref[rec:PSC-I-D-PACKET-CONDITION-NUMBER]{\nolinkurl{PSC-I-D-PACKET-CONDITION-NUMBER}}; \hyperref[rec:PSC-I-P-CONDITIONED-COORDINATES]{\nolinkurl{PSC-I-P-CONDITIONED-COORDINATES}}}
\Needspace{6\baselineskip}\subsection*{\nolinkurl{EXPAND_OPERATOR_PATHS}}
\field{Runtime implemented}{No}
\field{Modes}{\nolinkurl{EXACT_FIXTURE_ONLY}}
\field{Records}{\hyperref[rec:PSC-I-C-BLOCK-POWERS]{\nolinkurl{PSC-I-C-BLOCK-POWERS}}}
\Needspace{6\baselineskip}\subsection*{\nolinkurl{EXPORT_PSCI_INTERFACE}}
\field{Runtime implemented}{No}
\field{Modes}{\nolinkurl{PLANNER_ONLY}}
\field{Records}{\hyperref[rec:PSC-I-EXPORT]{\nolinkurl{PSC-I-EXPORT}}}
\Needspace{6\baselineskip}\subsection*{\nolinkurl{EXTRACT_PACKET_COMMUTANT}}
\field{Runtime implemented}{No}
\field{Modes}{\nolinkurl{PLANNER_ONLY}}
\field{Records}{\hyperref[rec:PSC-I-P-BLOCK-DIAGONAL-COMMUTANT]{\nolinkurl{PSC-I-P-BLOCK-DIAGONAL-COMMUTANT}}}
\Needspace{6\baselineskip}\subsection*{\nolinkurl{FACTOR_ACTION}}
\field{Runtime implemented}{No}
\field{Modes}{\nolinkurl{PLANNER_ONLY}}
\field{Records}{\hyperref[rec:PSC-I-P-ACTION-FACTORIZATION]{\nolinkurl{PSC-I-P-ACTION-FACTORIZATION}}}
\Needspace{6\baselineskip}\subsection*{\nolinkurl{LOWER_INTERTWINER}}
\field{Runtime implemented}{Yes}
\field{Modes}{\nolinkurl{RUNTIME_EXECUTABLE_EXACT_FINITE}}
\field{Records}{\hyperref[rec:PSC-I-T-BLOCK-INTERTWINING]{\nolinkurl{PSC-I-T-BLOCK-INTERTWINING}}}
\Needspace{6\baselineskip}\subsection*{\nolinkurl{PRESENT_AS_FREE_QUOTIENT}}
\field{Runtime implemented}{No}
\field{Modes}{\nolinkurl{EXACT_FIXTURE_ONLY}, \nolinkurl{PLANNER_ONLY}}
\field{Records}{\hyperref[rec:PSC-I-C-MINIMAL-FREE-COVER]{\nolinkurl{PSC-I-C-MINIMAL-FREE-COVER}}; \hyperref[rec:PSC-I-C-ONE-SEED-PRESENTATION]{\nolinkurl{PSC-I-C-ONE-SEED-PRESENTATION}}; \hyperref[rec:PSC-I-T-FREE-PRESENTATION]{\nolinkurl{PSC-I-T-FREE-PRESENTATION}}}
\Needspace{6\baselineskip}\subsection*{\nolinkurl{PROPAGATE_PACKET_SUPPORT}}
\field{Runtime implemented}{No}
\field{Modes}{\nolinkurl{EXACT_FIXTURE_ONLY}}
\field{Records}{\hyperref[rec:PSC-I-C-PACKET-SUPPORT-PROPAGATION]{\nolinkurl{PSC-I-C-PACKET-SUPPORT-PROPAGATION}}}
\Needspace{6\baselineskip}\subsection*{\nolinkurl{PUSH_GENERATORS_THROUGH_QUOTIENT}}
\field{Runtime implemented}{No}
\field{Modes}{\nolinkurl{EXACT_FIXTURE_ONLY}, \nolinkurl{PLANNER_ONLY}}
\field{Records}{\hyperref[rec:PSC-I-P-SCALAR-QUOTIENT-WITNESS]{\nolinkurl{PSC-I-P-SCALAR-QUOTIENT-WITNESS}}; \hyperref[rec:PSC-I-T-QUOTIENT-MONOTONICITY]{\nolinkurl{PSC-I-T-QUOTIENT-MONOTONICITY}}}
\Needspace{6\baselineskip}\subsection*{\nolinkurl{REALIZE_GRAMMAR}}
\field{Runtime implemented}{No}
\field{Modes}{\nolinkurl{PLANNER_ONLY}}
\field{Records}{\hyperref[rec:PSC-I-D-COLLAPSED-REALIZATION]{\nolinkurl{PSC-I-D-COLLAPSED-REALIZATION}}; \hyperref[rec:PSC-I-T-TYPED-MODULE-EQUIVALENCE]{\nolinkurl{PSC-I-T-TYPED-MODULE-EQUIVALENCE}}}
\Needspace{6\baselineskip}\subsection*{\nolinkurl{REDUCE_TO_SEMISIMPLE_TOP}}
\field{Runtime implemented}{No}
\field{Modes}{\nolinkurl{PLANNER_ONLY}}
\field{Records}{\hyperref[rec:PSC-I-C-LOCAL-ALGEBRA-FORMULA]{\nolinkurl{PSC-I-C-LOCAL-ALGEBRA-FORMULA}}; \hyperref[rec:PSC-I-T-MODULE-TOP]{\nolinkurl{PSC-I-T-MODULE-TOP}}; \hyperref[rec:PSC-I-T-SEMISIMPLE-SEED-RANK]{\nolinkurl{PSC-I-T-SEMISIMPLE-SEED-RANK}}}
\Needspace{6\baselineskip}\subsection*{\nolinkurl{RESOLVE_PACKET_EXACT_SEQUENCE}}
\field{Runtime implemented}{No}
\field{Modes}{\nolinkurl{EXACT_FIXTURE_ONLY}, \nolinkurl{PLANNER_ONLY}}
\field{Records}{\hyperref[rec:PSC-I-C-TOPOLOGICAL-PACKET-EXACT-SEQUENCE]{\nolinkurl{PSC-I-C-TOPOLOGICAL-PACKET-EXACT-SEQUENCE}}; \hyperref[rec:PSC-I-T-PACKET-EXACT-SEQUENCE]{\nolinkurl{PSC-I-T-PACKET-EXACT-SEQUENCE}}}
\Needspace{6\baselineskip}\subsection*{\nolinkurl{TENSOR_PACKETS}}
\field{Runtime implemented}{No}
\field{Modes}{\nolinkurl{EXACT_FIXTURE_ONLY}}
\field{Records}{\hyperref[rec:PSC-I-T-TENSOR-PACKET-SYSTEM]{\nolinkurl{PSC-I-T-TENSOR-PACKET-SYSTEM}}}
\Needspace{6\baselineskip}\subsection*{\nolinkurl{TRANSPORT_MORPHISM}}
\field{Runtime implemented}{No}
\field{Modes}{\nolinkurl{EXACT_FIXTURE_ONLY}, \nolinkurl{PLANNER_ONLY}}
\field{Records}{\hyperref[rec:PSC-I-D-PACKET-MORPHISM]{\nolinkurl{PSC-I-D-PACKET-MORPHISM}}; \hyperref[rec:PSC-I-D-SEED-MORPHISM]{\nolinkurl{PSC-I-D-SEED-MORPHISM}}; \hyperref[rec:PSC-I-P-SEED-SYSTEM-CATEGORY]{\nolinkurl{PSC-I-P-SEED-SYSTEM-CATEGORY}}; \hyperref[rec:PSC-I-T-NATURALITY]{\nolinkurl{PSC-I-T-NATURALITY}}}
\Needspace{6\baselineskip}\subsection*{\nolinkurl{TRANSPORT_SEED_RANK}}
\field{Runtime implemented}{No}
\field{Modes}{\nolinkurl{PLANNER_ONLY}}
\field{Records}{\hyperref[rec:PSC-I-C-ISOMORPHISM-INVARIANCE]{\nolinkurl{PSC-I-C-ISOMORPHISM-INVARIANCE}}; \hyperref[rec:PSC-I-P-SCALAR-EXTENSION-MONOTONICITY]{\nolinkurl{PSC-I-P-SCALAR-EXTENSION-MONOTONICITY}}}
\Needspace{6\baselineskip}\subsection*{\nolinkurl{TRANSPORT_THEOREM}}
\field{Runtime implemented}{Yes}
\field{Modes}{\nolinkurl{RUNTIME_EXECUTABLE_SEMANTIC_GATE}}
\field{Records}{\hyperref[rec:PSC-I-P-REALIZATION-EXPORT]{\nolinkurl{PSC-I-P-REALIZATION-EXPORT}}; \hyperref[rec:PSC-I-T-REALIZATION-EXPORT]{\nolinkurl{PSC-I-T-REALIZATION-EXPORT}}}
\Needspace{6\baselineskip}\subsection*{\nolinkurl{VALIDATE_PACKET_SYSTEM}}
\field{Runtime implemented}{Yes}
\field{Modes}{\nolinkurl{RUNTIME_EXECUTABLE_EXACT_FINITE}}
\field{Records}{\hyperref[rec:PSC-I-D-PACKET-SYSTEM]{\nolinkurl{PSC-I-D-PACKET-SYSTEM}}}
\Needspace{6\baselineskip}\subsection*{\nolinkurl{VERIFY_PROPAGATION_HYPOTHESES}}
\field{Runtime implemented}{No}
\field{Modes}{\nolinkurl{PLANNER_ONLY}}
\field{Records}{\hyperref[rec:PSC-I-D-EDGE-SURJECTIVE]{\nolinkurl{PSC-I-D-EDGE-SURJECTIVE}}; \hyperref[rec:PSC-I-D-PACKET-PROPAGATING]{\nolinkurl{PSC-I-D-PACKET-PROPAGATING}}}
\clearpage\section{Operation and hypothesis vocabulary}
The frozen ontology separates structural operations from hypothesis classes. These names form a controlled vocabulary for profiles and source-derived annotations. They are not a complete formal semantics for every mathematical construction named.
\begin{longtable}{@{}p{0.60\textwidth}p{0.31\textwidth}@{}}\toprule Name & Class\\\midrule\endhead
\nolinkurl{identity} & structural\\
\nolinkurl{zero} & structural\\
\nolinkurl{finite_sum} & structural\\
\nolinkurl{additive_inverse} & structural\\
\nolinkurl{composition} & structural\\
\nolinkurl{scalar_action} & structural\\
\nolinkurl{direct_sum} & structural\\
\nolinkurl{finite_product} & structural\\
\nolinkurl{biproduct} & structural\\
\nolinkurl{kernel} & structural\\
\nolinkurl{cokernel} & structural\\
\nolinkurl{image} & structural\\
\nolinkurl{quotient} & structural\\
\nolinkurl{tensor_product} & structural\\
\nolinkurl{dual} & structural\\
\nolinkurl{packet_idempotents} & structural\\
\nolinkurl{transfer_blocks} & structural\\
\nolinkurl{support_relation} & structural\\
\nolinkurl{path_composition} & structural\\
\nolinkurl{polynomial_action} & structural\\
\nolinkurl{jacobson_radical} & structural\\
\nolinkurl{semisimple_decomposition} & structural\\
\nolinkurl{primary_decomposition} & structural\\
\nolinkurl{universal_property} & structural\\
\nolinkurl{functoriality} & structural\\
\nolinkurl{exact_sequence} & structural\\
\nolinkurl{closure} & structural\\
\nolinkurl{completion} & structural\\
\nolinkurl{topology} & structural\\
\nolinkurl{norm} & structural\\
\nolinkurl{adjoint} & structural\\
\nolinkurl{dimension} & structural\\
\nolinkurl{finite_indexing} & hypothesis\\
\nolinkurl{surjectivity} & hypothesis\\
\nolinkurl{injectivity} & hypothesis\\
\nolinkurl{isomorphism} & hypothesis\\
\nolinkurl{continuity} & hypothesis\\
\nolinkurl{boundedness} & hypothesis\\
\nolinkurl{closedness} & hypothesis\\
\nolinkurl{density} & hypothesis\\
\nolinkurl{orthogonality} & hypothesis\\
\nolinkurl{self_adjointness} & hypothesis\\
\nolinkurl{positivity} & hypothesis\\
\nolinkurl{integrality} & hypothesis\\
\nolinkurl{unconditional_convergence} & hypothesis\\
\nolinkurl{hausdorff} & hypothesis\\
\nolinkurl{complete_field} & hypothesis\\
\nolinkurl{invariant_subspace} & hypothesis\\
\nolinkurl{commutation} & hypothesis\\
\nolinkurl{transitivity} & hypothesis\\
\nolinkurl{edge_surjectivity} & hypothesis\\
\nolinkurl{packet_propagation} & hypothesis\\
\nolinkurl{nonzero_packet_spaces} & hypothesis\\
\bottomrule\end{longtable}\clearpage\section{Historical Lean record map}
This is the eight-entry mapping attached to the archived runtime. The evidence descriptions and source-match categories are reproduced from that map. A fresh build performed elsewhere must be documented separately.
\Needspace{7\baselineskip}\subsection*{Finite packet system}
\hyperref[rec:PSC-I-D-PACKET-SYSTEM]{\nolinkurl{PSC-I-D-PACKET-SYSTEM}}
\field{Source match}{\nolinkurl{structural_exact}}
\field{Stored certificate status}{\nolinkurl{compiled_in_verified_tranche}}
\field{Declarations}{\nolinkurl{Komorebi.Foundation.LinearProjectorFamily}}
ProjectorCore.lean is part of the v0.2.0 kernel-accepted project; structure axioms encode idempotence, pairwise annihilation, and completeness.
\Needspace{7\baselineskip}\subsection*{Packet analysis and synthesis}
\hyperref[rec:PSC-I-D-ANALYSIS-SYNTHESIS]{\nolinkurl{PSC-I-D-ANALYSIS-SYNTHESIS}}
\field{Source match}{\nolinkurl{stronger_typed_realization}}
\field{Stored certificate status}{\nolinkurl{kernel_accepted}}
\field{Declarations}{\nolinkurl{Komorebi.Foundation.LinearProjectorFamily.PacketSpace}; \nolinkurl{Komorebi.Foundation.LinearProjectorFamily.rangeAnalysis}; \nolinkurl{Komorebi.Foundation.LinearProjectorFamily.rangeSynthesis}}
ExactCoordinates.lean defines the packet space as the product of projector ranges, so both transforms carry range certificates.
\Needspace{7\baselineskip}\subsection*{Exact packet-coordinate isomorphism}
\hyperref[rec:PSC-I-T-EXACT-PACKET-COORDINATES]{\nolinkurl{PSC-I-T-EXACT-PACKET-COORDINATES}}
\field{Source match}{\nolinkurl{stronger_typed_realization}}
\field{Stored certificate status}{\nolinkurl{kernel_accepted}}
\field{Declarations}{\nolinkurl{Komorebi.Foundation.LinearProjectorFamily.exactPacketCoordinates}}
Explicitly listed as Kernel accepted in VERIFICATION\_REPORT.md; proves both inverse laws for range-restricted packet coordinates.
\Needspace{7\baselineskip}\subsection*{Packetwise equality}
\hyperref[rec:PSC-I-C-PACKETWISE-EQUALITY]{\nolinkurl{PSC-I-C-PACKETWISE-EQUALITY}}
\field{Source match}{\nolinkurl{one_direction_plus_trivial_reverse}}
\field{Stored certificate status}{\nolinkurl{compiled_in_verified_tranche}}
\field{Declarations}{\nolinkurl{Komorebi.Foundation.LinearProjectorFamily.eq_of_projectors_eq}}
The nontrivial direction is proved in ProjectorCore.lean; the reverse direction follows by congruence from equality.
\Needspace{7\baselineskip}\subsection*{Transfer block}
\hyperref[rec:PSC-I-D-TRANSFER-BLOCK]{\nolinkurl{PSC-I-D-TRANSFER-BLOCK}}
\field{Source match}{\nolinkurl{generalized_domain_codomain}}
\field{Stored certificate status}{\nolinkurl{compiled_in_verified_tranche}}
\field{Declarations}{\nolinkurl{Komorebi.Foundation.transferBlock}}
Definition allows distinct source and target packet systems; the PSC-I endomorphism case is a specialization.
\Needspace{7\baselineskip}\subsection*{Exact operator reconstruction}
\hyperref[rec:PSC-I-T-OPERATOR-RECONSTRUCTION]{\nolinkurl{PSC-I-T-OPERATOR-RECONSTRUCTION}}
\field{Source match}{\nolinkurl{generalized_domain_codomain}}
\field{Stored certificate status}{\nolinkurl{kernel_accepted}}
\field{Declarations}{\nolinkurl{Komorebi.Foundation.transferBlock_reconstruction}}
Explicitly listed as Kernel accepted; reconstructs any linear map from all source-target blocks using completeness of both packet families.
\Needspace{7\baselineskip}\subsection*{Packet matrix multiplication}
\hyperref[rec:PSC-I-T-BLOCK-MULTIPLICATION]{\nolinkurl{PSC-I-T-BLOCK-MULTIPLICATION}}
\field{Source match}{\nolinkurl{generalized_composition_form}}
\field{Stored certificate status}{\nolinkurl{kernel_accepted}}
\field{Declarations}{\nolinkurl{Komorebi.Foundation.transferBlock_comp}}
Explicitly listed as Kernel accepted; inserts the complete middle packet family and yields the local composition sum.
\Needspace{7\baselineskip}\subsection*{Support-pattern algebra criterion}
\hyperref[rec:PSC-I-T-SUPPORT-PATTERN-ALGEBRA]{\nolinkurl{PSC-I-T-SUPPORT-PATTERN-ALGEBRA}}
\field{Source match}{\nolinkurl{partial_sufficient_direction}}
\field{Stored certificate status}{\nolinkurl{kernel_accepted_related}}
\field{Declarations}{\nolinkurl{Komorebi.Algebra.MatrixSupported.mul}}
Verified tranche proves multiplicative closure for transitive support relations; it is related evidence, not a source-equivalent certificate for every direction of the PSC-I criterion.
\clearpage\section{Preserved verification receipt}
The following is the recorded successful runtime verification run. The detailed exact-fixture output, command lines, interpreter versions, earlier missing-dependency explanation, and frozen-byte audit are retained in the build results.

\begin{lstlisting}
ENTRYPOINT_VERSION=PASS
RECORD_COUNT=PASS
DEPENDENCY_ACYCLIC=PASS
DEPENDENCIES_RESOLVE=PASS
AST_COVERAGE=PASS
ACTION_REGISTRY=PASS
SOURCE_HASH=PASS
API_STATUS=PASS
API_RECORD=PASS
API_DEPS=PASS
API_AST=PASS
EXECUTE_BLOCKIFY=PASS
EXECUTE_PACKET_VALIDATE=PASS
TRANSPORT_GATE_STRUCTURED=PASS
EXACT_FIXTURE_REPLAY=PASS
PSC_I_MACHINE_V1_RELEASE_GATE=PASS

\end{lstlisting}
\subsection*{Interpreter provenance}

\begin{lstlisting}
3.12.14 (main, Aug 25 2026, 14:00:49) [Clang 22.1.3 ]
SymPy=1.14.0

\end{lstlisting}
\subsection*{Source and package preservation}
The companion's integrity audit compared each of the 65 original archived files with the copied payload. Every comparison passed. Generated Python bytecode caches, if any, are not part of the archived source and are excluded from the delivered package.

\textbf{Catalogue provenance.} The 108 source statements and 69 proof bodies are generated from the archived mathematical inventory. Their printed hashes identify that source wording. The catalogue removes outer statement environments and translates cross-references into links to stable records. Formal acceptance is documented separately in the verification index.
\end{document}
